\documentclass[oneside,a4paper,german,parskip=half,draft]{scrbook}
\usepackage[latin1]{inputenc}
\pdfcompresslevel=9

%Neue Deutsche Rechtschreibung
\usepackage{ngerman}

%Seitenheader
\usepackage{fancyhdr}
\pagestyle{fancy}
\rhead{\nouppercase{\rightmark}}
\lhead{\nouppercase{\leftmark}}

%Standartfont
\usepackage[T1]{fontenc}
%\usepackage{fourier}
\usepackage{libertine}

%Zusatzpaket f�r mathematische Ausdr�cke
\usepackage{amsmath}

%Zusatzfonts f�r mathbb usw.
\usepackage{amsfonts}
\usepackage{mathrsfs}

%Verbessertes Ref
\usepackage[german]{varioref}

%Links
\usepackage{hyperref}

%Stichwortverzeichnis
\usepackage{makeidx}
\makeindex

%Acronyme
\usepackage[nolist,nohyperlinks]{acronym}

%Anpassbare Enumerates/Itemizes
\usepackage{enumitem}

%Tikz/PGF Zeichnenpaket
\usepackage{pgf}
\usepackage{tikz}
\usetikzlibrary{mindmap,trees,decorations,decorations.pathreplacing,decorations.pathmorphing,calc,arrows,automata,through,backgrounds,shapes.geometric}

%F�r graphische Todos
\usepackage[colorinlistoftodos, obeyDraft]{todonotes}

%Paket zum Berechnen von Textbreiten und H�hen
\usepackage{calc}

%BibTeX
%\usepackage{cite}
%\usepackage{bibgerm}
%\bibliographystyle{gerplain}

%F�r das manipulieren von Captions in Figuren
\usepackage{caption}

%F�r das anpassen von theoremen
\usepackage{amsthm}

%F�r die farbigen Boxen
\usepackage{xcolor}
\usepackage{framed}

%Quellcode Einbindung
\usepackage{listings}

%Um die Seite in mehrere Spalten aufzuteilen
\usepackage{multicol}

%F�r Widerspruchszeichen (Blitz)
\usepackage{ulsy}

%F�r zwischenr�ume in W�rtern
\usepackage{xspace}

%F�r das umstellen zwischen Doppelten, EinEinhalben und Einzelnen Zeilenabst�nden
\usepackage{setspace}

%F�r Commandos mit zwei opotionalen Argumenten
\usepackage{twoopt}

{\catcode`\!=8 % funny catcode so ! will be a delimiter 
\catcode`\Q=3 % funny catcode so Q will be a delimiter 
\long\gdef\given#1{88\fi\Ifbl@nk#1QQQ\empty!} 
\long\gdef\blank#1{88\fi\Ifbl@nk#1QQ..!}% if null or spaces 
\long\gdef\nil#1{\IfN@Ught#1* {#1}!}% if null 
\long\gdef\IfN@Ught#1 #2!{\blank{#2}} 
\long\gdef\Ifbl@nk#1#2Q#3!{\ifx#3}% same as above 
}

%<Commandos>
%In geschweifte Klammern setzen
\newcommand{\gklamm}[1]{\ensuremath{\left\{#1\right\}}}

%In eckige Klammern setzen
\newcommand{\eklamm}[1]{\ensuremath{\left[#1\right]}}

%In runde Klammern setzen
\newcommand{\rklamm}[1]{\ensuremath{\left(#1\right)}}

%In Betragsstriche Setzen
\newcommand{\betrag}[1]{\ensuremath{\left|#1\right|}}

%Script zum Eingeben von l�ngeren Beispielen
\newcommand{\bsp}[3][]
{
\if\blank{#3}\textalign{\underline{#2}}{#1}\else
\underline{#2} #1
\par
\begingroup
\leftskip=1.28em
\setlist[1]{labelindent=1.28em, leftmargin=*}
#3
\par
\endgroup\fi
}

%Script um Text einzur�cken
\newcommand{\textalign}[2]{
\begin{minipage}[b]{\widthof{#1} + \widthof{\space}}
#1
\end{minipage}
\begin{minipage}[t]{\linewidth-\widthof{#1}-\widthof{\space}}
#2
\end{minipage}
}

%Script um Text einzur�cken ohne Ausgabe
\newcommand{\textfakealign}[3]{
\begin{minipage}[b]{\widthof{#1} + \widthof{\space}}
\if\blank{#2}$ $\else#2\fi
\end{minipage}
\begin{minipage}[t]{\linewidth-\widthof{#1}-\widthof{\space}}
#3
\end{minipage}
}

%Indexe
%Normaler Index
\newcommand{\indexn}[2][]{#2\if\blank{#1}\index{#2}\else\index{#1}\fi}
%Unterstrichener Index
\newcommand{\indexu}[2][]{\underline{#2}\if\blank{#1}\index{#2}\else\index{#1}\fi}
%Kursiver Index
\newcommand{\indexi}[2][]{\textit{#2}\if\blank{#1}\index{#2}\else\index{#1}\fi}
%Fetter Index
\newcommand{\indexb}[2][]{\textbf{#2}\if\blank{#1}\index{#2}\else\index{#1}\fi}

%Definitionen, Beispiele, S�tze
\newenvironment{fshaded}[2]{%
\def\FrameCommand{\fcolorbox{#1}{#2}}%
\MakeFramed {\FrameRestore}}%
{\endMakeFramed}

\theoremstyle{definition}
\newtheorem{definition}{\ac{Def.}}[chapter]
\newenvironment{fdefinition}[1][]{
\definecolor{shadecolor_definition}{rgb}{.85882,.89803,.94509}
%OLD: \definecolor{shadecolor_definition}{rgb}{.87,.94,.96}
\definecolor{framecolor_definition}{rgb}{.30980,.50588,.74117}
%OLD: \definecolor{framecolor_definition}{rgb}{0,0,0}
\begin{fshaded}{framecolor_definition}{shadecolor_definition}\begin{definition}[#1]}{\end{definition}\end{fshaded}}

\newtheorem{lemma}[definition]{Lemma}

\newtheorem{satz}{Satz}[chapter]
\newenvironment{fsatz}[1][]{
\definecolor{shadecolor_satz}{rgb}{.91764,.94509,.86666}
%OLD: \definecolor{shadecolor_satz}{rgb}{.87,.96,.56}
\definecolor{framecolor_satz}{rgb}{.60784,.73333,.34901}
%OLD: \definecolor{framecolor_satz}{rgb}{0,0,0}
\begin{fshaded}{framecolor_satz}{shadecolor_satz}\begin{satz}[#1]}{\end{satz}\end{fshaded}}

\newtheoremstyle{beweis}% name
	{3pt}			%Space above
	{3pt}			%Space below
	{}				%Body font
	{}				%Indent amount (empty = no indent, \parindent = para indent)
	{\sc}	%Thm head font
	{:}			%Punctuation after thm head
	{.5em}		%Space after thm head: " " = normal interword space; \newline = linebreak
	{}				%Thm head spec (can be left empty, meaning `normal')

\theoremstyle{beweis}
\newtheorem{beweis}{Beweis}[chapter]
\newenvironment{fbeweis}[1][]{
\definecolor{shadecolor_beweis}{rgb}{.94901,.85882,.85882}
%OLD: \definecolor{shadecolor_beweis}{rgb}{1.00,.96,.92}
\definecolor{framecolor_beweis}{rgb}{.75294,.31372,.30196}
%OLD: \definecolor{framecolor_beweis}{rgb}{0,0,0}
\begin{fshaded}{framecolor_beweis}{shadecolor_beweis}\begin{beweis}[#1]}{\end{beweis}\end{fshaded}}

\theoremstyle{remark}
\newtheorem{notation}{Notation}[chapter]

\newtheoremstyle{note}% name
	{3pt}			%Space above
	{3pt}			%Space below
	{}				%Body font
	{}				%Indent amount (empty = no indent, \parindent = para indent)
	{\bfseries}	%Thm head font
	{:}			%Punctuation after thm head
	{.5em}		%Space after thm head: " " = normal interword space; \newline = linebreak
	{}				%Thm head spec (can be left empty, meaning `normal')

\newtheoremstyle{notelite}% name
	{3pt}			%Space above
	{3pt}			%Space below
	{}				%Body font
	{}				%Indent amount (empty = no indent, \parindent = para indent)
	{\itshape}	%Thm head font
	{:}			%Punctuation after thm head
	{.5em}		%Space after thm head: " " = normal interword space; \newline = linebreak
	{}				%Thm head spec (can be left empty, meaning `normal')

\theoremstyle{note}
\newtheorem*{beispiel}{Beispiel}

\theoremstyle{notelite}
\newtheorem*{anmerkung}{Anmerkung}
\newtheorem*{beobachtung}{Beobachtung}
\newtheorem*{bekannt}{Bekannt}
\newtheorem*{idee}{Idee}
\newtheorem*{frage}{Frage}
\newtheorem*{problem}{Problem}
\newtheorem*{also}{Also}
\newtheorem*{bemerkung}{Bemerkung}
\newtheorem*{ziel}{Ziel}
%%%%%%%%Arabische in R�mische Zahl umwandeln
\newcommand{\RM}[1]{\ensuremath{\mbox{\MakeUppercase{\romannumeral #1}}}}
%</Commandos>

%<Abk�rzungen>
\newcommand{\Ra}{\ensuremath{\Rightarrow}}
\newcommand{\ra}{\ensuremath{\rightarrow}}
\newcommand{\La}{\ensuremath{\Leftarrow}}
\newcommand{\la}{\ensuremath{\leftarrow}}
\newcommand{\Lra}{\ensuremath{\Leftrightarrow}}
\newcommand{\lra}{\ensuremath{\leftrightarrow}}
\newcommand{\hra}{\ensuremath{\hookrightarrow}}
\newcommand{\mal}{\ensuremath{\cdot}}
\newcommand{\klgl}{\ensuremath{\leq}}
\newcommand{\grgl}{\ensuremath{\geq}}
\newcommand{\oder}{\ensuremath{\vee}}
\newcommand{\und}{\ensuremath{\wedge}}
\newcommand{\ggq}[1]{"`#1"'}
\newcommand{\tx}[1]{\ensuremath{\text{#1}}}
\newcommand{\Sigs}{\ensuremath{\Sigma^*}}
\newcommand{\SigB}{\ensuremath{\Sigma_{\tx{Bool}}}}
\newcommand{\SigBs}{\ensuremath{\SigB^*}}
\newcommand{\ub}[2]{\ensuremath{\underbrace{#1}_{#2}}}
\newcommand{\tc}[2][red]{\textcolor{#1}{#2}}
\newcommand{\R}{\ensuremath{\mb{R}}}
\newcommand{\N}{\ensuremath{\mb{N}}}
\newcommand{\seq}{\ensuremath{\subseteq}}
\newcommand{\fnewline}{\mbox{}\par}
\newcommand{\ural}[2]{\ensuremath{\displaystyle \matholongrightarrow_{#1}^{#2}}}
\newcommand{\rkl}[1]{\ensuremath{\rklamm{#1}}}
\newcommand{\un}{\ensuremath{\infty}}
\newcommand{\notsubseteq}{\ensuremath{\not\subseteq}}
\newcommand{\qacc}{\ensuremath{q_{\tx{accept}}}}
\newcommand{\qrej}{\ensuremath{q_{\tx{reject}}}}
\newcommand{\txcol}[2]{\textcolor{#1}{#2}}

%Underbrace mit viel text
\newcommand{\ubl}[3][]{
\ub{#2}{\begin{minipage}{\if\blank{#1}\widthof{\ensuremath{#2}}\else#1\fi}\begin{center}\scriptsize{#3}\end{center}\end{minipage}}
}

%Stack mit viel text
\newcommand{\stackl}[3][]{
\stack{\begin{minipage}{\if\blank{#1}\widthof{\ensuremath{#3}}\else#1\fi}\begin{center}\scriptsize{#2}\end{center}\end{minipage}}{#3}
}
%Sin, Cos, Tan, Dim, Span, Arccos
\DeclareMathOperator{\spano}{span}
\DeclareMathOperator{\bin}{Bin}
\DeclareMathOperator{\ld}{ld}
\DeclareMathOperator{\prog}{Prog}
\DeclareMathOperator*{\vd}{\vdash}
\DeclareMathOperator*{\matholongrightarrow}{\longrightarrow}
\DeclareMathOperator{\Konf}{Konf}
\DeclareMathOperator{\Def}{Def}
\DeclareMathOperator{\Code}{Code}

\newcommand{\sinx}[1]{\ensuremath{\sin{\left(#1\right)}}}
\newcommand{\cosx}[1]{\ensuremath{\cos{\left(#1\right)}}}
\newcommand{\tanx}[1]{\ensuremath{\tan{\left(#1\right)}}}
\newcommand{\dimx}[1]{\ensuremath{\dim{\left(#1\right)}}}
\newcommand{\spanx}[1]{\ensuremath{\spano{\left(#1\right)}}}
\newcommand{\arccosx}[1]{\ensuremath{\arccos{\left(#1\right)}}}
\newcommand{\stack}[2]{\ensuremath{\stackrel{#1}{#2}}}

%Todo, Insert, Mark, Img, Solved
\newcommand{\Todo}[1]{\todo[inline]{TODO: #1}}
\newcommand{\Insert}[1]{\todo[inline, color=green!40]{INSERT: #1}}
\newcommand{\Mark}[1]{\todo[inline, nolist, color=blue!40]{MARK: #1}}
\newcommand{\Img}[1]{\todo[inline, nolist, color=black!40]{IMG: #1}}
\newcommand{\Solved}[2][]{\todo[inline, nolist]{TODO: #2\if\empty{#1} -- DONE\else\\SOLVED: #1\fi}}

%Schriften
\newcommand{\mb}[1]{\ensuremath{\mathbb{#1}}}
\newcommand{\mc}[1]{\ensuremath{\mathcal{#1}}}

\newcommand{\matrixp}[1]{\ensuremath{\begin{pmatrix} #1 \end{pmatrix}}}

%Tikz Abk�rzungen
%%F�r Automaten Graphen
\newcommand{\tikzedge}[5]{\path[->] (#1) edge[#4] node[#5] {$#3$} (#2);}
\newcommand{\tikzloop}[4]{\tikzedge{#1}{}{#2}{loop #3}{#4};}

%%F�r Mathematische Graphen
%%%Coordinatensystem
\newcommandtwoopt{\tikzcoor}[4][x][y]{\draw[<->,thick] (0,#4) node[right] {$#2$} -- (0,0) -- (#3,0) node[below] {$#1$};}
%</Abk�rzungen>


%<Einstellung f�r Hyperref>%%%%%%%%%%%%%%%%%%%%%%%%%%%%%%%%%%%%%%%%%%%%%%%%%%%%
\hypersetup{colorlinks=false, linkcolor=black, breaklinks=true, bookmarksdepth=3,unicode=true,bookmarksnumbered=true,
pdftitle={Hefter f�r Mathematik 2 -- Stand: \today},
pdfauthor={Thaller Alexander},
pdfsubject={Hefter f�r das Fach Mathematik 2 aus der Vorlesung von Professorin Gabriela Gr�niger (Gr�g) f�r das 2. Semester der Informatik im Sommersemester 2009 an der Hochschule Regensburg.},
pdfkeywords={informatik,studium,hefter,mathematik2,sommersemester,2009}}
%</Einstellung f�r Hyperref>%%%%%%%%%%%%%%%%%%%%%%%%%%%%%%%%%%%%%%%%%%%%%%%%%%%

%<Dokument Daten>%%%%%%%%%%%%%%%%%%%%%%%%%%%%%%%%%%%%%%%%%%%%%%%%%%%%%%%%%%%%%%
\titlehead{\includegraphics[width=34mm]{FH-Logo.pdf}}
\subject{Informatik (Bachelor) 2. Semester}
\author{Alexander Thaller}
\title{Mathematik 2\footnote{Gefundene Fehler oder Verbesserungsvorschlge bitte hier \href{http://fhrein2ss09.codeplex.com/WorkItem/List.aspx}{http://fhrein2ss09.codeplex.com/WorkItem/List.aspx} berichten oder alternativ mir eine E-Mail schreiben \href{mailto:alexander.thaller@stud.fh-regensburg.de}{alexander.thaller@stud.fh-regensburg.de}. Vielen Dank.}}
\subtitle{Hefter des Sommersemesters 2009}
\date{stand \today}
\publishers{Aus der Vorlesung von Professorin Gabriela Gr�niger (Gr�g)}
%</Dokument Daten>%%%%%%%%%%%%%%%%%%%%%%%%%%%%%%%%%%%%%%%%%%%%%%%%%%%%%%%%%%%%%
\begin{document}
%\begin{acronym}
	\acro{acr}{Akronym}
	\acro{z.B.}[\ensuremath{\tx{z.\,B.}\xspace}]{zum Beispiel}
	\acro{Def.}{Definition}
	\acro{elem.}{elementige}
	\acro{Teilm.}{Teilmenge}
	\acro{Anz.}{Anzahl}
	\acro{z.z.}[\ensuremath{\tx{z.\,z.}\xspace}]{zu zeigen}
	\acro{bzw.}{beziehungsweise}
	\acro{etc.}{et cetera}
	\acro{d.h.}[\ensuremath{\tx{d.\,h.}\xspace}]{daher}
	\acro{OBdA}{Ohne Beschr�nkung der Allgemeinheit}
	\acro{differenz.}{differenzierbar}
\end{acronym}

%\frontmatter
%\maketitle
%\setcounter{tocdepth}{1}
%\setcounter{secnumdepth}{2}
%\clearpage
%\pdfbookmark[0]{\contentsname}{tocanc}
%\tableofcontents
%\newpage
%\listoftodos
%\newpage
%<Text>%%%%%%%%%%%%%%%%%%%%%%%%%%%%%%%%%%%%%%%%%%%%%%%%%%%%%%%%%%%%%%%%%%%%%%%%
%\section*{�bersicht}
\begin{description}
\item[Professor:] Jorg Striegnitz
\item[Homepage:] \href{http://homepages.fh-regensburg.de/~stj39817/index.html}{http://homepages.fh-regensburg.de/~stj39817/index.html}
\item[E-Mail:] \href{mailto:joerg.striegnitz@informatik.fh-regensburg.de}{joerg.striegnitz@informatik.fh-regensburg.de}
\item[Zeitplan:]\mbox{}\par
		\begin{itemize}
		\item 4 Stunden pro Woche
		\item 2 Stunden Vorlesung
		\item 2 Stunden �bungen
		\end{itemize}
\item[Klausur:] Zulassungsvorrausetzungen: 50\% in den �bungen und eine Aufgabe Vorrechnen
\item[Tutoren:]\mbox{}\par
		\begin{itemize}
		\item Herr Bauer (Mo. 4. Stunde)
		\item Herr Lainer (Do. 6. Stunde)
		\item Frau Sporrer (U4M)
		\item Frau Franziska Brandstetter (Di. 2. Stunde) (\href{mailto:franziska.brandstetter@web.de}{franziska.brandstetter@web.de})
		\item Herr Brillner (U612) (\href{mailto:billner.michael@web.de}{billner.michael@web.de})
		\end{itemize}
\end{description}

%\mainmatter
%\part{Einf�hrung} \pagestyle{fancy}
%\Mark{Part 1}
%\chapter{Fakult�ten und Binominialkoeffizienten}
%MARK: Chapter 1
\begin{fdefinition}[Fakult�t]
%MARK: Definition 1.1
F�r $n \in \mathbb{N}_0$ definiert man die Fakult�t $n!$ von $n$ durch:
%MA2-18.03.2009-FORM1
\[n! = \begin{cases}
1 \text{ f�r } n = 0\\
(n - 1)! \mal n \text{ f�r } n \grgl 1
\end{cases}\]
\end{fdefinition}

\section{Anwendung der Fakult�t}
\begin{enumerate}
\item Anzahl aller Anordnungen einer Menge von Elementen $N = \gklamm{1, 2, 3, \dots, n}$ ist $n!$\\
		Eine Anordnung einer Menge nennt man auch \indexi{Permutation}\\
		\ac{z.B.} $n = 3$, $N = \gklamm{1, 2, 3}$

		\begin{figure}[h]
		\centering
		\begin{tikzpicture}[level 2/.style={sibling distance=10mm}]
			\node{1 \dots}[grow=right]
			child[->]{node{3 \dots}
				child[->]{node{32} 
					child[->]{node{321} edge from parent coordinate (desc3)}
				edge from parent 
				coordinate (desc2)}
				child[->]{node{31}
					child[->]{node{312}}}
			edge from parent 
			coordinate (desc1)}
%
			child[->]{node{2 \dots}
				child[->]{node{23}
					child[->]{node{231}}}
				child[->]{node{21}
					child[->]{node{213}}}}
%
			child[->]{node{1 \dots}
				child[->]{node{13}
					child[->]{node{132}}}
				child[->]{node{12}
					child[->]{node{123}}}};

		\draw[->] (desc1) +(-0.5,-0.5) node[below] {3} -- (desc1);
		\draw[->] (desc2) +(-0.5,-0.5) node[below] {2} -- (desc2);
		\draw[->] (desc3) +(0,-0.5) node[below] {1 = $3! = 6$} -- (desc3);
		\end{tikzpicture}
		\Img{MA2-18.03.2009-DIA2}
		\end{figure}

\item Anzahl M�glichkeiten aus einer $n$-\ac{elem.} Grundmenge $N = \gklamm{1, 2, 3, \dots, n}$ eine $k$-\ac{elem.} Teilmenge auszuw�hlen $(k \grgl n)$
		\ac{z.B.} $n = 49$, $k = 6$ (Lotto)\\
		$\text{Anzahl} = \frac{49 \mal 48 \mal 47 \mal 46 \mal 45 \mal 44}{6!} = 13938816 \approx 14 \text{Mio}$\\
		\underline{$n, k$} $\Ra$ $\text{Anzahl} = \frac{n \mal (n - 1) \mal (n - 2) \mal \dots \mal (n - k + 1)}{k!}$\\
\end{enumerate}

\begin{fdefinition}[Binominialkoeffizient]
%MARK: Definition 1.2
F�r $n, k \in \mathbb{N}_0$ und $0 \klgl k \klgl n$ definiert man den \indexu[Binominialkoeffizient]{Binominialkoeffizienten} $\binom{n}{k}$ (''$n$ �ber $k$'' oder ''$n$ tief $k$'') durch
\[\binom{n}{k} = \frac{n!}{k! \mal (n - k)!}\]
\end{fdefinition}

\begin{bemerkung}
\mbox{}\par
\[\frac{n \mal (n - 1) \mal (n - 2) \mal \dots \mal (n - k + 1)}{k!} \mal \frac{(n - k) \mal (n - k - 1) \mal \dots \mal 2 \mal 1}{(n - k) \mal (n - k - 1) \mal \dots \mal 2 \mal 1} = \frac{n!}{k! \mal (n - k)!}\]
\end{bemerkung}

\section{Aufgabe 1.1}
\label{sec:Fakultaet_und_Binominalkoeffizient_A1_1}
%MARK: Aufgabe 1.1
\begin{enumerate}[label=\alph*)]
\item Um von Punkt $0 = (0 \vert 0)$ zum Punkt $P = (6 \vert 4)$ zu gelangen, wollen wir uns nur auf Wegst�cken bewegen, die durch ganzzahlige Gitterpunkte des $\mathbb{R}^2$ gehen und parallel zu einer Koordinatenachse verlaufen.\\
		(Beispiel eines solchen Weges von 0 nach  P)
		\begin{figure}[h]
		\centering
		\begin{tikzpicture}
			\draw[step=1,gray,very thin] (0,0) grid (6,4);
			\draw (0,0) rectangle (6,4);
			\draw[ultra thick] (0,0) node[below]{$0$} -- (5,0) -- (5,3) -- (6,3) -- (6,4) node[above]{$P$};
		\end{tikzpicture}
		\end{figure}
		Auf wie vielen Wegen der minimalen L�nge $l = 6 + 4$ kann man von $0$ nach $P$ gelangen? (Vom Gitterpunkt $(i \vert  j)$ bewegt man sich also entweder zum Gitterpunkt $(i + 1 \vert j)$ oder $(i \vert j + 1)$.

\item Wie viele $12$-stellige Zahlen, bei denen die Ziffer $9$ genau zweimal auftritt, lassen sich aus den Ziffern $1, 2, \dots, 9$ bilden?\\
		(Das Produkt muss nicht ausmultipliziert werden.)
\end{enumerate}

L�sung siehe \vref{sec:Fakultaet_und_Binominalkoeffizient_A1_1L}.

\begin{fsatz}[Eigenschaften von Binominialkoeffizienten (Teil 1)]
%MARK: Satz 1.3 Teil 1
\label{satz:Fakultaet_und_Binominalkoeffizient_Satz1_3T1}
Es seien $n, k \in \mb{N}_0$, dann gilt:
\begin{enumerate}[label=\roman*)]
\item $\binom{n}{0} = \binom{n}{n} = 1$
\item $0 \klgl k \klgl n \binom{n}{k} = \binom{n}{n - k}$
\item $0 \klgl k \klgl k \klgl n - 1 \binom{n}{k} + \binom{n}{k + 1} = \binom{n + 1}{k + 1}$
		$\binom{n}{k} + \binom{n}{k + 1} = \frac{n!}{k! (n - k)!} + \frac{n!}{(k + 1)! \mal (n - k - 1)!}$\\
		$= \frac{n!}{k! \mal (n - k)!} \mal \left(\frac{1}{1 \mal (n - k)} + \frac{1}{(k + 1) \mal 1}\right)$\\
		$= \frac{n!}{k! \mal (n - k - 1)!} \mal \frac{k + 1 + n - k}{(n - k) (k + 1)} = \frac{(n + 1)!}{(k + 1)! \mal (n - k)!} = \binom{n + 1}{k + 1}$
\end{enumerate}
\end{fsatz}

\subsection*{Anschauliche Bedeutung von Satz \ref{satz:Fakultaet_und_Binominalkoeffizient_Satz1_3T1}}
F�r i), ii) betrachteten wir eine Menge $\gklamm{1, 2, 3, \dots, n}$
\begin{enumerate}[label=\roman*)]
\item Es gibt nur eine $0$-\ac{elem.} Teilmenge, n�mlich $\emptyset$ also ist $\binom{n}{0} = 1$\\
		Es gibt nur eine $n$-\ac{elem.} Teilmenge n�mlich $\gklamm{1, 2, \dots, n}$, damit $\binom{n}{n} = 1$

\item Wenn wir aus der Menge $\gklamm{1, 2, \dots, n}$ eine $k$-\ac{elem.} Teilmenge w�hlen, bleibt eine $n$-$k$-\ac{elem.} Teilmenge �brig.\\
		$n = 5$, $k = 3$ $\gklamm{1, 2, 3, 4 ,5}$
		\begin{table}[h]
		\centering
		\begin{tabular}{ccc}
		3-\ac{elem.} \ac{Teilm.} & & 2-\ac{elem.} \ac{Teilm.}\\
		$\gklamm{1,2,3}$ & $\ral$ & $\gklamm{4,5}$\\
		$\gklamm{1,2,4}$ & $\ral$ & $\gklamm{3,5}$\\
		$\gklamm{1,2,5}$ & $\ral$ & $\gklamm{3,4}$\\
		$\gklamm{1,3,4}$ & $\ral$ & $\gklamm{2,5}$\\
		$\gklamm{1,3,5}$ & $\ral$ & $\gklamm{2,4}$\\
		$\gklamm{1,4,5}$ & $\ral$ & $\gklamm{2,3}$\\
		$\gklamm{2,3,4}$ & $\ral$ & $\gklamm{1,5}$\\
		$\gklamm{2,3,5}$ & $\ral$ & $\gklamm{1,4}$\\
		$\gklamm{2,4,5}$ & $\ral$ & $\gklamm{1,3}$\\
		$\gklamm{3,4,5}$ & $\ral$ & $\gklamm{1,2}$
		\end{tabular}
		\end{table}
		$\binom{5}{3} = \frac{5 \mal 2 \mal 3}{1 \mal 2 \mal 3} = 10 = \binom{5}{2}$

\item Urne mit $n$ wei�en Kugeln mit Nummern $1, 2, \dots, n$ und einer gelben Kugeln mit Nummer $n + 1$.\\
		Wenn man aus dieser Urne $k + 1$ Kugeln zieht, gibt es daf�r $\binom{n + 1}{k + 1}$ M�glichkeiten.
		\begin{description}
		\item[Fall 1:] Gelbe Kugel wird gezogen, damit m�ssen noch $k$ wei�e Kugeln gezogen werden.\\
				\ac{Anz.} M�glichkeiten $\binom{n}{k}$

		\item[Fall 2:] Gelbe Kugel wird \underline{nicht} gezogen, damit m�ssen $k + 1$ wei�e Kugeln gezogen werden \ac{Anz.} M�glichkeiten $\binom{n}{k + 1}$
		\end{description}
		da Fall 1 und 2 nicht gleichzeitig auftreten ist die totale Anzahl $\binom{n}{k} + \binom{n}{k + 1} = \binom{n + 1}{k + 1}$
\end{enumerate}

\section{Pascalsches Dreieck}
Mit Hilfe von $\binom{n}{0} = \binom{n}{n} = 1$ und $\binom{n}{k} + \binom{n}{k + 1} = \binom{n + 1}{k + 1}$ kann man den Binominialkoeffizienten in einem Deiecksschema darstellen.\\
Die im Inneren liegenden Zahlen sind sind die Summe der jeweils dar�ber liegenden beiden Zahlen.
\begin{figure}[h]
\centering
\begin{tikzpicture}[scale=0.5]
	\foreach{\xa/\xb/\xc}in{0/8/2,-2/6/4,0/6/6,2/6/4,-1/5/10,1/5/10,-4/4/6,0/4/20,4/4/6}{\draw(\xa,\xb)node{\xc};}
	\draw(0,10)node (p1){1};
	\foreach{\xa/\xb}in{-1/1,1/1}{\draw(\xa,9)node(p2){\xb};}
	\foreach{\xa/\xb}in{-2/1,2/1}{\draw(\xa,8)node(p3){\xb};}
	\foreach{\xa/\xb}in{-3/1,-1/3,1/3,3/1}{\draw(\xa,7)node(p4){\xb};}
	\foreach{\xa/\xb}in{-4/1,4/1}{\draw(\xa,6)node(p5){\xb};}
	\foreach{\xa/\xb}in{-5/1,-3/5,3/5,5/1}{\draw(\xa,5)node(p6){\xb};}
	\foreach{\xa/\xb}in{-6/1,-2/15,2/15,6/1}{\draw(\xa,4)node(p7){\xb};}

	\foreach{\xa/\xb/\xc}in{0/10/0,-1/9/1,-2/8/2,-3/7/3,-4/6/4,-5/5/5,-6/4/6}{
		\node(n) at (\xa,\xb){};
		\draw[->] (n.west) +(-2,0) node[left] {$n = \xc$} -- (n.west);}

	\draw[->] (p1.north east) +(1,1) node[above] {$\binom{n}{0}$} -- (p1.north east);
	\draw[->] (p2.north east) +(1,1) node[above] {$\binom{n}{1}$} -- (p2.north east);
	\draw[->] (p3.north east) +(1,1) node[above] {$\binom{n}{2}$} -- (p3.north east);
	\draw[->] (p4.north east) +(1,1) node[above] {$\binom{n}{3}$} -- (p4.north east);
	\draw[->] (p5.north east) +(1,1) node[above] {$\binom{n}{4}$} -- (p5.north east);
	\draw[->] (p6.north east) +(1,1) node[above] {$\binom{n}{5}$} -- (p6.north east);
	\draw[->] (p7.north east) +(1,1) node[above] {$\binom{n}{6}$} -- (p7.north east);
\end{tikzpicture}
\end{figure}
\Img{MA2-19.03.2009-IMG1}

\subsection{Implementierung von Binominialkoeffizienten}
\begin{itemize}
\item falls man \underline{alle} Koeffizienten bis zu einem gegebenen $n$ haben will, erfolgt die Implementierung analog zum Pascalschen Dreieck.
\item falls man \underline{einen} Binominialkoeffizienten $\binom{n}{k}$ haben m�chte.\\
		Implementierung in Maple:
		\begin{lstlisting}
		n := ...; k := ...;
		b := 1:
		for i from 1 to n do
			b := b * (n - 1 + 1)/i;
		end do:
		print (b);
		\end{lstlisting}

		\begin{bemerkung}
		Alle Zwischenresultate $b$ sind $\in \mb{N}$
		\begin{itemize}
		\item $\frac{n}{1} \in \mb{N}$
		\item $\frac{n}{1} \mal \frac{(n - 1}{2} \in \mb{N}$, da entweder $n$ oder $n - 1$ gerade
		\item $\frac{n}{1} \mal \frac{(n - 1}{2} \mal \frac{n - 2}{3}$ da (genau) eine der drei Zahlen $n$, $n -1$, $n - 2$ durch $3$ teilbar
		\end{itemize}
		\end{bemerkung}
\end{itemize}

\begin{fsatz}[Binomischer Lehrsatz (Binomische Formel)]
%MARK: Satz 1.4
\label{satz:Fakultaet_und_Binominalkoeffizient_Satz1_4}
F�r $a, b \in \mb{R}$ und $n \in \mb{N}_0$ gilt:
\[\left(a + b\right)^n = \sum_{k = 0}^n \binom{n}{k} \mal a^{n - k} \mal b^k\]
\end{fsatz}

\begin{fbeweis}[Vollst�ndige Induktion]
\mbox{}\par
%MA2-19.03.2009-BEW2
\begin{description}
\item[(IA)] $n = 0$
		\begin{align*}
			LS &= (a + b)^0 = 1\\
			RS &= \sum_{k = 0}^0 \binom{0}{k} \mal a^{0 - k} \mal b^k = 1 \mal 1 \mal 1 = 1 \Ra \text{Wahr}
		\end{align*}

\item[(IS)] \ac{z.z.}: $\forall n \in \mb{N}_0$
		\[(a + b)^n = \sum_{k = 0}^n \mal a^{n - k} \mal b^k \Ra (a + b)^{n + 1} \binom{n + 1}{k} \mal a^{n + 1 - k} \mal b^k\]
		\begin{align*}
			(a + b)^{n + 1} &= (a + b) \mal (a + b)^n \stackrel{\text{(IV)}}{=} (a + b) \sum_{k = 0}^n \binom{n}{k} \mal a^{n - k} \mal b^k\\
			&= \sum_{k = 0}^n \binom{n}{k} \mal a^{n + 1 - k} \mal b^k + \sum_{k = 0}^n \binom{n}{k} \mal a^{n - k} \mal b^{k + 1}\\
			&= a^{n + 1} + \sum_{k = 0}^n \binom{n}{k} \mal a^{n + 1 - k} \mal b^k + b^{n + 1} + \sum_{k = 0}^n \binom{n}{k} \mal a^{n - k} \mal b^{k + 1}\\
			&= a^{n + 1} + b^{n + 1} + \sum_{k = 0}^n \binom{n}{k} \mal a^{n + 1 - k} \mal b^k + \sum_{\underbrace{k = 0}_{\begin{array}{cc}&l = k + 1\\\Lra & k = l - 1\end{array}}}^n \binom{n}{l - 1}\\
			&= \binom{n + 1}{0} a^{n + 1} + \binom{n + 1}{n + 1} b^{n + 1} + \sum_{k = 1}^n \underbrace{\left(\binom{n}{k} + \binom{n}{k - 1}\right)}_{\binom{n + 1}{k}} \mal a^{n + 1 - k} \mal b^k\\
			&= \sum_{k = 0}^{n + 1} \mal a^{n + 1 - k} \mal b^k
		\end{align*}
\end{description}
\end{fbeweis}

\subsection*{Anschauliche Begr�ndung f�r Satz \ref{satz:Fakultaet_und_Binominalkoeffizient_Satz1_4}}
$(a + b)^n = \underbrace{(a + b)}_{\text{1. Klammer}} \mal \underbrace{(a + b)}_{\text{2. Klammer}} \mal (a + b) \mal \dots \mal \underbrace{(a + b)}_{\text{$n$-te Klammer}}$
\begin{itemize}
\item Summe der Exponenten von $a$ und $b$ muss $n$ sein
\item $\binom{n}{k}$ \ac{Anz.} m�glicher Wahlen von $k$ Klammern aus der ''$b$'' gew�hlt wird.
\end{itemize}

\begin{beispiel}[Absch�tzung von Gr��enordnungen]
$(1,07)^{10} = 1,067$\\
$(1,07)^{10} = \left(1 + \frac{7}{100}\right)^{10} = \sum_{k = 0}^{10} \underbrace{\binom{10}{k} \mal 1^{10 - k} \mal \left(\frac{7}{100}\right)^k}_{\klgl 0} = *$\\
\begin{align*}
* &\klgl \binom{10}{0} \mal 1^{10} \mal \left(\frac{7}{100}\right)^0 + \binom{10}{1} \mal 1^9 \mal \left(\frac{7}{100}\right)^1\\
&= 1 + 0,7 = 1,7\\
* &\klgl 1 + 0,7 + \underbrace{\binom{10}{2}}_{45} \mal 1^8 \mal \underbrace{\left(\frac{7}{100}\right)^2}_{\frac{49}{10000}} \approx 1,7 + \underbrace{\frac{225}{10000}}_{0,225} = 1,925
\end{align*}
\end{beispiel}

\begin{fsatz}[Eigenschaften von Binominialkoeffizienten (Teil 2)]
\label{satz:Fakultaet_und_Binominalkoeffizient_Satz1_3T2}
Teil 1 siehe Satz \vref{satz:Fakultaet_und_Binominalkoeffizient_Satz1_3T1}
\begin{enumerate}[label = \roman*), start=4]
\item $\sum_{k = 0}^n \binom{n}{k} = (1 + 1)^n = 2^n$
\item $\sum_{k = 0}^n (-1)^k \binom{n}{k} = (1 + (-2))^n = 0$\\$n \grgl 1$
\end{enumerate}
%MARK: Satz 1.3 Teil 2
\end{fsatz}

\subsection*{Interpretation von Satz \ref{satz:Fakultaet_und_Binominalkoeffizient_Satz1_3T2} f�r Mengen}
$\binom{n}{k}$ = Anzahl $k$-\ac{elem.} Teilmengen einer $n$-\ac{elem.} Grundmenge\\
\begin{align*}
\Ra \sum_{k = 0}^n \binom{n}{k} &= \text{ Anzahl aller m�glichen Teilmengen einer $n$-\ac{elem.} Grundmenge}\\
&= \text{ Anzahl Elemente der Potenzmenge einer $n$-\ac{elem.} Menge}\\
&= 2^n
\end{align*}
\[\sum_{k = 0}^n (-1)^k \binom{n}{k} = \underbrace{\sum_{k = 0}^n \binom{n}{k}}_{k \text{ gerade}} - \underbrace{\sum_{k = 0}^n \binom{n}{k}}_{k \text{ ungerade}} = 0\]
\begin{itemize}[label=\Ra]
\item Anzahl Teilmengen der $n$-\ac{elem.} Grundmenge, die eine gerade Anzahl an Elementen enth�lt, ist gleich der Anzahl der Teilmengen mit einer ungeraden Anzahl an Elementen.
\end{itemize}

\section{L�sungen}
\subsection{Aufgabe 1.1}
\label{sec:Fakultaet_und_Binominalkoeffizient_A1_1L}
L�sung zu Aufgabe \vref{sec:Fakultaet_und_Binominalkoeffizient_A1_1}.

\begin{enumerate}[label=\alph*)]
\item Es werden 10 Schritte gelaufen, davon 4 vertikal (und 6 horizontal).\\
		Anzahl Pfade = Anzahl M�glichkeiten die 4 Schritte nach oben aus den 10 auszuw�hlen.
		\[= \binom{10}{4} = \frac{10 \mal 9 \mal 8 \mal 7}{4 \mal 3 \mal 2 \mal 1} = 210\]

\item $8^{10}$ = Anzahl der M�glichkeiten an den Stellen ohne 9.\\
		$\binom{12}{2}$ = Anzahl M�glichkeiten die Stellen f�r die 9en auszusuchen.\\
		Die Formel lautet damit wie folgt:
		\[8^{10} \mal \binom{12}{2} \approx 2,95 \mal 10^{10}\]
\end{enumerate}

%\chapter{Ungleichungen}
%MARK: Chapter 2
\begin{fdefinition}[Ordnungsaxiome]
%MARK: Definition 1.6
\label{def:Fakultaet_und_Binominalkoeffizient_Definition_1_6}
Es gibt eine Relation ''$<$'' in $\mb{R}$ mit den Eigenschaften
\begin{enumerate}[label=\roman*)]
\item f�r $a, b \in \mb{R}$ gilt \underline{genau eine} der Aussagen
		\[a < b~~a = b~~b < a\]
		
\item f�r $a, b, c \in \mb{R}$:
		\[a < b \und b < c \Ra a < c \text{ (Transitivit�t)}\]

\item f�r $a, b, c \in \mb{R}$:
		\[a < c \und b < d \Ra a + b < c + d\]
		Monotonie bez�glich der Addition
		
\item f�r $a, b, c \in \mb{R}$:
		\[a < b \und 0 < c \Ra ac < bc\]
		Monotonie bez�glich der Multiplikation
\end{enumerate}
\end{fdefinition}

\begin{bemerkung}
Es gelten folgende Schreibweisen und Vereinbarungen
\begin{enumerate}
\item $a \klgl b \Lra a < b \oder a = b$
\item $a > b \Lra b < a$
\item $a \grgl b \Lra b \klgl a$
\item $a < b \klgl c \Lra a < b \und b \klgl c$
\item Eine reelle Zahl $a$ hei�t \indexu{positiv}, falls $a > 0$
\item Eine reelle Zahl $a$ hei�t \indexu{negativ}, falls $a < 0$
\item Eine reelle Zahl $a$ hei�t \indexu{nichtnegativ} falls $a \grgl 0$
\item Eine reelle Zahl $a$ hei�t \indexu{nichtpositiv}, falls $a \klgl 0$
\end{enumerate}
\end{bemerkung}

\begin{bemerkung}Mit Definition \ref{def:Fakultaet_und_Binominalkoeffizient_Definition_1_6} und obiger Bemerkung kann man zeigen, dass ''$\klgl$'' auf $\mb{R}$ eine Ordnungsrelation darstellt, also gilt:
\begin{enumerate}[label=\roman*)]
\item reflexiv: $\forall a \in \mb{R}: a \klgl a$
\item transitiv: $\forall a, b, c \in \mb{R}: a \klgl b \und b \klgl c \Ra a \klgl c$
\item antitsymmetrisch: $\forall a, b \in \mb{R}: a \klgl b \und b \klgl a \Ra a = b$
\end{enumerate}
\end{bemerkung}

\begin{fsatz}[Rechenregel f�r Ungleichungen]
%MARK: Satz 1.7
\label{satz:Fakultaet_und_Binominalkoeffizient_Satz_1_7}
\mbox{}\par
\begin{enumerate}[label=\roman*)]
\item $a > 0 \Lra -a < 0$
\item $a^2 \grgl 0$
\item $a < b \und c < 0 \Ra ac > bc$
\item $0 < a < b \Lra 0 < \frac{1}{b} < \frac{1}{a}$
\item f�r $a, b \grgl 0: a < b \Lra a^2 < b^2$
\end{enumerate}
\end{fsatz}

\begin{fbeweis}[Nachrechnen]\mbox{}\par
\begin{enumerate}[label = \roman*), start=2]
\item \begin{description}
		\item[Fall 1:] $a = 0$
		\[a^2 = a \mal a = 0 \mal 0 = 0\]

		\item[Fall 2:] $a > 0$
		\begin{align*}
		&a > 0~~~~\vert \mal a \text{ Monotonie bez�glich Multiplikation}\\
		\Lra &a^2 > 0
		\end{align*}

		\item[Fall 3:] $a < 0$
		\begin{align*}
		&a < 0~~\vert -a \text{ Monotonie bez�glich Addition}\\
		\Lra &0 < -a~~~~\vert \mal (-a) \text{ Monotonie bez�glich Multiplikation}\\
		\Lra &0 > (-a)(-a)\\
		\Lra &0 < \underbrace{(-1)^2}_{= 1} \mal a^2
		\end{align*}
		\end{description}
\end{enumerate}\qed
\end{fbeweis}

\section{Aufgabe 1.3}
\label{sec:Ungleichungen_A1_3}
F�r welche Werte von $x$ und $y$ gilt: $xy > x$?

L�sung siehe \vref{sec:Ungleichungen_A1_3L}.

\section{Aufgabe 1.4}
\label{sec:Ungleichungen_A1_4}
Zeigen Sie: F�r positive Zahlen $x$ und $y$ gilt:
\[\frac{x}{y} + \frac{y}{x} \grgl 2\]
wobei Gleichheit genau dann Eintritt, wenn $x = y$ ist.

L�sung siehe \vref{sec:Ungleichungen_A1_4L}.

\begin{fdefinition}[Intervalle]
Die folgenden Mengen von reellen Zahlen bezeichnet man als Intervalle.
\begin{center}
\begin{multicols}{2}
$x \in (a, b) \Lra a < x < b$\\
$x \in [a, b) \Lra a \klgl x < b$\\
$x \in (a, b] \Lra a < x \klgl b$\\
$x \in [a, b] \Lra a \klgl x \klgl b$\\
$x \in (a, \infty) \Lra x > a$\\
$x \in (-\infty, b) \Lra x < b$\\
$x \in [a, \infty) \Lra x \grgl a$\\
$x \in (-\infty, b] \Lra x \klgl b$
\end{multicols}
\end{center}
\begin{itemize}
\item Intervalle der Form $(a, b)$ hei�en offene Intervalle
\item Intervalle der Form $[a, b]$ hei�en abgeschlossene Intervalle
\item Intervalle der Form $(a, b]$ \ac{bzw.} $[a, b)$ hei�en halboffene Intervalle
\item Die Intervalle auf der linken Seite hei�en endliche Intervalle
\end{itemize}
Statt $(a, b)$, $(a, b]$, $[a, b)$ kann man auch schreiben: $]a, b[$, $]a,b]$, $[a, b[$
%MARK: Definition 1.8
\end{fdefinition}

\begin{fsatz}[Ungleichung zwischen arithmetischem und geometrischem Mittel]
Seien $a, b > 0$, dann gilt:
\begin{itemize}
\item $\underbrace{\sqrt{a \mal b}}_{\text{geometrische Mittel}} \klgl \underbrace{\frac{1}{2} (a + b)}_{\text{arithmetisches Mittel}}$
\item $\sqrt{a \mal b} = \frac{1}{2} (a + b) \Lra a = b$
\end{itemize}
%MARK: Satz 1.9
\end{fsatz}

\textbf{Anschauliche Erkl�rung:}

\begin{tikzpicture}
	\draw[draw=gray] (0,0) arc (180:0:2);
	\draw[draw=green] (60:2) -- (1,0) node[right, midway, text=green] {$h$};
	\draw[draw=blue] (2,2) -- (2,0) node[right, near start, text=blue] {$r$};
	\draw (1,0) -- (1,-0.1);
	\draw[serif cm-] (0,0) -- (1,0) node[below, midway] {$a$};
	\draw[-serif cm] (1,0) -- (4,0) node[below, midway] {$b$};
	\draw (4,0) -- (60:2) -- (0,0);
\end{tikzpicture}

H�hensatz: $a \mal b = h^2$
\begin{itemize}[label=\Ra]
\item $\underbrace{h = \sqrt{ab}}_{\text{geometrisches Mittel}}$\\
		$r = \frac{1}{2} (a + b)$
\item $h \klgl r$ mit Gleichheit, falls $a = b$
\end{itemize}

\begin{fbeweis}
$a, b > 0$
%MA-25.03.2009-FORM1
\begin{align*}
\sqrt{ab} &\stackrel{\klgl}{!} \frac{1}{2} (a + b)~~~~\vert \space^2\\
\Lra ab &\stackrel{\klgl}{!} \frac{1}{4} (a + b)^2\\
\Lra ab &\stackrel{\klgl}{!} \frac{1}{4} \left(a ^2 + 2ab + b^2\right)~~~~\vert \mal 4\\
\Lra 4ab &\stackrel{\klgl}{!} a^2 + 2ab + b^2~~~~\vert -4ab\\
\Lra 0 &\stackrel{\klgl}{!} a^2 - 2ab + b^2\\
\Lra 0 \klgl (a - b)^2
\end{align*}
\begin{description}
\item[Fall] $a = b \Lra (a - b)^2 = 0 \Lra \dots \Lra \sqrt{ab} = \frac{1}{2} (a + b)$
\end{description}
\qed
\end{fbeweis}

\begin{bemerkung}
\mbox{}\par
\begin{itemize}
\item arithmetisches Mittel wird am meisten verwendet\\
		\ac{z.B.} Gewicht eines ''Standardpassagiers'' beim Fliegen

\item geometrisches Mittel wird verwendet, wenn gro�e Ausrei�er keine zu gro�e Rolle spielen sollen\\
		\ac{z.B.} gewisse Armutma�e basieren auf geometrischem Mittel
\end{itemize}
\end{bemerkung}

\begin{beispiel}
$a_1 = 1$, $a_2 = 1$, $a_3 = 1$, $a_4 = 10000$
\begin{align*}
	&\text{aritmetisches: }\frac{1}{4} \left(a_1 + a_2 + a_3 + a_4\right) = \frac{10003}{4} \approx 2500\\
	&\text{geometrisches: }\sqrt[4]{a_1 \mal a_2 \mal a_3 \mal a_4} = 10
\end{align*}
\end{beispiel}

\begin{fsatz}[Bernoulli-Ungleichungen]
%MARK: Satz 1.10
F�r Alle $x \grgl -1$ und alle $n \in \mb{N}$ gilt:
\[(1 + x)^n \grgl 1 + n \mal x\]
\end{fsatz}

\begin{fbeweis}
(f�r $x \grgl 0$ folgt das aus binomischem Lehrsatz)
\begin{description}
\item[(IA)] $n = 0$\\
		\[\left. \begin{array}{l}\text{LS} = (1 + x)^0 = 1\\\text{RS} = 1 + 0 \mal x = 1\end{array}\right\}\checkmark\]

\item[(IS)] \ac{z.z.} $\forall n \in \mb{N}_0$: $\underbrace{(1 + x)^n \grgl 1 + n \mal x}_{\text{(IV)}}$ \Ra $(1 + x)^{n + 1} \grgl 1 + (n + 1) \mal x$
		%MA-25.03.2009-FORM2
		\begin{align*}
		(1 + x)^{n + 1} &= (1 + x) \mal (1 + x)^n\\
		&\stackrel{\grgl}{\text{(IV)}} (1 + x) (1 + n \mal x)\\
		&= 1 +nx + x + n \mal x^2 = 1 + (n + 1) \mal x + \underbrace{\underbrace{n}_{\grgl 0} \mal \underbrace{x^2}_{\grgl 0}}_{\grgl 0}\\
		&\grgl 1 + (n + 1) \mal x
		\end{align*}
\end{description}
\end{fbeweis}

\begin{fdefinition}
F�r $x \in \mb{R}$ definiert man den (Absolut-) Betrag durch
\[\betrag{x} = \begin{cases}x \text{ f�r } x \grgl 0\\x \text{ f�r } x < 0\end{cases}\]
\end{fdefinition}

\begin{bemerkung}Es gilt auch
		\begin{itemize}
		\item $\betrag{x} = \max(x, -x)$
		\item $\betrag{x} = \sqrt{x^2}$ entspricht Norm f�r Vektoren der Dimension 1
		\end{itemize}
\end{bemerkung}

\begin{fsatz}[Rechnen mit Betr�gen]
%MARK: Satz 1.12
$a \in \mb{R}$
\begin{enumerate}[label=\roman*)]
\item $\betrag{a} \grgl 0$; $\betrag{a} = 0 \Lra a = 0$
\item $\betrag{a \mal b} = \betrag{a} \mal \betrag{b}$
\item $\betrag{a + b} \klgl \betrag{a} + \betrag{b}$ $\Delta$-Ungleichung
\item $\betrag{a - b} \grgl \betrag{\betrag{a} - \betrag{b}}$
\item $\betrag{a} \klgl b \Lra -b \klgl a \klgl b$
\item $\betrag{x - x_0} \klgl r \Lra x_0 - r x \klgl x_0 + r$
\end{enumerate}
\end{fsatz}

\begin{fbeweis}
\mbox{}\par
\begin{enumerate}
\item[\rM{1}), \rM{2})] folgt aus \ac{Def.} 1.11, \ac{bzw.} aus Mathe 1, Satz 6.2, 6.4
\item[\rM{3})]
		\begin{enumerate}[label=\textcircled{\arabic*}]
		\item $\begin{array}{c}
				a \klgl \betrag{a}\\
				b \klgl \betrag{b}\\
				\hline
				a + b \klgl \betrag{a} + \betrag{b}
				\end{array}$
		\item $\begin{array}{c}
				-a \klgl \betrag{a}\\
				-b \klgl \betrag{b}\\
				\hline
				-a - b \klgl \betrag{a} + \betrag{b}
				\end{array}$
		\end{enumerate}
		\[\betrag{a + b} = \max (a + b, a - b) \underbrace{\klgl}_{\text{\textcircled{1}, \textcircled{2}}} \betrag{a} + \betrag{b}\]

\item[\rM{4})]
		\begin{enumerate}[label=\textcircled{\arabic*}]
		\item $\betrag{a} = \betrag{( a - b) + b} \underbrace{\klgl}_{\Delta \text{- Ungl.}} \betrag{a - b} + \betrag{b}$\\
				$\Lra \betrag{a - b} \grgl \betrag{a} - \betrag{b}$

		\item $\betrag{b} = \betrag{b - a + a} \klgl \betrag{b - a} + \betrag{a}$\\
				$\Lra \betrag{b - a} \grgl \betrag{b} - \betrag{a}$\\
				$\Lra \betrag{a - b} \grgl \betrag{b} - \betrag{a}$\\
				\begin{align*}
				\Ra \betrag{\betrag{a} - \betrag{b}} &= \max\left(\betrag{a} - \betrag{b}, \betrag{b} - \betrag{a}\right)\\
				&\underbrace{\klgl}_{\text{\textcircled{1}, \textcircled{2}}} \betrag{a - b}
				\end{align*}
		\end{enumerate}
\item[\rM{5})] $\betrag{a} \klgl b$
		\begin{align*}
		\Lra &\max (a, -a) \klgl b\\
		\Lra &a \klgl b \und -a \klgl b\\
		\Lra &a \klgl b \und a \grgl -b\\
		\Lra &-b \klgl a \klgl b
		\end{align*}

\item[\rM{6})] $\betrag{x - x_0} \klgl r$
		\begin{align*}
		\stackrel{V}{\Lra} &-r \klgl x - x_0 \klgl r\\
		\Lra &x_0 - r \klgl x \klgl x_0 + r
		\end{align*}
\end{enumerate}
\qed
\end{fbeweis}

\begin{bemerkung}
\mbox{}\par
\begin{description}
\item[zu iii)] $n$ als $\Delta$-Ungleichung bezeichnet, da Spezialfall von $\|a + b\| \klgl \|a\| + \|b\|$

		%MA-25.03.2009-IMG4
		\begin{tikzpicture}
			\draw[serif cm->] (0,0) -- (3,0) node[below, midway]{$\vec{a}$};
			\draw[serif cm->] (3,0) -- (4,2.5) node[right, midway]{$\vec{b}$};
			\draw[serif cm->] (0,0) -- (4,2.5) node[midway, left=1em]{$\vec{a} + \vec{b}$};
		\end{tikzpicture}

		w�re die Ungleichung verletzt, k�nnten die 3 Vektoren $\vec{a}, \vec{b}$, $\vec{a} + \vec{b}$ kein Dreieck bilden.

\item[zu vi)] $x_0$ interpretiert als fester Punkt (\ra Mittelpunkt und $r$ als Radius)

		%MA-25.03.2009-IMG5
		\begin{tikzpicture}[decoration=brace]
		\draw[thick] (-1,0) -- (4,0);
		\draw[->, thick] (4.2,0) -- (5,0) node[right] {$x$};
		\draw (0.05,0.1) -- (0,0.1) -- (0,-0.1) -- (0.05,-0.1);
		\draw (3.95,0.1) -- (4,0.1) -- (4,-0.1) -- (3.95,-0.1);
		\draw (2,0.2) -- (2,-0.2) node[below] {$x_0$};
		\draw[dashed] (0,0) arc (180:360:2);
		\draw[decorate] (2,0.5) -- (4,0.5) node[above,midway] {$r$};
		\end{tikzpicture}

		$\betrag{x - x_0}$\\
		die Menge aller Punkte die h�chstens um $r$ von $x$ entfernt sind.
\end{description}
\end{bemerkung}

\section{Aufgabe 1.5}
\label{sec:Ungleichungen_A1_5}
L�sen sie die Ungleichung
\[\betrag{2 - \frac{x}{2}} < \frac{1}{2}\]

L�sung siehe \vref{sec:Ungleichungen_A1_5L}.

\section{Aufgabe 1.6}
\label{sec:Ungleichungen_A1_6}
Bestimmen sie alle L�sungen $x$ der Gleichung
\[\sqrt{(x^2 - 5x + 6)^2} = x^2 - 5x + 6\]

L�sung siehe \vref{sec:Ungleichungen_A1_6L}.

\section{L�sungen}
\subsection{Aufgabe 1.3}
\label{sec:Ungleichungen_A1_3L}
L�sung zu Aufgabe \vref{sec:Ungleichungen_A1_3}.

\begin{description}
\item[Fall 1:] $x > 0$
		\begin{align*}
		&xy > x~~~~\vert \mal x\\
		\Lra &y > 1
		\end{align*}

\item[Fall 2:] $x = 0$
		\[0 \mal y > 0 \blitza\]
		
\item[Fall 3:] $x < 0$
		\begin{align*}
		&xy > x~~~~\vert \mal x\\
		\Lra &y < 1
		\end{align*}
\end{description}
\[(x, y) \in \left((0, \infty) \times (1, \infty)\right) \cup \left((-\infty, 0) \times (-\infty, 1)\right)\]

\subsection{Aufgabe 1.4}
\label{sec:Ungleichungen_A1_4L}
L�sung zu Aufgabe \vref{sec:Ungleichungen_A1_4}.

$x, y >0$:
\begin{align*}
&\frac{x}{y} + \frac{y}{x} \stackrel{!}{\grgl} 2~~~~\vert \mal xy\\
\Lra &x^2 + y^2 \stackrel{!}{\grgl} 2xy~~~~\vert -2xy\\
\Lra &x^2 - 2xy + y^2 \stackrel{!}{\grgl} 0\\
\Lra & (x - y)^2 \grgl 0
\end{align*}

\subsection{Aufgabe 1.5}
\label{sec:Ungleichungen_A1_5L}
L�sung zu Aufgabe \vref{sec:Ungleichungen_A1_5}.

$\betrag{2 - \frac{x}{2}} < \frac{1}{2}$
\begin{align*}
\Lra -\frac{1}{2} < 2 - \frac{x}{2} < \frac{1}{2}~~~~\vert -2\\
\Lra -\frac{5}{2} < -\frac{x}{2} < - \frac{3}{2}~~~~\vert \mal (-2)\\
\Lra 5 > x > 3\\
\Lra L = (3,5)
\end{align*}

\subsection{Aufgabe 1.6}
\label{sec:Ungleichungen_A1_6L}
L�sung zu Aufgabe \vref{sec:Ungleichungen_A1_6}.

$\sqrt{(x^2 - 5x + 6)^2} = x^2 - 5x + 6$
\begin{align*}
\Lra \betrag{x^2 - 5x + 6} = x^2 - 5x + 6\\
\Lra x^2 -5x + 6 \grgl 0\\
\Lra (x - 3) \mal (x - 2) \grgl 0\\
\end{align*}

\begin{description}
\item[$x = 0$] $0^2 - 5 \mal 0 + 6 > 0$\\
		\Ra Parabel nach oben ge�ffnet\\
		$L = \mb{R} \backslash (2, 3) = (- \infty, 2] \cup [3, \infty)$
\end{description}

%
%\part{Folgen und Reihen}
%\Mark{Part 2}
%\chapter{Konvergente Folgen}
%Chapter: 3
\begin{fdefinition}[Folge]
Eine Folge ist eine Abbildung von
\[a: \mb{N} \ra \mb{R}\]
\[n \mapsto a(n) = a_n\]
Man schreibt $(a_1, a_2, a_3, \dots)$ oder $(a_n)_{n \in \mb{N}}$ oder $(a_n)$
\end{fdefinition}

\section{Spezielle Folgen}
\begin{itemize}
\item arithmetische Folge: $a_n = a_0 + n \mal d$ $d \in \mb{R}$, also $(a_0, a_0 + d, a_0 + 2d, a_0 + 3d, \dots)$, \ac{d.h.} die Folge der Differenzen $(a_1 - a_0, a_2 - a_1, a_3 - a_2, \dots) = (d, d, d, d, \dots)$ ist konstant.
		\begin{beispiel}
		\mbox{}\par
		\begin{itemize}
		\item $(1, 2, 3, 4, \dots)$ \Ra $a_0 = 1$, $d = 1$
		\item $(1, 3, 5, 7, \dots)$ \Ra $a_0 = 1$, $d = 2$
		\end{itemize}
		\end{beispiel}

\item arithmetische Folge 2. Ordnung: $b_n = b_o + \frac{n \mal (n - 1)}{2} \mal d$ die Folge der Differenzen $(b_1 - b_0, b_2 - b_2, \dots)$ ist eine arithmetische Folge (1. Ordnung).
		\begin{beispiel}
		$(1, 4, 9, 16, 25, 36, 49, 64, 81, 100, \dots)$\\
		\Ra Folge der Differenzen $(3, 5, 7, 9, 11, \dots)$
		\end{beispiel}

\item geometrische Folge: $c_n = c_0 \mal q^n$ $q \in \mb{R}$, also $(c_0, c_0 \mal q, c_0 \mal q^2, c_0 \mal q^3, \dots)$, daher die Folge der Quotienten $\left(\frac{c_1}{c_0}, \frac{c_2}{c_1}, \frac{c_3}{c_2}, \dots\right) = (q, q, q, \dots)$ ist konstant.
		\begin{beispiel}
		$(1, 2, 4, 8, 16, 32, \dots)$ \Ra $c_0 = 1$, $q = 2$
		\end{beispiel}
\end{itemize}
\begin{bemerkung}
\mbox{}\par
\begin{enumerate}
\item $a_0 = a$, $a_1 = \frac{a + b}{2}$, $a_2 = b$ bilden den Anfang einer arithmetischen Folge
		\begin{align*}
		d &= a_1 - a_0 = \frac{a + b}{2} - a = \frac{a + b - 2a}{2} = \frac{b - a}{2}\\
		&= a_2 - a_1 = b - \frac{a + b}{2} = \frac{2b - a - b}{2} = \frac{b - a}{2}
		\end{align*}

\item $c_0 = a$, $c_1 = \sqrt{ab}$, $c_2 = b$ bilden den Anfang einer geometrischen Folge
		\begin{align*}
		q &= \frac{c_1}{c_0} = \frac{\sqrt{ab}}{a} = \sqrt{\frac{ab}{a^2}} = \sqrt{\frac{b}{a}}\\
		&= \frac{c_2}{c_1} = \frac{b}{\sqrt{ab}} = \sqrt{\frac{b^2}{ab}} = \sqrt{\frac{b}{a}}
		\end{align*}
\end{enumerate}
\end{bemerkung}

\begin{fdefinition}[Eigenschaften von Folgen]
Eine Folge $(a_n)_{n \in \mb{N}}$ hei�t
\begin{enumerate}
\item monoton wachsend, falls $\forall n \in \mb{N}: a_{n + 1} \grgl a_n$
\item monoton fallend, falls $\forall n \in \mb{N}: a_{n + 1} \klgl a_n$
\item streng monoton wachsend, falls $\forall n \in \mb{N}: a_{n + 1} > a_n$
\item streng monoton fallend, falls $\forall n \in \mb{N}: a_{n + 1} < a_n$
\item von oben (oder nach oben) beschr�nkt, falls $\exists M \in \mb{R}: \forall n \in \mb{N}: a_n \klgl M$
\item von unten (oder nach unten) beschr�nkt, falls $\exists m \in \mb{R}: \forall n \in \mb{N}: a_n \grgl m$
\item beschr�nkt, falls sie nach oben und nach unten beschr�nkt ist, daher es existiert $K > 0$ so dass $\betrag{a_n} \klgl K$
\end{enumerate}
\end{fdefinition}

\begin{bemerkung}
In einigen B�chern \ac{etc.} wird verwendet
\begin{itemize}
\item statt monoton wachsend, monoton nichtfallend
\item statt monoton fallend, monoton nichtwachsend
\item statt streng monoton wachsend, monoton wachsend
\item statt streng monoton fallend, monoton fallend
\end{itemize}
\end{bemerkung}

\begin{beispiel}
\mbox{}\par
\begin{enumerate}
\item arithmetische Folge $(a_n)_{n \in \mb{N}_0}$ mit $d > 0$
		\begin{itemize}
		\item $(a_n)$ ist streng monoton wachend
		\item $(a_n)$ ist nach unten beschr�nkt ($m = a_0$)
		\end{itemize}

\item geometrische Folge $(c_n)_{n \in \mb{N}_0}$ mit $q > 1$
		\begin{itemize}
		\item $(c_n)$ ist streng monoton wachsend
		\item $(c_n)$ ist nach unten beschr�nkt ($m = c_0$)
		\end{itemize}

\item geometrische Folge $(c_n)_{n \in \mb{N}_0}$ mit $0 < q < 1$, $c_0 > 0$
		\begin{itemize}
		\item $(c_n)$ ist streng monoton fallend
		\item $(c_n)$ ist beschr�nkt ($M = c_0$, $m = 0$ \Ra $k = c_0$)
		\end{itemize}
\end{enumerate}
\end{beispiel}

\begin{fdefinition}[Konvergente Folge]
Eine Folge $(a_n)_{n \in \mb{N}}$ hei�t \indexu{konvergent} gegen den Grenzwert $a \in \mb{R}$, falls $\forall \epsilon > 0 \exists N(\epsilon)$, so dass $\forall n \grgl N(\epsilon)$:
\[\betrag{a_n - a} < \epsilon\]
In diesem Fall schreibt man
\[\lim_{n \ra \infty} a_n = a\]
Ist eine Folge nicht konvergent nennt man sie \indexu{divergent}.
\end{fdefinition}

\begin{beispiel}
\mbox{}\par
\begin{enumerate}
\item $a_n \coloneqq \frac{1}{n}$, \ac{d.h.} mit $N(\epsilon) = \lceil \frac{1}{\epsilon} + 1 \rceil$ gilt
		\[\betrag{{a_n} - \underbrace{a}_{= 0}} = \betrag{\frac{1}{n} - 0} = \frac{1}{n} < \epsilon \text{ f�r } n \grgl N(\epsilon)\]
\item geometrische Folge mit $c_0 > 0$, $0 < q < 1$, \ac{d.h.} mit $N(\epsilon) = \lceil \frac{\ln(\epsilon)}{\ln(q)} + 1\rceil$ gilt:
		\[\betrag{c_n - 0} = \betrag{c_0 \mal q^n - 0} = c_0 \mal q^n < \epsilon\]
\end{enumerate}
\end{beispiel}

\begin{fsatz}[Rechenregeln f�r konvergente Folgen]
Seien $(a_n)_{m \in \mb{N}}$ und $(b_m)_{n \in \mb{N}}$ konvergente Folgen mit $\lim_{n \ra \infty} a_n = a$ und $\lim_{n \ra \infty} b_n = b$, dann gilt:
\begin{enumerate}[label=\roman*)]
\item $\lim_{n \ra \infty} (a_n \pm b_n) ) a \pm b$
\item $\lim_{n \ra \infty} (a_n \mal b_n) = a \mal b$
\item falls $a_n < b_n \forall n \in \mb{N} \Ra a \klgl b$
\item falls $b_n \neq 0 \forall n \in \mb{N} \Ra \lim_{n \ra \infty} \frac{a_n}{b_n} = \frac{a}{b}$
\end{enumerate}
\end{fsatz}

\begin{fbeweis}
\mbox{}\par
\begin{enumerate}[label=\roman*)]
\item f�r ''$+$''\\
		wir wissen 
		\[\lim_{n \ra \infty} a_n = a \Ra \forall \frac{\epsilon}{2} > 0 \exists N_1 \left(\frac{\epsilon}{2}\right): \forall n \grgl N_1 \left(\frac{\epsilon}{2}\right): \betrag{a_n - a} < \frac{\epsilon}{2}\]
		und
		\[\lim_{n \ra \infty} b_n = b \Ra \forall \frac{\epsilon}{2} > 0 \exists N_2 \left(\frac{\epsilon}{2}\right): \forall n \grgl N_2 \left(\frac{\epsilon}{2}\right): \betrag{b_n - b} < \frac{\epsilon}{2}\]
		Sei nun $n \grgl \max\left(N_1 \left(\frac{\epsilon}{2}\right), N_2 \left(\frac{\epsilon}{2}\right)\right)$\\
		\[\betrag{(a_n + b_n) - (a + b)} = \betrag{(a_n - a) + (b_n - b)} \underbrace{\klgl}_{\Delta\text{-Ung.}} \betrag{a_n - a} + \betrag{b_n - b} < \frac{\epsilon}{2} + \frac{\epsilon}{2} = \epsilon\]
\end{enumerate}a
\end{fbeweis}

\begin{bemerkung}[zu Satz 2.4 iii)]
$a_n \ceq \frac{1}{n^2 + 1}$, $b_n \ceq \frac{1}{n}$ $n \in \mb{N}$
\[a_n < b_n \forall n \in \mb{N}\]
$a = \lim_{n \ra \infty} (a_n) = 0$, $b = \lim_{n \ra \infty} b_n) = 0$
\end{bemerkung}

\begin{bemerkung}
Wenn $f$ eine Funktion ohne Sprungstellen ist, (\ra $f$ stetig) und $(a_n)_{n \in \mb{N}}$ eine konvergente Folge mit $\lim_{n \ra \infty} a_n = a$ ist, gilt
\[\lim_{n \ra \infty} f(a_n) = f\left(\lim_{n \ra \infty} a_n\right) = f(a)\]
\begin{beispiel}
$f(x) = x^3$, 
$f(x) = e^x$, 
$f(x) = \sinx{x}$, 
$f(x) = \sqrt[5]{x}$
\end{beispiel}
\end{bemerkung}

\begin{beispiel}
\mbox{}\par
\begin{align*}
&\lim_{n \ra \infty} \left(\frac{n^2}{n - 1} - \frac{n^2 + 1}{n}\right) = \lim_{n \ra \infty} \left(\frac{n^3 - (n - 1) \mal (n^2 + 1)}{(n - 1) \mal n}\right)\\
= &\lim_{n \ra \infty} \frac{n^3 - (n^3 + n - n ^2 - 1)}{(n - 1) \mal n} = \lim_{n \ra \infty} \frac{-n + n^2 + 1}{(n - 1) \mal n}\\
= &\lim_{n \ra \infty} \frac{n^2 - n + 1}{n^2 - n} = \lim_{n \ra \infty} \frac{1 - \frac{1}{n} + \frac{1}{n^2}}{1 - \frac{1}{n}}\\
= &\frac{\lim_{n \ra \infty}\left(1 - \frac{1}{n} + \frac{1}{n^2}\right)}{\lim_{n \ra \infty}\left(1 - \frac{1}{n}\right)} = \frac{1}{1} = 1
\end{align*}
\end{beispiel}

\section{Aufgabe 2.1}
\label{sec:KonvergenteFolgen_A2_1}
Bestimmen sie die folgenden Grenzwerte:
\begin{enumerate}[label=\alph*)]
\item $\lim_{n \ra \infty} \frac{2}{3} \mal \left(\frac{1}{3^n} + \frac{n^2 + 1}{n^2 + 2n}\right)$
\item $\lim_{n \ra \infty} \frac{(n - 2)^3}{(2n + 1)^3}$
\end{enumerate}

L�sung siehe \vref{sec:KonvergenteFolgen_A2_1L}.

\begin{fsatz}[Hinreichendes Konvergenzkriterium]
%MARK: Satz 2.5
Jede monotone und beschr�nkte Folge ist konvergent.
\begin{beweis}
-
\end{beweis}
\end{fsatz}

\begin{bemerkung}
\mbox{}\par
\begin{enumerate}
\item Die Beschr�nktheit alleine ist eine notwendige Bedingung
		\begin{itemize}[label=\Lra]
		\item Jede Folge die nicht beschr�nkt ist, ist auch nicht konvergent.
		\item Jede konvergente Folge ist beschr�nkt.
		\end{itemize}
\item Satz 2.5 ist gleichbedeutend mit der Vollst�ndigkeit der reellen Zahlen (\ac{d.h.} es existieren keine L�cken)
		\[a_n = \sqrt{2} - \frac{1}{n}\]
		beschr�nkt mit $M = \sqrt{2}$, $m = 0$ (streng) monoton wachsend.

		%MA-27.03.2009-IMG-1
		\begin{tikzpicture}
		\draw[<->, thick] (0,4.5) -- (0,0) -- (6.5,0) node[below] {$n$};
		\draw (-0.2,3) -- (6,3) node[right] {$\sqrt{2}$};
		\draw (0.7,0.2) parabola[bend pos=0.3] bend(5.5,2.9) (6,2.9);
		\end{tikzpicture}

		\[\lim_{n \ra \infty} a_n = \sqrt{2} \notin \mb{Q}\]
\end{enumerate}
\end{bemerkung}

\begin{beispiel}[Konvergente Folgen]
\mbox{}\par
\begin{enumerate}
\item $\lim_{n \ra \infty} q^n = \begin{cases}1 \text{ f�r } q = 1\\0 \text{ f�r } \betrag{q} < 1\end{cases}$
\item $\lim_{n \ra \infty} \frac{q^n}{n!} = 0 \text{ f�r } q \in \mb{R}$
\item $\lim_{n \ra \infty} \frac{1}{n^{\alpha}} = \begin{cases}1 \text{ f�r } \alpha = 0\\0 \text{ f�r } \alpha > 0\end{cases}$
\item $\lim_{n \ra \infty} \frac{n^{\alpha}}{n!} = 0 \text{ f�r } \alpha \in \mb{R}$
\item $\lim_{n \ra \infty} \sqrt[n]{q} = 1 \text{ f�r } q > 0$
\item $\lim_{n \ra \infty} \sqrt[n]{n} = 1$
\item $\lim_{n \ra \infty} \rklamm{1 + \frac{1}{n}}^n = e \approx 2,718$
\end{enumerate}
\end{beispiel}

\begin{fbeweis}
\mbox{}\par
\begin{enumerate}
\item bereits erledigt
\item $\lim_{n \ra \infty} \frac{q^n}{n!}$
		\begin{align*}
		&\lim_{n \ra \infty} \rklamm{\underbrace{\frac{q}{n}}_{> 1} \mal \underbrace{\frac{q}{n - 1}}_{> 1} \mal \underbrace{\frac{q}{n - 2}}_{> 1} \mal \dots \mal \underbrace{\frac{q}{m + 1}}_{< 1} \mal \underbrace{\frac{q}{m}}_{> 1} \mal \dots \mal \underbrace{\frac{q}{2}}_{> 1} \mal \underbrace{\frac{q}{1}}_{> 1}}\\
		&= \frac{q^m}{m!} \mal \lim_{n \ra \infty} \rklamm{\frac{q}{n} \mal \frac{q}{n + 1} \mal \dots \mal \underbrace{\frac{q}{m + 1}}_{\tx{der gr��te der Faktoren}}} \klgl \frac{q^m}{m!} \mal \lim_{n \ra \infty} \rklamm{\frac{q}{m + 1}}^{n - m}\\
		&= \underbrace{\frac{q^m}{m!} \mal \rklamm{\frac{m + 1}{q}^m}}_{\tx{enth�lt}} \mal \underbrace{\lim_{n \ra \infty} \underbrace{\rklamm{\frac{q}{m + 1}}^n}_{< 1}}_{= 0} = 0
		\end{align*}
\item -
\item -
\item $\lim_{n \ra \infty} \sqrt[n]{q}$
		\[\lim_{n \ra \infty} q^{\frac{1}{n}} = q^{\lim_{n \ra \infty} \rklamm{\frac{1}{n}}} = q^0 = 1\]

\item $\lim_{n \ra \infty} \sqrt[n]{n}$
		\begin{align*}
		&\lim_{n \ra \infty} n^{\frac{1}{n} = \lim_{n \ra \infty} e^{\frac{1}{n} \ln (n)}}\\
		&= e^{\underbrace{\lim_{n \ra \infty} \frac{1}{n} \mal \ln (n)}_{= 0}} = e^0 = 1
		\end{align*}

\item -
\end{enumerate}
\end{fbeweis}

\begin{beispiel}
$\lim_{n \ra \infty} \rklamm{\sqrt[n]{3^n \mal \rklamm{n^2 + 1}}}$
\begin{align*}
&\lim_{n \ra \infty} \rklamm{\underbrace{\sqrt[n]{3^n}}_{= 3} \mal \sqrt[n]{n^2 + 1}}\\
&= 3 \mal \lim_{n \ra \infty} \sqrt[n]{n^2 + 1}\\
\Ra & \grgl 3 \mal \lim_{n \ra \infty} \underbrace{\sqrt[n]{n^2}}_{n^{\frac{2}{n}}} = 3 \mal \lim_{n \ra \infty} \rklamm{\sqrt[n]{n}}^2 = 3 \mal \rklamm{\lim_{n \ra \infty} \sqrt[n]{n}}^2 = 3\\
\Ra & \underbrace{\klgl}_{n \grgl 2} 3 \mal \lim_{n \ra \infty} \sqrt[n]{n^3} = 3 \rklamm{\lim_{n \ra \infty} \sqrt[n]{n}}^3 = 3\\
& \Ra \lim_{n \ra \infty} \sqrt[3]{3^n \rklamm{n^2 + 1}} = 3
\end{align*}
\end{beispiel}

\begin{definition}[Nullfolge]
%MARK: Definition 2.6
Eine Folge $(a_n)_{n \in \mb{N}}$ nennt man \indexu{Nullfolge}, wenn $\lim_{n \ra \infty} a_n = 0$
\end{definition}

\section{L�sungen}
\subsection{Aufgabe 2.1}
\label{sec:KonvergenteFolgen_A2_1L}
L�sung zu Aufgabe \vref{sec:KonvergenteFolgen_A2_1}.

\begin{enumerate}[label=\alph*)]
\item $\lim_{n \ra \infty} \frac{2}{3} \mal \left(\frac{1}{3^n} + \frac{n^2 + 1}{n^2 + 2n}\right)$
		\begin{align*}
		&\frac{2}{3} \mal \lim_{n \ra \infty} \left(\frac{1}{3^n} + \frac{n^2 + 1}{n^2 + 2n}\right)\\
		=&\frac{2}{3} \mal \lim_{n \ra \infty} \frac{1}{3^n} + \lim_{n \ra \infty} \frac{n^2 + 1}{n^2 + 2n}\\
		=&\frac{2}{3} \left(0 + \lim_{n \ra \infty} \frac{1 + \frac{1}{n^2}}{1 + \frac{2}{n}}\right) = \frac{2}{3}
		\end{align*}

\item $\lim_{n \ra \infty} \frac{(n - 2)^3}{(2n + 1)^3}$
		\begin{align*}
		&\lim_{n \ra \infty} \left(\frac{n - 2}{2n + 1}\right)^3\\
		=&\left(\lim_{n \ra \infty} \left(\frac{n - 2}{2n + 1}\right)\right)^3 = \left(\lim_{n \ra \infty} \frac{1 - \frac{2}{n}}{2 + \frac{1}{n}}\right)^3\\
		=&\rklamm{\frac{\lim_{n \ra \infty} \rklamm{1 - \frac{2}{n}}}{\lim_{n \ra \infty}}} \rklamm{2 + \frac{1}{n}}^3 = \left(\frac{1}{2}\right)^3 = \frac{1}{8}
		\end{align*}
\end{enumerate}

%\chapter{Landau-Symbole}
%Chapter: 4
Die Landau-Symbole (\ac{bzw.} Landausche Symbole) beschreiben das Wachstum von Folgen (\ac{bzw.} auch Funktionen) f�r $n \ra \infty$ (\ac{bzw.} $x \ra x$)\\
Sie werden beispielsweise bei der Laufzeitanalyse von Algorithmen eingesetzt. Die Laufzeit eines Algorithmuses  kann man als Folge $(a_n)_{n \in \mb{N}}$ betrachten, wobei der Index $n$ f�r die Gr��e der Eingabe steht (oft in Anzahl bit).

\begin{fdefinition}[Landau-Symbole]
%MARK: Definition 2.7
Es seien zwei Folgen $(a_n)_{n \in \mb{N}}$ und $(b_n)_{n \in \mb{N}}$ gegeben
\begin{enumerate}
\item Die Folge $(a_n)_{n \in \mb{N}}$ ist "`gro�-O"' von $(b_n)_{n \in \mb{N}}$ (geschrieben $a_n = 0 (b_n)$), selten auch $a_n \in O (b_n)$ falls die Quotientenfolge $\rklamm{\frac{a_n}{b_n}}_{n \in \mb{N}}$ beschr�nkt ist (\ac{d.h.} es $\exists K > 0: \forall n \in \mb{N}$ $\betrag{\frac{a_n}{b_n}} \klgl K$)

		\textit{Anschaulich:} $(a_n)_{n \in \mb{N}}$ w�chst h�chstens so schnell wie $(b_n)_{n \in \mb{N}}$
\item Die Folge $(a_n)_{n \in \mb{N}}$ ist "`klein o"' von $(b_n)_{n \in \mb{N}}$ (geschrieben $a_n = o(b_n)$, selten auch $a_n \in o(b_n)$), falls die Quotientenfolge $\rklamm{\frac{a_n}{b_n}}_{b \in \mb{N}}$ eine Nullfolge ist.

		\textit{Anschaulich:} $(a_n)_{n \in \mb{N}}$ w�chst langsamer als $(b_n)_{n \in \mb{N}}$
\end{enumerate}
\end{fdefinition}

\begin{beispiel}
\mbox{}\par
\begin{enumerate}
\item $a_n = 7 \mal n^4 - 8 n^3 + 1738n - 17$
		\begin{multicols}{2}
		\begin{itemize}
		\item $a_n = O(7n^5)$
		\item $\underbrace{a_n = O(n^4)}_{\text{bestm�gliche}}$
		\item $a_n = o(n^5)$
		\item $a_n = o(n^{4,1})$
		\item $a_n = o(n^4 \log(n))$
		\end{itemize}
		\end{multicols}
\item $a_n = n^3 + (\log(n))^7$\\
		$a_n = O(n^3)$

		$a_n = n^3 + 2^{\frac{n}{100}}$\\
		$a_n = O \rklamm{2^{\frac{n}{100}}}$
\end{enumerate}
\end{beispiel}

\begin{bemerkung}
\mbox{}\par
\begin{itemize}
\item $a_n = o(b_n) \Ra a_n = O(b_n)$
\item $a_n = O\rklamm{n^k} \Ra a_n = o\rklamm{n^{k + \epsilon}}$ f�r $\epsilon > 0$ beliebig
\end{itemize}
\end{bemerkung}

\begin{bemerkung}[F�r Landau-Symbole gilt]
\mbox{}\par
\begin{itemize}
\item Term niedrigerer Ordnung sind unwichtig
\item Konstante Faktoren sind unwichtig
\end{itemize}
\end{bemerkung}

\begin{achtung}
Landau-Symbole machen Aussagen f�r $n \ra \infty$, \ac{d.h.} f�r endliche $n$ k�nnen Vorfaktoren oder Terme niedrigerer Ordnung durchaus eine Rolle spielen. Zum Beispiel:
\begin{align*}
a_n &= 10^{100} n = O(n)\\
a_n &= n + 1000 \sqrt{n}
\end{align*}
\end{achtung}

\begin{bemerkung}
\mbox{}\par
\begin{enumerate}
\item \begin{itemize}
		\item Man nennt einen Algorithmus \indexu{effizient} wenn seine Komplexit�t/Laufzeit im "`worst case"' h�chstens polynomial w�chst.
		\item Ist die Komplexit�t im "`worst case"' exponentiell nennt man den Algorithmus \indexu{ineffizient}
		\item Es werden in der Praxis auch ineffiziente Algorithmen verwendet, weil \ac{z.B.} kein effizienter Algorithmus existiert oder der "`worst case"' nicht oder nur selten auftritt.
		\end{itemize}

\item \begin{enumerate}
		\item Fester Prozessor/Vergr��erung des Problems um Faktor 10
		\begin{table}[h]
		\centering
		\begin{tabular}{c|c}
		Komplexit�t des Algorithmuses & Verl�ngerung der Laufzeit um Faktor\\
		\hline
		$O(n)$ & 10\\
		$O(n^2)$ & 100\\
		$La = c \mal 2^n$ & $2^{9n} (n = 100 \Ra 2^{900} = 10^{270})$\\
		$Ln = c \mal 2^{10n}$ & \\
		$\frac{Ln}{La} = \frac{c \mal 2^{10n}}{c \mal 2^n} = 2^{9n}$ & \\
		\end{tabular}
		\end{table}

		\item Feste Laufzeit, um Faktor 100 besserer Prozessor
		\begin{table}[h]
		\centering
		\begin{tabular}{c|c}
		Komplexit�t des Algorithmuses & Vergr��erung des Problems um Faktor\\
		\hline
		$O(n)$ & 100\\
		$O(n^2)$ & 10\\
		$O(2^n)$ & $\frac{\lb (100)}{n} + 1 (n = 100 \Ra 1{,}06)$
		\end{tabular}
		\end{table}
		\begin{align*}
		&L_a = c \mal 2^n\\
		&L_n = c \mal 2^{xn} = 100 \mal L_a = 100 \mal c \mal 2^n\\
		\Lra &2^{xn} = 100 \mal 2^n\\
		\Lra &2^{n (x - 1)} = 100\\
		\Lra &n(x - 1) = \lb(100)\\
		\Lra &x - 1 = \frac{\lb(100)}{n}\\
		\Lra &x = \frac{\lb(100)}{n} + 1
		\end{align*}
		\end{enumerate}

\item Die Komplexit�t eines (mathematischen) Problems definiert man durch die Komplexit�t des besten (bestm�glichen) Algorithmuses.
\end{enumerate}
\end{bemerkung}

\section{Aufgabe 2.2}
\label{sec:Landau_Symbole_A2_2}
Bestimmen sie die Grenzwerte f�r $n$ gegen $\infty$ (falls diese existieren)
\begin{enumerate}[label=\alph*)]
\item $a_n \ceq \frac{\rklamm{2^n - 6}^n}{2 \mal 4^n + 1}$
\item $b_n \ceq \rklamm{1 + \frac{1}{n}}^n + \rklamm{-1 - \frac{1}{n}}^n$
\item $c_n \ceq \sqrt[n]{2^n + 3^n} \mal \rklamm{1 + \frac{1}{2n}}^{4n}$
\end{enumerate}

L�sung siehe \vref{sec:Landau_Symbole_A2_2L}.

\section{L�sungen}
\subsection{Aufgabe 2.2}
\label{sec:Landau_Symbole_A2_2L}
L�sung zu Aufgabe \vref{sec:Landau_Symbole_A2_2}.

\begin{enumerate}[label=\alph*)]
\item $a_n \ceq \frac{\rklamm{2^n - 6}^n}{2 \mal 4^n + 1}$
		\begin{align*}
		&\lim_{n \ra \infty} \frac{\rklamm{2^n}^2 - 2 \mal 2^n \mal 6 + 36}{2 \mal 4^n + 1} = \lim_{n \ra \infty} \frac{\rklamm{2^2}^n - 12 \mal 2^n + 36}{2 \mal 4^n + 1}\\
		= &\lim_{n \ra \infty} \frac{1 - 12 \rklamm{\frac{2}{4}}^n + \frac{36}{4n}}{2 + \frac{1}{4^n}} = \frac{\lim_{n \ra \infty} \rklamm{1 - 12 \rklamm{\frac{1}{2}}^n + \frac{36}{4^n}}}{\lim_{n \ra \infty} \rklamm{2 + \frac{1}{4^n}}} = \frac{1}{2}
		\end{align*}

\item $b_n \ceq \rklamm{1 + \frac{1}{n}}^n + \rklamm{-1 - \frac{1}{n}}^n$
		\begin{align*}
		&\rklamm{1 + \frac{1}{n}}^n + \rklamm{-1 - \frac{1}{n}}^n = \rklamm{1 + \frac{1}{n}}^n + (-1)^n \mal \rklamm{1 + \frac{1}{n}}^n\\
		= &\rkl{1 + \frac{1}{n}}^n \mal \rkl{1 + (-1)^n} = \begin{cases}0 \tx{ f�r } n \tx{ ungerade}\\2 \underbrace{\rkl{1 + \frac{1}{n}}^n}_{\overrightarrow{n \ra \infty} \mbox{ $e$}} \tx{ f�r } n \tx{ gerade}\end{cases}
		\end{align*}
		unterschiedliche Grenzwerte f�r gerade \ac{bzw.} ungerade $n$
		\begin{itemize}[label=\Ra]
		\item $(b_n)_{n \in \mb{N}}$ divergiert
		\end{itemize}

\item $c_n \ceq \sqrt[n]{2^n + 3^n} \mal \rklamm{1 + \frac{1}{2n}}^{4n}$
\[\underbrace{\lim_{n \ra \infty} \rkl{\sqrt[n]{2^n + 3^n}}}_{= \tx{(1)}} \mal \underbrace{\lim_{n \ra \infty} \rkl{1 + \frac{1}{2n}}^{4n}}_{= \tx{(2)}}\]
\begin{align*}
\tx{(1)} \grgl & \lim_{n \ra \infty} \sqrt[n]{3^n} = 3\\
\tx{(1)} \klgl & \lim_{n \ra \infty} \sqrt[n]{2 \mal 3^n} = \overbrace{\lim_{n \ra \infty} \rkl{\sqrt[n]{2}}}^{= 1} \mal \overbrace{\lim_{n \ra \infty} \sqrt[n]{2^n}}^{= 3} = 3
\end{align*}
\begin{align*}
\tx{(2)} = & \lim_{n \ra \infty} \rkl{\rkl{1 + \frac{1}{2n}}^{2n}}^2 = \rkl{\lim_{n \ra \infty} \rkl{1 + \frac{1}{n}}^{2n}}^2\\
= & \rkl{\lim_{\substack{m \ra \infty\\m = 2n}} \rkl{1 + \frac{1}{m}}^m}^2 = e^2\\
\Ra & \lim_{n \to \infty} c_n = 3e^2
\end{align*}
\end{enumerate}

%\chapter{Konvergente Reihen}
%Chapter: 5
\begin{fsatz}[Endliche geometrische Reihe]
F�r alle $q \in \mb{R} \backslash \gklamm{1}$ und $n \in \mb{N}_0$ gilt:
\[\sum_{k = 0}^n q^k = \frac{1 - q^{n +1}}{1 - q}\]
\end{fsatz}

\begin{fbeweis}[Vollst�ndige Induktion]
\mbox{}\par
\begin{description}
\item[(IA)] $n = 0$
		\[\left.\begin{array}{l}\tx{LS } = \sum_{k = 0}^0 q^k = q^0 = 1\\\tx{RS } = \frac{1 - q^1}{1 - q} = 1\end{array}\right\} \checkmark\]
\item[(IS)] \ac{z.z.} $\forall n \in \N_0: \underbrace{\rkl{\sum_{k = 0}^n = q^k = \frac{1 - q^{n + 1}}{1 - q}}}_{\tx{(IV)}} \Ra \sum_{k = 0}^{n + 1} q^k = \frac{1 - q^{n + 2}}{1 - q}$
\begin{align*}
&\sum_{k = 0}^{n + 1} q^k = q^{n + 1} + \sum_{k = 0}^n q^k \stack{=}{\tx{(IV)}} q^{n + 1} + \frac{1 - q^{n + 2}}{1 - q}\\
= & \frac{q^{n + 1} - q^{n + 2} + 1 - q^{n + 1}}{1 - q} = \frac{1 - q^{n + 2}}{1 - q}
\end{align*}
\end{description}\qed
\end{fbeweis}

\begin{fdefinition}[Unendliche Reihen]
%MARK: Definition 2.9
Sei die Folge $(A_k)_{k \in \mb{N}}$ gegeben. Die Folge $(S_n)_{n \in \mb{N}}$ der \indexu{Partialsummen} ist definiert durch
\[S_n \ceq \sum_{k = 1}^n A_k\]
Man nennt $(S_n)_{n \in \mb{N}}$ auch \indexu{unendliche Reihe} mit den Gliedern $A_k$ (geschrieben $\sum_{k = 1}^{\infty} A_k$ oder $\sum A_k$).\\
Wenn die Folge $(S_n)_{n \in \mb{N}}$ konvergiert, nennt man auch die Reihe $\sum_{k = 1}^{\infty} A_k$ konvergent und $S \ceq \lim_{n \ra \infty} S_n = \sum_{k = 1}^{\infty} A_k$ ist der Wert der Summe.
\end{fdefinition}

\begin{bemerkung}
Unendliche Reihen k�nnen bei einem beliebigen Index $k \in \mb{Z}$ beginnen ($\sum_{k = j}^{\infty} A_k$) meistens $j = 0$ oder $j = 1$.
\end{bemerkung}

\begin{fsatz}[Unendliche geometrische Reihe]
F�r die (unendliche) geometrische Reihe gilt:
\[\sum_{k = 0}^{\infty} q^k \text{ ist divergent f�r } \betrag{q} \grgl 1\]
\[S \ceq \sum_{k = 0}^{\infty} q^k = \frac{1}{1 - q} \tx{ f�r } \betrag{q} < 1\]
\end{fsatz}

\begin{fbeweis}
\mbox{}\par
\begin{description}
\item[Fall 1] $\betrag{q} < 1$
		\[S = \lim_{n \ra \infty} S_n = \lim_{n \ra \infty} \sum_{k = 0}^n = A_k \underbrace{=}_{\text{Satz 2.8}} \lim_{n \ra \infty} \frac{1 - q^{n + 1}}{1 - q} = \frac{1}{1 - q}\]
\item[Fall 2] $\betrag{q} > 1$
		\[S_n = \sum_{k = 0}^n A_k = \frac{1 - q^{n + 1}}{1 - q}\]
		ist nicht beschr�nkt, da $q^{n + 1}$ f�r $\betrag{q} > 1$ divergent
		\begin{itemize}[label=\Ra]
		\item $(S_n)_{n \in \mb{N}}$ divergent
		\item $\sum_{k = 0}^{\infty} A_k$ divergent
		\end{itemize}

\item[Fall 3] $q = 1$
		\[ S_n = \sum_{k = 0}^n 1^k = n + 1 \stack{n \ra \infty}{\longrightarrow} \infty\]
		\begin{itemize}[label=\Ra]
		\item $(S_n)_{n \in \mb{N}}$ divergent
		\item $\sum_{k = 0}^{\infty} A_k$ divergent
		\end{itemize}

\item[Fall 4] $q = -1$
		\[S_n = \sum_{k = 0}^n (-1)^k = \begin{cases}1 \text{ f�r } n \tx{ gerade}\\0 \tx{ f�r } n \tx{ ungerade}\end{cases}\]
		\begin{itemize}[label=\Ra]
		\item $(S_n)_{n \in \mb{N}}$ divergent
		\item $\sum_{k = 0}^{\infty} A_k$ divergent
		\end{itemize}
\end{description}\qed
\end{fbeweis}

\begin{fsatz}[Konvergenzkriterien f�r Reihen]
\mbox{}\par
\begin{enumerate}[label=\roman*)]
\item Notwendiges Kriterium\\
		\[\sum_{k = 0}^{\infty} A_k \tx{ konvergent } \Ra \lim_{k \to \infty} A_k = 0\]
		Dies ist aber kein zwingendes Kriterium f�r konvergenz eher f�r das Beweisen der divergenz geeignet.

\item Leibnitz-Kriterium (f�r alternierende Reihen)\\
		Wenn $(A_k)_{k \in \N}$ eine monoton fallende Nullfolge ist, dann konvergiert die alternierende Reihe
		\[\sum_{k = 1}^{\infty} (-1)^k \mal A_k\]

\item Vergleichskriterien\\
		Sei $\betrag{A_k} \klgl B_k \forall k \grgl k_0$ dann gilt:
		\begin{enumerate}[label=\alph*)]
		\item Majorantenkriterium\\
				$\sum_{k = 0}^{\infty} B_k$ konvergent
				\begin{itemize}[label=\Ra]
				\item $\sum_{k = 0}^{\infty} \betrag{A_k}$ konvergent
				\item $\sum_{k = 0}^{\infty} A_k$ konvergent
				\end{itemize}

		\item Minorantenkriterium\\
				$\sum_{k = 0}^{\infty} \betrag{A_k}$ divergent
				\begin{itemize}[label=\Ra]
				\item $\sum_{k = 0}^{\infty} B_k$ divergent
				\end{itemize}
		\end{enumerate}

\item Wurzelkriterium
		\[\lim_{k \ra \infty} \sqrt[k]{\betrag{A_k}} = w = \begin{cases}< 1 & \Ra \sum_{k = 0}^{\infty} A_k, \sum_{k = 0}^{\infty} \betrag{A_k} \tx{ konvergent}\\> 1 & \Ra \sum_{k = 0}^{\infty} A_k, \sum_{k = 0}^{\infty} \betrag{A_k} \tx{ divergent}\end{cases}\]

\item Quotientenkriterium
		\[\lim_{k \ra \infty} \betrag{\frac{A_k + 1}{A_k}} = q = \begin{cases}< 1 & \Ra \sum_{k = 0}^{\infty} A_k, \sum_{k = 0}^{\infty} \betrag{A_k} \tx{ konvergent}\\> 1 & \Ra \sum_{k = 0}^{\infty} A_k, \sum_{k = 0}^{\infty} \betrag{A_k} \tx{ divergent}\end{cases}\]
\end{enumerate}
\end{fsatz}

\begin{fbeweis}
\mbox{}\par
\begin{enumerate}[label=\roman*)]
\item Kontraposition \ac{z.z.}
		\[\neg \rkl{\lim_{k \ra \infty} A_k = 0} \Ra \sum A_k \tx{ divergent}\]
		$A_k$ keine Nullfolge
		\begin{align*}
		&\exists \epsilon > 0~\forall N(\epsilon)~\exists n \grgl N(\epsilon):~\betrag{A_k - 0} > \epsilon\\
		\Lra &\exists \epsilon > 0 ~ \forall \N_0 \in \N ~ \exists n \grgl \N_0:~\betrag{A_n} > \epsilon\\
		&\betrag{S_n - S_{n - 1}} = \betrag{A_n} > \epsilon\\
		&\Ra \forall S \in \mb{R}:~\max \rkl{\betrag{S - S_n}, \betrag{S - S_{n - 1}}} > \frac{\epsilon}{2}
		\end{align*}

		%MA-02.04.2009-IMG-2
		\begin{tikzpicture}[decoration=brace]
		\foreach \x in {0, 0.5, 1,..., 5}
			\draw (\x,0) rectangle +(0.3,0.2);
		
		\draw[decorate] (5.3,-0.2) -- (3.5,-0.2) node[below,midway] {m�gl. positionen};
		\draw[decorate] (4.8,-1) -- (1.5,-1) node[below,midway] {m�gl. positionen};
		\end{tikzpicture}

\item $A_{k + 1} \klgl A_k$ $S_n = \sum_{k = 1}^n (- 1)^k A_k$

		%MA-07.04.2009-IMG-1
		\begin{tikzpicture}
		\draw[->, thick] (-0.3,0) node[left] {$S_n$} -- (6.5,0) node[right] {$n$};
		\draw[<-, thick] (0,0.3) -- (0,-3);
		
		\foreach \x in {0.5, 1.5,...,5.5}
			\draw (\x,0.05) -- (\x,-0.05);
		
		\foreach \x in {1, 2,...,6}
			\draw (\x,0.05) node[above] {$\x$} -- (\x,-0.05);
		
		\draw[<->] (1,-2.75) node[below] {$S_1$} -- (2,-0.5) node[right] {$S_2$} -- (3,-2) node[below] {$S_3$} -- (4,-0.5) node[right] {$S_4$} -- (5,-1.5) node[below] {$S_5$} -- (6,-1) node[above] {$S_6$} -- (7,-1.25) node[below] {$S_7$};
		
		\draw[<->, draw=red] (1,0) -- (1,-2.75) node[left,midway,color=red] {$A_1$};
		\draw[<->, draw=red] (2,-0.5) -- (2,-2.75) node[right,midway,color=red] {$A_2$};
		\draw[<->, draw=red] (3,-0.5) -- (3,-2) node[right,near start,color=red] {$A_3$};
		\end{tikzpicture}

		\[\underbrace{S_1 \klgl S_3 \klgl S_5 \klgl \dots \klgl S_{2n + 1}}_{\begin{minipage}{35mm}\begin{center}\scriptsize{monoton wachsende $n$ nach oben beschr�nkt Folge mit $T \ceq \lim_{n \ra \infty} s_{2n + 1}$}\end{center}\end{minipage}} \klgl S \klgl \underbrace{S_{2n} \klgl \dots \klgl S_6 \klgl S_4 \klgl S_2}_{\begin{minipage}{32mm}\begin{center}\scriptsize{monoton fallende und nach unten beschr�nkte Folge mit $U \ceq \lim_{n \ra \infty} s_{2n}$}\end{center}\end{minipage}}\]
		\begin{align*}
		U - T &= \lim_{n \ra \infty} \rklamm{S_{2n}} - \lim_{n \ra \infty} \rklamm{S_{2 n + 1}} = \lim_{n \ra \infty} \rklamm{S_{2n} - S_{2n + 1}}\\
		&= \lim_{n \ra \infty} \rklamm{\sum_{k = 1}^{2n} (-1)^k A_k - \sum_{k = 1}^{2n + 1} (-1)^k A_k} = \lim_{n \ra \infty} -(-1)^{2n + 1} \mal A_{2n + 1} = \lim_{n \ra \infty} A_{2n + 1} = 0\\
		&\Ra \lim_{n \ra \infty} S_n = S(= T = U)
		\end{align*}

\item \begin{enumerate}[label=\alph*)]
		\item Majorantenkriterium
				\begin{itemize}
				\item es gilt $\sum_{k = 0}^{\infty} B_k$ konvergent und $\forall k \grgl k_0$ $0 \klgl \betrag{A_k} \klgl B_k$
				\item $\sum B_k$ konvergent
						\begin{align*}
						\Lra &\forall \epsilon > 0 \exists N(\epsilon) \forall n \grgl N(\epsilon)\\
						&\betrag{\underbrace{\sum_{k = 0}^n B_k - \sum_{k = 0}^{\infty}}_{\sum_{k = n + 1}^{\infty} B_k}} < \epsilon
						\end{align*}
						$n \grgl \max(N, k_0)$
						\[\betrag{\sum_{k = 0}^n A_k - \sum_{k = 0}^{\infty} A_k} = \betrag{\sum_{k = n + 1}^{\infty} A_k} \klgl \sum_{k = n + 1}^{\infty} \betrag{A_k} \klgl \sum_{k = n + 1}^{\infty} B_k \klgl \epsilon\]
						\Ra $\sum_{k = 0}^{\infty} \betrag{A_k}$ konvergent; $\sum_{k = 0}^{\infty} A_k$ konvergent\\
						Nebenergebnis:
						$\sum \betrag{A_k}$ konvergent \Ra $\sum A_k$ konvergent
				\end{itemize}
		\item Minorantenkriterium\\
%				$\sum \betrag{A_k} \tx{ divergent}}_{\Lra \sum \betrag{A_k} = \infty}$ und $\forall k \grgl k_0 \mal \betrag{A_k} \klgl B_k$
				\[\sum_{k = 0}^{\infty} B_k = \sum_{k = k_0}^{\infty} B_k \grgl \sum_{k = k_0}^{\infty} \betrag{A_k} = \infty\]
				\Ra $\sum_{k = 0}^{\infty} B_k$ divergent
		\end{enumerate}

\item Wurzelkriterium
		\begin{description}
		\item[Fall 1] $w < 1$
				\begin{align*}
				\lim_{k \ra \infty} \sqrt[k]{\betrag{A_k}} = w < 1 \Lra &\forall \epsilon > 0 \exists N(\epsilon) \forall n \grgl N (\epsilon)\\
				&\betrag{\sqrt[n]{\betrag{a_n}} - w} < \epsilon
				\end{align*}
				\begin{align*}
				\Ra &\sqrt[n]{\betrag{A_n}} - w < \epsilon\\
				\Lra &\sqrt[n]{\betrag{A_n}} < \epsilon + w = q\\
				\Ra & \betrag{A_n} < q^n\\
				&\tx{Sei $\epsilon$ so klein, dass $q < 1$}
				\end{align*}
				\begin{align*}
				\Ra 0 &\klgl \sum_{k = 0}^{\infty} \betrag{A_k} = \sum_{k = 0}^{N - 1} \betrag{A_k} + \sum_{k = N}^{\infty} \betrag{A_k} \klgl \sum_{k = 0}^{N - 1} \betrag{A_k} + \sum_{k = N}^{\infty} q^k\\
				&\klgl \underbrace{\sum_{k = 0}^{N - 1} \betrag{A_k}}_{\begin{minipage}{18mm}\begin{center}\scriptsize{End Summe \Ra konvergent}\end{center}\end{minipage}} + \underbrace{\sum_{k = 0}^{\infty} q^n}_{\begin{minipage}{18mm}\begin{center}\scriptsize{konvergent (geometrische Reihe mit $q < 1$)}\end{center}\end{minipage}} < \infty
				\end{align*}
				\Ra $\sum \betrag{A_k}$ konvergent \Ra $\sum A_k$ konvergent

		\item[Fall 2] $w > 1$
				\begin{align*}
				\lim_{k \ra \infty} \sqrt[k]{\betrag{A_k}} = w > 1 \Lra &\forall \epsilon > 0 \exists N(\epsilon): \forall n \grgl N(\epsilon)
				&\betrag{\sqrt[q]{\betrag{A_n}} - w} < \epsilon
				\end{align*}
				\begin{align*}
				\Ra & - \epsilon < \sqrt[n]{\betrag{A_n}} - w\\
				\Lra & \underbrace{w - \epsilon}_{= Q} < \sqrt[n]{\betrag{A_n}}\\
				\Ra & Q^n < \betrag{A_n}\\
				&\tx{$\epsilon$ so klein, dass $Q = w - \epsilon > 1$}
				\end{align*}
				\[\sum_{k = 0}^{\infty} \betrag{A_k} = \sum_{k = N - 1}^{\infty} \betrag{A_k} + \sum_{k = N}^{\infty} \betrag{A_k} \grgl \sum_{k = N}^{\infty} \betrag{A_k} \grgl \sum_{k = N}^{\infty} Q^k = \infty\]
				\Ra $\sum \betrag{A_k}$ divergent (Rechnug f�r $\sum_{A_k}$ analog)
		\end{description}

\item Beweis Quotientenkriterium �hnlich wie Wurzelkriterium \qed
\end{enumerate}
\end{fbeweis}

\begin{fsatz}[Cauchycher Verdichtungssatz]
Sei $(A_k)_{k \in \N}$ monoton fallend ($A_K > 0$), dann gilt
\[\sum_{k = 1}^{\infty} \tx{ konvergent } \Lra \sum_{k = 0}^{\infty} 2^k \mal A_{2k} \tx{ konvergent}\]
\end{fsatz}

\begin{fbeweis}
Zu zeigen $\sum_{k = 1}^{2^{n + 1} - 1} A_k \klgl \sum_{k = 0}^n 2^k A_{2k} \klgl 2 \sum_{k = 1}^{2n} A_k$
\end{fbeweis}

\begin{beispiel}
\mbox{}\par
\begin{enumerate}
\item Harmonische Reihe: $\sum_{k = 1}^{\infty} \underbrace{\frac{1}{k}}_{A_k}$ divergent, denn $\sum 2^k \mal \frac{1}{2^k} = \sum_{k = 1}^{\infty} 1 = \infty$ \Ra $\sum 2^k \mal \frac{1}{2^k}$ divergent. Damit ist nach dem Verdichtungssatz auch $\sum_{k = 1}^{\infty} \frac{1}{k}$ divergent
\item Alternierende harmonische Reihe
		\[\sum_{k = 1}^{\infty} (-1)^k \mal \frac{1}{k} = \mathbf{- \ln(2)}\]
		konvergent nach Leibnitz, da $A_k \frac{1}{k}$ monoton fallend mit $\lim_{k \ra \infty} A_k = \lim_{k \ra \infty} \frac{1}{k} = 0$

\item Verallgemeinerte harmonische Reihe
		\[\sum_{k = 1}^{\infty} \frac{1}{k^{\alpha}} = \begin{cases}\alpha > 1 & \tx{konvergent}\\\alpha \klgl 1 & \tx{divergent}\end{cases}\]
		Wir betrachten
		\[\sum_{k = 0}^{\infty} 2^k \mal \frac{1}{\rklamm{2^k}^{\alpha}} = \sum_{k = 0}^{\infty} 2^{k (1 - \alpha)} = \sum_{k = 0}^{\infty} \rklamm{2^{1 - \alpha}}^k\]
		\begin{itemize}
		\item divergent f�r $\alpha \klgl 1$ (da $^{1 - \alpha} \grgl 1$ und geometrische Reihe daf�r divergent)
		\item konvergent f�r $\alpha > 1$ (da $2^{1 - \alpha} < 1$ und geometrische Reihe f�r $q < 1$ konvergent)
		\end{itemize}
		\[\sum_{k = 1}^{\infty} \frac{1}{k^2} = \frac{\pi^2}{6}\]

\item $\sum_{k = 0}^{\infty} \frac{x^k}{k!}$, $x \in \mb{R}$\\
		Quotientenkriterium
		\begin{align*}
		&\lim_{k \ra \infty} \betrag{\frac{A_{k + 1}}{A_k}} = \lim_{k \ra \infty} \betrag{\frac{x^{k + 1}}{(k + 1)!} \mal \frac{k!}{x^k}}
		=& \lim_{k \ra \infty} \betrag{\frac{x}{k + 1}} = \betrag{x} \mal \lim_{k \ra \infty} \frac{1}{k + 1} = 0 < 1
		\end{align*}
		\Ra $\sum_{k = 0}^{\infty} \frac{x^k}{k!}$
\end{enumerate}
\end{beispiel}

\section{Aufgabe 2.3}
\label{sec:Konvergente_Reihen_A2_3}
Untersuchen Sie die folgenden Reihen auf Konvergenz:
\begin{enumerate}[label=\alph*)]
\item $\sum_{n = 1}^{\infty} \frac{1}{2+ 3^n}$
\item $\sum_{n = 1}^{\infty} \sqrt[n]{n^2}$
\item $\sum_{k = 1}^{\infty} (-1)^{k + 1} \frac{k^2}{k^3 + 1}$
\end{enumerate}

L�sung siehe \vref{sec:Konvergente_Reihen_A2_3L}.

\section{Aufgabe 2.4}
\label{sec:Konvergente_Reihen_A2_4}
F�r welche Werte von $r$ konvergiert die Reihe
\[1+2r+r^2+2r^3+r^4+2r^5+r^6+\dots\]
Bestimmen Sie im Fall der Konvergenz den Grenzwert.

L�sung siehe \vref{sec:Konvergente_Reihen_A2_4L}.

\begin{fsatz}[Cauchy-Produkt]
%MARK: Satz 2.13
Seien $\sum_{k = 0}^{\infty} \betrag{A_k}$ und $\sum_{k = 0}^{\infty} \betrag{B_k}$ konvergente Reihen. Dann ist
\[\rklamm{\sum_{k = 0}^{\infty} A_k} \mal \rklamm{\sum_{k = 0}^{\infty} B_k} = \sum_{k = 0}^{\infty} C_k\]
mit $C_k = \sum_{l = 0}^{k} A_l \mal B_{k - l}$
\end{fsatz}

\begin{beweis}
-
\end{beweis}

\section{L�sungen}
\subsection{Aufgabe 2.3}
\label{sec:Konvergente_Reihen_A2_3L}
L�sung zu Aufgabe \vref{sec:Konvergente_Reihen_A2_3}.

\begin{enumerate}[label=\alph*)]
\item $\sum_{n = 1}^{\infty} \frac{1}{2+ 3^n}$
		\begin{itemize}
		\item notwendiges Kriterium
				\[\lim_{n \ra \infty} A_n = \lim_{n \ra \infty} \frac{1}{2 + 3^n} = 0 \Ra \sum A_n \tx{ kann konvergent sein}\]

		\item Majorantenkriterium
				\[\sum_{n = 1}^{\infty} \frac{1}{2 + 3^n} \klgl \underbrace{\sum_{n = 1}^{\infty} \rklamm{\frac{1}{3}}^n}_{\begin{minipage}{10mm}\begin{center}\scriptsize{konvergente geometrische Reihe}\end{center}\end{minipage}} \Ra \tx{ konvergent}\]

		\item Quotientenkriterium
				\begin{align*}
				\lim_{n \ra \infty} \betrag{\frac{A_{n + 1}}{A_n}} &= \lim_{n \ra \infty} \betrag{\frac{1}{2 + 3^{n + 1}} \mal \frac{2 + 3^n}{1}} = \lim_{n \ra \infty} \frac{2 + 3^n}{2 + 3^{n + 1}}\\
				&= \lim_{n \ra \infty} \frac{\frac{2}{3^n} + 1}{\frac{2}{3^n} + 3} = \frac{\lim_{n \ra \infty} \rklamm{\frac{2}{3^n} + 1}}{\lim_{n \ra \infty} \rklamm{\frac{2}{3^n} + 3}} = \frac{1}{3} < \Ra \tx{ konvergent}
				\end{align*}

		\item Wurzelkriterium
				\[0 \klgl \lim_{n \ra \infty} \sqrt[n]{\betrag{A_n}} = \lim_{n \ra \infty} \sqrt[n]{\betrag{\frac{1}{2 + 3^n}}} \klgl \lim_{n \ra \infty} \sqrt[n]{\frac{1}{3^n}} = \lim_{n \ra \infty} \frac{1}{3} = \frac{1}{3} \Ra \tx{ konvergent}\]

		\item Verdichtungssatz
				\[\sum_{n = 0}^{\infty} 2^n \mal \frac{1}{\underbrace{2 + 3^{2^n}}_{= A_{2^n}}} \klgl \sum_{n = 0}^{\infty} 2^n \mal \rklamm{\frac{1}{3^2}}^n = \underbrace{\sum_{n = 0}^{\infty} \rklamm{\frac{2}{9}}^n}_{\tx{konvergente geometrische Reihe}}\]
				\[\sum_{n = 0}^{\infty} 2^n \frac{2}{2 + 3^{2^n}} \tx{ konvergent } \Lra \sum_{n = 1}^{\infty} \frac{1}{2 + 3^n} \tx{ konvergent}\]
		\end{itemize}

\item $\sum_{n = 1}^{\infty} \sqrt[n]{n^2}$
		notwendiges Kriterium:
		\[\lim_{n \ra \infty} A_n = \lim_{n \ra \infty} \sqrt[n]{n^2} = \rklamm{\lim_{n \ra \infty} \sqrt[n]{n}}^2 = 1 \Ra \tx{ divergent}\]

\item $\sum_{k = 1}^{\infty} (-1)^{k + 1} \frac{k^2}{k^3 + 1}$
		\begin{description}
		\item[Leibnitz Kriterium] $A_K = \frac{k^2}{k^3 + 1}$
				\[\lim_{n \ra \infty} A_k = \lim_{k \ra \infty} \frac{k^2}{k^3 + 1} = \lim_{n \ra \infty} \frac{1}{k + \frac{1}{k^2}} = 0\]
				$(A_k)$ Nullfolge

		\item[Monoton fallend] zu zeigen $\frac{(k + 1)^2}{(k + 1)^3 + 1} < \frac{k^2}{k^3 + 1}$
				$\Ra \sum_{k = 1}^{\infty} (-1) \frac{k^2}{k^3 + 1}$ konvergent
		\end{description}
\end{enumerate}

\subsection{Aufgabe 2.4}
\label{sec:Konvergente_Reihen_A2_4L}
L�sung zu Aufgabe \vref{sec:Konvergente_Reihen_A2_4}.

$1+2r+r^2+2r^3+r^4+2r^5 + \dots$
\begin{align*}
&\sum_{k = 0}^{\infty} r^{2k} + 2 \sum_{k = 0}^{\infty} r^{2k + 1}\\
&= \sum_{k = 0}^{\infty} \rklamm{r^2}^k + 2r \sum_{k = 0}^{\infty} \rklamm{r^2}^k\\
\end{align*}
konvergent \Lra $r^2 < 1$ \Lra $r \in (-1, 1)$

%\chapter{Potenzreihen und elementare Funktionen}
\begin{fdefinition}[Potenzreihe]
%MARK: Definition 2.14
Die (formale) Reihe $\sum_{n = 0}^{\infty} a_n (x - x_0)^n$ hei�t \indexu{Potenzreihe} mit Entwicklungspunkt $x_0$ und Koeffizienten $(A_n)_{n \in \N_0}$
\end{fdefinition}

\begin{bemerkung}
Man setzt $0^0 = 1$
\end{bemerkung}

\begin{fsatz}
\label{satz:Potenzreihen_und_elementare_Funktionen_S2_15}
F�r die Potenzreihe $\sum_{n = 0}^{\infty} a_n (x - x_0)^n$ gilt:
\begin{enumerate}[label=\roman*)]
\item konvergiert die Reihe $a_n$ einer Stelle $x_1 \neq x_0$, so konvergiert sie f�r alle $x$ mit $\betrag{x - x_0} < \betrag{x_1 - x_0}$
\item divergiert die Reihe $a_n$ einer STelle $x_2 \neq x_0$, so divergiert sie f�r alle $x$ mit $\betrag{x - x_0} > \betrag{x_2 - x_0}$
\end{enumerate}
\Mark{Satz 2.15}
\end{fsatz}

\begin{fbeweis}
\mbox{}\par
\begin{enumerate}[label=\roman*)]
\item $\sum_{n = 0}^{\infty} a_n (x_1 - x_0)^n$ konvergent \Ra $\lim_{n \ra \infty} a_n (x_1 - x_0)^n = 0$
		\begin{align*}
		&\sum_{n = 0}^{\infty} a_n (x - x_0)^n\\
		\Lra & \forall \epsilon > 0 \exists N(\epsilon) \forall k \grgl N(\epsilon)\\
		&\betrag{a_n(x_1 - x_0)^n} < \epsilon\\
		\Ra & \sum_{n = 0}^{\infty} \betrag{a_n (x - x_0)^n} = \sum_{n = 0}^{N - 1} \betrag{a_n (x - x_0)^n} + \sum_{n = N}^{\infty} \betrag{_n (x - x_0)^n}\\
		= & \sum_{n = 0}^{N - 1} \underbrace{\betrag{a_n (x - x_0)^n}}_{\tx{endlich}} + \sum_{n = N}^{\infty} \betrag{\underbrace{a_n (x_1 - x_0)^n}_{< \epsilon}} \betrag{\underbrace{\rklamm{\frac{x - x_0}{x_1 - x_0}}^n}_{< 1}}\\
		= & \underbrace{\sum_{n = 0}^{N - 1} \betrag{a_n (x - x_0)}}_{\tx{endlich}} + \epsilon \underbrace{\sum_{n = N}^{\infty} \betrag{\underbrace{\frac{x - x_0}{x_1 - x_0}}_{= q < 1}}^n}_{\begin{minipage}{2cm}\begin{center}\scriptsize{konvergente geometrische Reihe}\end{center}\end{minipage}} \Ra \tx{ konvergent}
		\end{align*}

\item Widerspruchsbeweis
		\[\sum_{n = 0}^{\infty} a_n (x_2 - x_0)^n \tx{ divergent}\]
		Annahme: $\exists x_3$ mit $\betrag{x_3 - x_0} > \betrag{x_2 - x_0}$
		\[\sum_{n = 0}^{\infty} a_n (x_3 - x_0)^n \tx{ konvergent}\]
		\stack{i)}{\Ra} $\sum_{n = 0}^{\infty} a_n (x_2 - x_0)^n$ konvergent \blitza
\end{enumerate}
\end{fbeweis}

\begin{fsatz}[Konvergenzradius einer Potenzreihe]
F�r jede Potenzreihe $\sum_{n = 0}^{\infty} a_n (x - x_0)^n$ existiert ein eindeutig bestimmter \indexu{Konvergenzradius} $r$ ($0 \klgl r \klgl \infty$) f�r den gilt:
\begin{enumerate}
\item Falls $r = 0$ so konvergiert die Potenzreihe nur f�r $x = x_0$
\item Falls $0 < r < \infty$ so konvergiert die Potenzreihe f�r alle $x \in \mb{R}$ mit $\betrag{x - x_0} < r$ und divergiert f�r alle $x \in \mb{R}$ mit $\betrag{x - x_0} > r$.
\item Falls $r = \infty$, so konvergiert die Potenzreihe f�r alle $x \in \mb{R}$
\end{enumerate}
\Mark{Satz 2.16}
\end{fsatz}

\begin{fbeweis}
Folgt aus Satz \vref{satz:Potenzreihen_und_elementare_Funktionen_S2_15}.
\Solved{Insert ref to Satz 2.15}{}
\end{fbeweis}

\begin{fdefinition}[Konvergenzbereich]
%MARK: Definition 2.17
Der \indexu{Konvergenzbereich} einer Potenzreihe $\sum_{n = 0}^{\infty} a_n (x - x_0)^n$ ist $K = \gklamm{x \in \mb{R} \vert \sum_{n = 0}^{\infty} a_n (x - x_0)^n \tx{ konvergent}}$. Es gilt:
\[(x_0 - r, x_0 + r) \subseteq K \subseteq \eklamm{x_0 - r, x_0 + r}\]
\end{fdefinition}

\begin{beispiel}
geometrische Reihe $\sum_{n = 0}^{\infty} (x - x_0)^n$
\begin{itemize}[label=\Ra]
\item Konvergenzradius $r = 1$
\item $(x_0 - 1, x_0 + 1) \subseteq K \subseteq \eklamm{x_0 - 1, x_0 + 1}$
		\begin{description}
		\item[$x = x_0 - 1$] $\sum_{n = 0}^{\infty} \rklamm{\rklamm{x_0 - 1} - x_0}^n = \sum_{n = 0}^{\infty} (-1)^n$ divergent
		\item[$x = x_0 + 1$] $\sum_{n = 0}^{\infty} \rklamm{\rklamm{x_0 + 1} - x_0}^n = \sum_{n = 0}^{\infty} 1^n$ divergent
		\end{description}
\item $K = (x_0 -1, x_0 + 1)$
\end{itemize}
\end{beispiel}

\begin{fsatz}[Bestimmung des Konvergenzradius]
%MARK: Satz 2.18
Sei $\sum_{n = 0}^{\infty} a_n (x - x_0)^n$ gegeben
\begin{enumerate}[label=\roman*)]
\item Existiert $q = \lim_{n \ra \infty} \betrag{\frac{a_n}{a_{n + 1}}}$, so ist der Konvergenzradius $r = q$
\item Existiert $w = \lim_{n \ra \infty} \sqrt[n]{\betrag{a_n}}$, so ist der Konvergenzradius $r = \frac{1}{w}$\\
		Wobei "`$\frac{1}{0}$"'$ = \infty$ und "`$\frac{1}{\infty}$"'$ = 0$
\end{enumerate}
\end{fsatz}

\begin{fbeweis}
Anwendung des Quotienten \ac{bzw.} Wurzelkriteriums auf die Reiher $\sum_{n = 0}^{\infty} A_n$ mit
\[A_n = a_n (x - x_0)^n\]
\end{fbeweis}

\begin{beispiel}
\mbox{}\par
\begin{enumerate}
\item $\sum_{n = 0}^{\infty} \frac{x^n}{n!} = \sum_{n = 0}^{\infty} \ub{\frac{1}{n!}}{= a_n} \mal \rkl{x - \ub{0}{= x_0}}^n$\\
		Entwicklungspunkt $x_0 = 0$\\
		Quotientenmethode
		\begin{align*}
		\lim_{n \ra \infty} \betrag{\frac{a_n}{a_{n + 1}}} &= \lim_{n \ra \infty} \betrag{\frac{1}{n!} \mal \frac{(n + 1)!}{1}} = \lim_{n \ra \infty} \frac{(n + 1)!}{n!}\\
		&\lim_{n \ra \infty} (n + 1) = \infty
		\end{align*}
		\begin{itemize}[label=\Ra]
		\item Konvergenzradius $r = q = \infty$
		\item Konvergenzbereich $K = \mb{R}$
		\end{itemize}

\item $\sum_{n = 0}^{\infty} \frac{1}{2^n} x^n = \sum_{n = 0}^{\infty} \ub{\frac{1}{2^n}}{= a_n} \rkl{x - \ub{0}{= x_0}}^n$\\
		Entwicklungspunkt $x_0 = 0$\\
		Wurzelmethode
		\[w = \lim_{n \ra \infty} \sqrt[n]{\betrag{a_n}} = \lim_{n \ra \infty} \sqrt[n]{\betrag{\frac{1}{2^n}}} = \lim_{n \ra \infty} \frac{1}{2} = \frac{1}{2}\]
		\begin{itemize}[label=\Ra]
		\item Konvergenzradius $r = \frac{1}{w} = 2$\\
				$(-2, 2) \subseteq K \subseteq \eklamm{-2 , 2}$
				\begin{description}
				\item[$x = -2$] $\sum_{n = 0}^{\infty} \frac{1}{2^n} (-2)^n = \sum_{n = 0}^{\infty} (-1)^n$ divergent
				\item[$x = 2$] $\sum_{n = 0}^{\infty} \frac{1}{2^n} (2)^n = \sum_{n = 0}^{\infty} 1^n$ divergent
				\end{description}

		\item $K = (-2, 2)$
		\end{itemize}

\item $\sum_{n = 0}^{\infty} n^n (x + 2)^n = \sum_{n = 0}^{\infty} \ub{n^n}{a_n} \rkl{x - \ub{(-2)}{= x_0}}^n$\\
		Entwicklungspunkt $x_0 = -2$
		Wurzelmethode
		\[w = \lim_{n \ra \infty} \sqrt[n]{\betrag{a_n}} = \lim_{n \ra \infty} \sqrt[n]{n^n} \lim_{n \ra \infty} n = \infty\]
		\begin{itemize}[label=\Ra]
		\item Konvergenzradius $r = \frac{1}{w} = 0$
		\item $K =\gklamm{-2}$
		\end{itemize}
\end{enumerate}
\end{beispiel}

\section{Aufgabe 2.5}
\label{sec:Potenzreihen_und_elementare_Funktionen_A2_5}
Bestimmen Sie den Entwicklungspunkt, den Konvergenzradius und den Konvergenzbereich der folgenden Potenzreihen.
\begin{enumerate}[label=\alph*)]
\item $\sum_{n = 0}^{\infty} 3n (x + 1)^n$
\item $\sum_{n = 0}^{\infty} \frac{1}{n^n} \mal x^n$
\item $\sum_{k = 0}^{\infty} \frac{(2x + 1)^k}{k + 3}$
\end{enumerate}

L�sung siehe \vref{sec:Potenzreihen_und_elementare_Funktionen_A2_5L}.

\begin{bemerkung}
Die Potenzreihe $\sum_{n = 0}^{\infty} a_n \mal (x - x_0)^n$ definiert in ihrem Konvergenzbereich eine Funktion
\begin{align*}
f: &K \ra \mb{R}\\
&x \mapsto f(x) = \sum_{n = 0}^{\infty} a_n \mal (x - x_0)^n
\end{align*}
Auf diese Weise kann man einige elementare Funktionen definieren. Zum Beispiel Exponentialfunktion, Sinus, Cosinus.
\end{bemerkung}

\begin{fdefinition}[Exponentialfunktion]
%MARK: Definition 2.20
Die Funktion
\begin{align*}
\exp: \mb{R}& \ra \mb{R}\\
x& \mapsto \exp (x) = \sum_{n = 0}^{\infty} \frac{1}{n!} \mal x^n = e^x
\end{align*}
hei�t Exponentialfunktion

\begin{tikzpicture}[domain=-2.5:1.2]
\draw[->,thick] (-2.5,0) -- (4.5,0) node[below] {$x$};
\draw[->,thick] (0,-0.5) -- (0,3.5);
\node at (0,0) [below right] {$0$};
\draw[thick] (-0.05,1) -- (0.05,1) node[left] {$1$};

\foreach \x in {1,2,3,4}
	\draw[thick] (\x,-0.05) -- (\x,0.05);

\draw plot[smooth] (\x,{exp(\x)}) node[right] {$\exp (x)$};
\end{tikzpicture}
\Img{MA-16.04.2009-IMG-1}

\end{fdefinition}

\begin{fsatz}[Eigenschaften der Exponentialfunktion]
%MARK: Satz 2.21
\mbox{}\par
\begin{enumerate}[label=\roman*)]
\item $\forall x, y \in \mb{R}$: $e^x \mal e^y = e^{x + y}$ Funktionalgleichung
\item $e^0 = 1$, $e^1 = e$, $e^{-x} = \frac{1}{e^x}$
\item $x < 0$: $0 < e^x < 1$\\
		$x > 0$: $e^x > 1$

\item $\lim_{x \ra -\infty} e^x = 0$, $\lim_{x \ra \infty} e^x = \infty$\\
		$\forall n \in \N$: $\lim_{x \ra \infty} \frac{x^n}{e^x} = 0$
\end{enumerate}
\end{fsatz}

\begin{fbeweis}
\mbox{}\par
\begin{enumerate}[label=\roman*)]
\item $e^x \mal e^y$
		\[e^x \mal e^y = \rkl{\sum_{n = 0}^{\infty} \ub{\frac{x^n}{x!}}{= A_n}} \mal \rkl{\sum_{n = 0}^{\infty} \ub{\frac{y^n}{n!}}{= B_n}} = \sum_{n = 0}^{\infty} C_n\]
		\begin{align*}
		C_n &= \sum_{l = 0}^{n} A_l \mal B_{n - l} = \sum_{l = 0}^{n} \frac{x^l}{l!} \mal \frac{y^{n - l}}{(n - l)!} = \frac{1}{n!} \sum_{l = 0}^{n} \ub{\frac{n!}{l! \mal (n - l)!}}{\binom{n}{l}} \mal x^l \mal y^{n - l}\\
		&= \frac{1}{n!} (x + y)^n\\
		\Ra& e^x \mal e^y = \sum_{n = 0}^{\infty} \frac{1}{n!} \mal (x + y)^n = e^{x + y}
		\end{align*}

\item Verwenden Funktionalgleichung
		\begin{align*}
		&1 = e^0 = e^{x - x} = e^x \mal e^{-x}\\
		\Lra & e^{-x} = \frac{1}{e^x}
		\end{align*}

\item $x > 0$
		\[\sum_{n = 0}^{\infty} \ub{\frac{x^n}{n!}}{\grgl 0} \grgl \frac{x^0}{0!} + \frac{x^1}{1!} = 1 + \ub{x}{\grgl 0} > 1\]
		$x < 0$
		\begin{align*}
		y &= -x > 0\\
		0 \stack{*}{<} e^x &= e^{-y} = \frac{1}{e^y} < 1
		\end{align*}
		$*$ verwendet, dass $e^x < \infty$

\item \begin{itemize}
		\item $\lim_{x \ra \infty} e^x \grgl \lim_{x \ra \infty} (1 + x) = \infty$
		\item $y = -x$
				\[\lim_{n \ra -\infty} e^x = \lim_{y \ra \infty} e^{-y} = \lim_{y \ra \infty} \frac{1}{e^y} = 0\]

		\item \ac{z.z.} $\forall n \in \N$: $\lim_{x \ra \infty} \frac{x^n}{e^x} = 0$\\
				\Lra \ac{z.z.} $\forall n \in \N$: $\lim_{x \ra \infty} \frac{e^x}{x^n} = \infty$
				\begin{align*}
				\lim_{x \ra \infty} \frac{e^x}{x^n} &= \lim_{x \ra \infty} \sum_{k = 0}^{\infty} \frac{1}{x^n} \mal \frac{x^k}{k!} = \lim_{x \ra \infty} \sum_{k = 0}^{\infty} \ub{\frac{x^{k - n}}{k!}}{\grgl 0} \grgl \lim_{x \ra \infty} \sum_{k = n}^{\infty} \frac{x^{k - n}}{k!}\\
				&\grgl \lim_{x \ra \infty} \rkl{\frac{x^0}{n!} + \frac{x^1}{(n + 1)!}} = \infty
				\end{align*} \qed
		\end{itemize}
\end{enumerate}
\end{fbeweis}

\begin{fdefinition}[Trigonometrische Funktionen]
%MARK: Definition 2.21
\mbox{}\par
\begin{enumerate}
\item Cosinus
		\begin{align*}
		\cos:~\R& \ra \R\\
		x& \mapsto \cosx{x} = \sum_{n = 0}^{\infty} (-1)^n \frac{x^{2n}}{(2n)!}
		\end{align*}

\item Sinus
		\begin{align*}
		\sin:~\R& \ra \R\\
		x& \mapsto \sinx{x} = \sum_{n = 0}^{\infty} (-1)^n \mal \frac{x^{2n + 1}}{(2n + 1)!}
		\end{align*}

\item Tangens
		\begin{align*}
		\tan:~\R& \backslash \gklamm{x \in \R \vert x = (2k + 1) \frac{\pi}{2} K \in \mb{Z}} \ra \R\\
		x& \mapsto \tanx{x} = \frac{\sinx{x}}{\cosx{x}}
		\end{align*}
\end{enumerate}

\begin{tikzpicture}[domain=-pi:pi]
\draw[->,thick] (-3.5,0) -- (3.5,0) node[below] {$x$};
\draw[->,thick] (0,-1) -- (0,2);
\node at (0,0) [below right] {$0$};
\node at (-3.5,1) {$\sin$};

\foreach \x/\y in {-pi/-\pi,-1.570/-\frac{\pi}{2},1.570/\frac{\pi}{2},pi/\pi}
	\draw[thick] (\x,-0.05) -- (\x,0.05) node[above] {$\y$};

\draw plot[smooth] (\x,{sin(\x r)});
\end{tikzpicture}

\begin{tikzpicture}[domain=-pi:pi]
\draw[->,thick] (-3.5,0) -- (3.5,0) node[below] {$x$};
\draw[->,thick] (0,-1) -- (0,2);
\node at (0,0) [below right] {$0$};
\node at (-3.5,1) {$\cos$};

\foreach \x/\y in {-pi/-\pi,-1.570/-\frac{\pi}{2},1.570/\frac{\pi}{2},pi/\pi}
	\draw[thick] (\x,-0.05) -- (\x,0.05) node[above] {$\y$};

\draw plot[smooth] (\x,{cos(\x r)});
\end{tikzpicture}

\begin{tikzpicture}
\draw[->,thick] (-3.5,0) -- (3.5,0) node[below] {$x$};
\draw[->,thick] (0,-1) -- (0,2);
\node at (0,0) [below right] {$0$};
\node at (-3.5,1) {$\tan$};

\foreach \x/\y in {-pi/-\pi,-1.570/-\frac{\pi}{2},1.570/\frac{\pi}{2},pi/\pi}
	\draw[thick] (\x,-0.05) -- (\x,0.05) node[above] {$\y$};

\draw plot[domain=0:1.2,smooth] (\x,{tan(\x r)});  \draw plot[domain=-1.2:0,smooth] (\x,{-tan(-\x r)});
\end{tikzpicture}
\Insert{MA-16.04.2009-IMG-2}

\end{fdefinition}

\begin{fsatz}[Eigenschaften der trigonometrischen Funktionen]
\mbox{}\par
\begin{enumerate}[label=\roman*)]
\item $\cosx{0} = 1, \sinx{0} = 0$
\item $\cosx{-x} = \cosx{x}$ gerade Funktion\\
		$\sinx{-x} = -\sinx{x}$ ungerade Funktion

\item $\cosx{x + y} = \cosx{x} \mal \cosx{y} - \sinx{x} \mal \sinx{y}$\\
		$\sinx{x + y} = \sinx{x} \mal \cosx{y} + \cosx{x} \mal \sinx{y}$

\item $\cosx{x - y} = \cosx{x} \mal \cosx{y} + \sinx{x} \mal \sinx{y}$\\
		$\sinx{x - y} = \sinx{x} \mal \cosx{y} - \cosx{x} \mal \sinx{y}$

\item $\cosx{x} \mal \cosx{y} = \frac{1}{2} \rkl{\cosx{x + y} + \cosx{x - y}}$\\
		$\sinx{x} \mal \sinx{y} = \frac{1}{2} \rkl{\cosx{x - y} - \cosx{x + y}}$\\
		$\sinx{x} \mal \cosx{y} = \frac{1}{2} \rkl{\sinx{x + y} + \sinx{x - y}}$

\item $\cosx{2x} = \cos^2(x) - \sin^2(x)$\\
		$\sinx{2x} = 2 \sinx{x} \mal \cosx{x}$

\item $\cos^2(x) + \sin^2(x) = 1$\\
		$-1 \klgl \sinx{x}, \cosx{x} \klgl 1$

\item F�r alle $k \in \Z$ gilt:
		\[\left.\begin{array}{l}
		\cosx{x + 2\pi k} = \cosx{x}\\
		\sinx{x + 2\pi k} = \sinx{x}
		\end{array}\right\} 2 \pi\tx{-periodisch}\]
		\[\cosx{(2k + 1) \mal \frac{\pi}{2}} = 0\]
		\[\cosx{k \mal \pi} = (-1)^k\]
		\[\sinx{k \mal \pi} = 0\]
		\[\sinx{(2 \mal k + 1) \mal \frac{\pi}{2}} = (-1)^k\]
%TODO: Align to left
\end{enumerate}
\end{fsatz}

\section{Aufgabe 2.6}
\label{sec:Potenzreihen_und_elementare_Funktionen_A2_6}
Zu zeigen: $\cos^2(x) + \sin^2(x) = 1$ (Verwenden Summendarstellung von $\sin$, $\cos$ und Cauchy-Produkt)

L�sung siehe \vref{sec:Potenzreihen_und_elementare_Funktionen_A2_6L}.

\begin{bemerkung}
\mbox{}\par
\begin{enumerate}
\item Die Zahl $\pi$ kann man als das doppelte der ersten positiven Nullstelle vom Cosinus definieren.
		\Insert{IMG-MA-09-04-21-1} Diese verwendet man zur einfacheren Berechnung (einfacher als �ber den Umfang von Vielecken) von $\pi$.

\item Eulersche Beziehung: $i = \sqrt{-1}$, $x \in \R$
		\begin{align*}
		e^{(ix)} &= \sum_{n = 0}^{\infty} \frac{(ix)^n}{n!} = \sum_{n = 0}^{\infty} \frac{(ix)^{2n}}{(2n)!} + \sum_{n = 0}^{\infty} \frac{(ix)^{2n + 1}}{(2n + 1)!}\\
		&= \sum_{n = 0}^{\infty} (i)^{2n} \frac{x^{2n}}{(2n)!} + \sum_{n = 0}^{\infty} i \mal \ub{(i)^{2n}}{= (-1)^n} \frac{x^{2n + 1}}{(2n + 1)!}\\
		&= \sum_{n = 0}^{\infty} (-1)^n \frac{x^{2n}}{(2n)!} + i \sum_{n = 0}^{\infty} (-1)^n \mal \frac{x^{2n + 1}}{(2n + 1)!}\\
		&= \cosx{x} + i \mal \sinx{x}
		\end{align*}
		\Insert{IMG-MA-09-04-21-2}

		$z = x + iy$ $x,y \in \R$\\
		$\overline{z} = x - iy$
		\Todo{Herausfinden ob es $\overline{z}$ oder nur $z$ ist}

\item $\frac{1}{2} \rkl{e^{ix} + e^{-ix}}$
		\begin{align*}
		= &\frac{1}{2} \rklamm{\cosx{x} + i \mal \sinx{x} + \rklamm{\cosx{-x} + i \mal \sinx{-x}}}\\
		= &\frac{1}{2} \rklamm{\cosx{x} + i \mal \sinx{x} + \rklamm{\cosx{x} - i \mal \sinx{x}}}\\
		= &\cosx{x}\\
		&\frac{1}{2} \rkl{e^{ix} - e^{-ix}} = \dots = \sinx{x}
		\end{align*}
\end{enumerate}

\section{L�sungen}
\subsection{Aufgabe 2.5}
\label{sec:Potenzreihen_und_elementare_Funktionen_A2_5L}
L�sung zu Aufgabe \vref{sec:Potenzreihen_und_elementare_Funktionen_A2_5}.

\begin{enumerate}[label=\alph*)]
\item $\sum_{n = 0}^{\infty} 3n (x + 1)^n$
		\[\sum_{n = 0}^{\infty} 3n (x + 1)^n = \sum_{n = 0}^{\infty} \ub{3n}{a_n} (x - \ub{(-1)}{x_0})^n\]
		Entwicklungspunkt $x_0 = -1$\\
		Wurzelmethode
		\begin{align*}
		\lim_{n \ra \infty} \sqrt[n]{\betrag{a_n}} &= \lim_{n \ra \infty} \sqrt[n]{3n} = \lim_{n \ra \infty} \sqrt[n]{3} \mal \sqrt[n]{n}\\
		&= \lim_{n \ra \infty} \sqrt[n]{3} \mal \lim_{n \ra \infty} \sqrt[n]{n} = 1 \mal 1 = 1
		\end{align*}
		Konvergenzradius $r = 1$
		$(x_0 -r , x_0 + r) = (-2, 0) \subseteq K \subseteq \eklamm{-2, 0}$
		\begin{description}
		\item[$x = -2$] $\sum_{n = 0}^{\infty} 3n (-1)^n$ divergent nach Kriterium
		\item[$x = 0$] $\sum_{n = 0}^{\infty} 3n \mal (1)^n (\lim_{n \ra \infty} 3_n = \infty \neq 0)$ nach dem notwendigen Kriterium divergent
		\end{description}
		\Ra $K = (-2, 0)$

\item $\sum_{n = 0}^{\infty} \frac{1}{n^n} \mal x^n$
		\[\sum_{n = 0}^{\infty} \frac{1}{n^n} \mal x^n = \sum_{n = 0}^{\infty} \ub{\frac{1}{n^n}}{= a_n} \mal (x - \ub{0}{x_0})^n\]
		Entwicklungspunkt $x_0 = 0$\\
		Wurzelmethode
		\[w = \lim_{n \ra \infty} \sqrt[n]{\betrag{a_n}} = \lim_{n \ra \infty} \sqrt[n]{\frac{1}{n^n}} = \lim_{n \ra \infty} \frac{1}{n} = 0\]
		Konvergenzradius $r = \infty$
		\Ra Konvergenbereich $K = \mb{R}$

\item $\sum_{k = 0}^{\infty} \frac{(2x + 1)^k}{k + 3}$
		\[\sum_{k = 0}^{\infty} \frac{(2x + 1)^k}{k + 3} = \sum_{k = 0}^{\infty} \frac{1}{k + 3} \rkl{2 \rkl{x - \rkl{- \frac{1}{2}}}}^k = \sum_{k = 0}^{\infty} \ub{\frac{2^k}{k + 3}}{= a_k} \rkl{x - \rkl{-\frac{1}{2}}}^k\]
		Entwicklungspunkt $x_0 = -\frac{1}{2}$
		Quotientenmethode
		\begin{align*}
		q &= \lim_{k \ra \infty} \betrag{\frac{a_k}{a_{k + 1}}} = \lim_{k \ra \infty} \frac{2^k}{k + 3} \mal \frac{(k + 4)}{2^{k + 1}} = \lim_{k \ra \infty} \frac{k + 3}{2(k + 4)}\\
		&= \frac{1}{2} \lim_{k \ra \infty} \frac{1 + \frac{3}{k}}{1 + \frac{4}{k}} = \frac{1}{2}\\
		&\Ra r = q \frac{1}{2}
		\end{align*}
		$(-1, 0) \subseteq K \subseteq \eklamm{-1, 0} \Ra K = \eklamm{-1, 0}$
\end{enumerate}

\subsection{Aufgabe 2.6}
\label{sec:Potenzreihen_und_elementare_Funktionen_A2_6L}
L�sung zu Aufgabe \vref{sec:Potenzreihen_und_elementare_Funktionen_A2_6}.

\begin{align*}
\sin^2(x) &= \rklamm{\sum_{k = 0}^{\infty} \ub{(-1)^k \mal \frac{x^{2k + 1}}{(2k + 1)!}}{= A_k}} \mal \rklamm{\sum_{k = 0}^{\infty} \ub{(-1)^k \mal \frac{2^{2k + 1}}{(2k + 1)!}}{= B_k}}\\
&\stack{\tx{CP}}{=} \sum_{k = 0}^{\infty} \rklamm{\sum_{l = 0}^{k} A_l \mal B_{k - l}} = \sum_{k = 0}^{\infty} \rklamm{\sum_{l = 0}^{\infty} (-1)^l \frac{x^{2l + 1}}{(2 l + 1)!} \mal (-1)^{k - l} \mal \frac{x^{2 (k - l) + 1}}{(2 (k - l) + 1)!}}\\
&= \sum_{k = 0}^{\infty} (-1)^k x^{2k + 2} \textcolor{blue}{\frac{1}{(2k + 2)!}} \sum_{l = 0}^{k} \ub{\frac{\textcolor{blue}{(2k + 2)!}}{(2l + 1)! \mal (2k - 2l + 1)}}{\binom{2k + 2}{2l + 1}}\\
&= \sum_{k = 0}^{\infty} (-1)^k \mal \frac{x^{2k + 2}}{(2k + 2)!} \mal \frac{2^{2k + 2}}{2}\\
&= \sum_{n = 0}^{\infty} (-1)^n \frac{x^{2n}}{(2n)!} \mal 2^{2n - 1}
\end{align*}

\begin{align*}
\cos^2(x) &= \rklamm{\sum_{n = 0}^{\infty} \ub{\rklamm{-1}^n \mal \frac{x^{2n}}{(2n)!}}{= A_n}} \mal \rkl{\sum_{n = 0}^{\infty} \ub{(-1)^n \mal \frac{x^{2n}}{(2n)!}}{= B_n}}\\
&\stack{\tx{CP}}{=} \sum_{n = 0}^{\infty} \mal \sum_{l = 0}^{n} A_l \mal B_{n - l}\\
&= \sum_{n = 0}^{\infty} \sum_{l = 0}^{n} (-1)^l \mal \frac{x^{2l}}{(2l)!} \mal (-1)^{n - l} \mal \frac{x^{2 (n - l)}}{(2 (n - l))!}\\
&= \sum_{n = 0}^{\un} (-1)^n \frac{x^{2n}}{(2n)!} \sum_{l = 0}^{n} \ub{\frac{(2n)!}{(2l)! \mal (2n - 2l)!}}{= \binom{2n}{2l}}\\
&= \rklamm{\sum_{n = 1}^{\un} (-1)^n \frac{x^{2n}}{(2n)!} \mal 2^{2n - 1}} + 1
\Ra & \sin^2(x) + \cos^2(x) = 1
\end{align*}
\end{bemerkung}

%
%\part{Grenzwerte und Stetigkeit}
%\Mark{Part 3}
%\chapter{Grenzwerte und Stetigkeit von Funktionen}
\Mark{Chapter 3.1}
Wir betrachten eine Funktion $f$: $D \ra \R$ in einer Umgebung des Punktes $x_0 \in D$.

\begin{fdefinition}[Grenzwerte]
\mbox{}\par
\begin{enumerate}[label=\alph*)]
\item Die Funktion $f$ hat im Punkt $x_0$ den \indexu{Grenzwert} $a$, $\lim_{x \ra x_0} f(x) = a$, falls gilt:
		\[\forall(x_n)_{n \in \N} \tx{ mit } x_n \in D, \lim_{n \ra \un} x_n = x_0:~\lim_{n \ra \un} f(x_n) = a\]

\item Die Funktion $f$ hat im Punkt $x_0$ den \indexu{linksseitigen Grenzwert} $a \lim_{x \nearrow x_0} f(x) = a$ (\ac{bzw.} $\lim_{x \uparrow x_0} f (x) = a$, $\lim_{x \ra x_0^-} f(x) = a$), falls gilt:
		\[\forall(x_n)_{n \in \N} \tx{ mit } x_n < x_0, x_n \in D, \lim_{n \ra \infty} x_n ? x_0:~\lim_{n \ra \infty} f(x_n) = a\]

\item Die Funktion $f$ hat im Punkt $x_0$ den \indexu{rechtsseitigen Grenzwert} $a$ $\lim_{x \searrow x_0} f (x) = a$ (\ac{bzw.} $\lim_{x \downarrow x_0} f(x) = a$ \ac{bzw.} $\lim_{x \ra x_0^+} f (x) = a$), falls gilt
		\[\forall(x_n)_{n \in \N} \tx{ mit } x_n \in D, x_n > x_0, \lim_{n \ra \infty} x_n = x_0\]

\item Die Funktion $f$ hat in $\pm \un$ den Grenzwert $a$ $\lim_{x \ra \pm \infty} f (x) = a$, falls gilt:
		\[\forall(x_n)_{n \in \N} \tx{ mit } x_n \in D, \lim_{n \ra \infty} = \pm \infty:~\lim_{n \ra \infty} f(x_n) = a\]

\item Die Funktion $f$ hat in $x_0$ den \indexu{uneigentlichen Grenzwert} $\pm \infty$, $\lim_{x \ra x_0} f (x) = \pm \un$ falls gilt:
		\[\forall(x_n)_{n \in \N} \tx { mit } x_n \in D, \lim_{n \ra \un} x_n = x_0:~\lim_{n \ra \un} f(x_n) = \pm \un\]
		(Links-/rechtsseitige uneigentliche Grenzwert werden analog definiert).
\end{enumerate}
\Mark{Definition 3.1}
\end{fdefinition}

\begin{bemerkung}
Es gilt $\lim_{x \ra x_0} f(x) = a \Lra \lim_{x \nearrow x_0} f(x) = \lim_{x \searrow x_0} f(x) = a$
\end{bemerkung}

\begin{fsatz}[Rechenregeln f�r Grenzwerte]
\label{satz:Grenzwerte_und_Stetigkeit_von_Funktionen_S3_2}
Falls $\lim_{x \ra x_0} f (x) = a$ und $\lim_{x \ra x_0} g (x) = b$ ($a, b \in \R$) gilt:
\begin{enumerate}[label=\roman*)]
\item $\lim_{x \ra x_0} (f(x) \pm g(x)) = \lim_{x \ra x_0} f(x) \pm \lim_{x \ra x_0} g(x) = a \pm b$
\item $\lim_{x \ra x_0} (f(x) \mal g(x)) = \lim_{x \ra x_0} f(x) \mal \lim_{x \ra x_0} g(x) = a \mal b$
\item $\lambda \in \R$: $\lim_{x \ra x_0} \lambda \mal f(x) = \lambda \mal \lim_{x \ra x_0} f(x) = \lambda \mal a$
\item falls $b \neq 0$: $\lim_{x \ra x_0} \frac{f(x)}{g(x)} = \frac{\lim_{x \ra x_0} f(x)}{\lim_{x \ra x_0} g(x)} = \frac{a}{b}$
\end{enumerate}
\Mark{Satz 3.2}
\end{fsatz}

\begin{beweis}
Folgt aus Rechenregeln f�r Grenzwerte von Folgen.
\end{beweis}

\begin{bemerkung}
Satz \vref{satz:Grenzwerte_und_Stetigkeit_von_Funktionen_S3_2} gilt analog f�r links-/rechtsseitige Grenzwerte.
\Solved{Referenziere zu Satz 3.2}{}
\end{bemerkung}

\begin{fdefinition}[Stetigkeit]
Eine Funktion $f$: $D \ra \R$ hei�t stetig in $x_0 \ra D$, wenn
\[\lim_{x \ra x_0} f(x) = f(x_0)\]
$f$ hei�t stetig in der Menge $A \subseteq D$, wenn $f$ in jedem Punkt $x_0 \in A$ stetig ist.
\Mark{Definition 3.3}
\end{fdefinition}

%\chapter{Wertannahme stetiger Funktionen}
\Mark{Chapter 3.2}
\begin{fsatz}
Die Funktion $f$: $D \ra \R$ sei stetig in $x_0$ und es gelte $f(x_0) > 0 (f(x_0)< 0)$. Dann gibt es eine Umgebung $U(x_0)$ von $x_0$ (\ac{d.h.} $\exists \delta > 0 U_{\mathbf{\delta}} (x_0) (x_0 - \delta, x_0 + \delta)$) so dass $\forall x \in U_{\delta} (x_0) \cap D$: $f(x) > 0$ ($f(x) < 0$).
\Mark{Satz 3.5}
\end{fsatz}

\begin{fbeweis}
Sei $f(x_0) > 0$ und $f$ in $x_0$ stetig. Nach \ac{Def.} der Stetigkeit gilt $\forall \epsilon > 0$ $\exists \delta (\epsilon, x_0) > 0$ so dass $\ub{\forall x \in D \tx{ mit } \betrag{x - x_0} < \delta}{\Lra x \in U_{\delta} (x_0)}$ \Ra $\ub{\betrag{f(x) - f(x_0)} < \epsilon}{\Lra f(x_0) - \epsilon < f(x) < f(x_0) + \epsilon}$
\begin{align*}
&\epsilon \ceq \frac{1}{2} f(x_0) > 0\\
&f(x) > f(x_0) - \epsilon = f(x_0) - \frac{1}{2} f(x_0) = \rac{1}{2} f(x_0)
\end{align*} \qed
\end{fbeweis}

\begin{fsatz}[Zwischenwertsatz]
\label{satz:Wertannahme_stetiger_Funktionen_S3_6}
Die Funktion $f$: $[a, b] \ra \R$ sei auf dem ganzen Intervall $[a, b]$ stetig und $f(a) \mal f(b) < 0$. Dann gibt es ein $\xi \in (a, b)$ mit $f(\xi) = 0$
\Insert{MA-23.04.2009-IMG-1}
\Mark{Satz 3.6}
\end{fsatz}

\begin{fbeweis}[Konstruktiver Beweis]
Sei \ac{OBdA} $f(a) < 0$, $f(b) > 0$\\
Folge von Intervallen $([x_n, y_n])_{n \in \N}$
\begin{enumerate}
\item $[x_n, y_n] \subset [x_{n - 1}, y_{n - 1}]$
\item $y_n - x_n = 2^{-n} (b -a)$
\item $f(x_n) \klgl 0$, $f(y_n > 0)$
\end{enumerate}
\begin{description}
\item[(IA)] $[x_0, y_0] = [a, b]$ erf�llt Bedingungen 2/3
\item[(IS)] $[x_n, y_n]$ sei bereits definiert und erf�lle die drei Bedingungen\\
		\[m \ceq \frac{x_n + y_n}{2}\]
		\begin{description}
		\item[Fall 1 $f(m) \klgl 0$] $[x_{n + 1}, y_n{n + 1}] = [m, y_n]$
		\item[Fall 2 $f(m) > 0$] $[x_{n + 1}, y_{n + 1}] = [x_n, m]$
		\end{description}
		Die drei Bedingungen sind f�r beide F�lle erf�llt.
		\begin{itemize}
		\item $(x_n)_{n \in \N_0}$ ist eine monoton wachsende und durch $b$ nach oben beschr�nkte Folge.
				\begin{itemize}[label=\Ra]
				\item $(x_n)_{n \in \N_0}$ ist konvergent
				\end{itemize}

		\item $(y_n)_{n \in \N_0}$ ist eine monoton fallende und durch $a$ nach unten beschr�nkte Folge.
				\begin{itemize}[label=\Ra]
				\item $(x_n)_{n \in \N_0}$ ist konvergent
				\end{itemize}
		\end{itemize}

		wegen 2) gilt:
		\begin{align*}
		&\lim_{n \ra \un} (y_n - x_n) = \lim_{n \ra \un} 2^{-n} (b - a) = 0\\
		=& \lim_{n \ra \un} y_n - \lim_{n \ra \un} x_n\\
		\Ra & \lim_{n \ra \un} (y_n) = \lim_{n \ra \un} (x_n) \eqc \xi
		\end{align*}
		Da $f$ stetig ist gilt:
		\begin{align*}
		\ub{\lim_{n \ra \infty} f(y_n)}{\ub{\lim_{y \searrow \xi} f(y)}{\grgl 0}} = f(\xi) = \ub{lim_{n \ra \infty} f(x_n)}{\ub{\lim_{x \nearrow \xi} f(x)}{\klgl 0}}
		\Ra f(\xi) = 0
		\end{align*} \qed
\end{description}
\end{fbeweis}

\begin{beispiel}
$f(x) = x^3 - 1$, $[a, b] = [0, 2]$
\begin{enumerate}[label=\arabic*), start=0]
\item $[x_0, y_0] = [0, 2]$
\item $f(1) = -1 \stack{\tx{Fall 1}}{\Ra} [x_1, y_1] = [1, 2]$
\item $f \rkl{\frac{3}{2}} = \frac{27}{8} - \frac{16}{8} = \frac{11}{8} > 0 \stack{\tx{Fall 2}}{\Ra} [x_2, y_2] = \eklamm{1, \frac{3}{2}}$
\item $f \rkl{\frac{5}{4}} = \frac{125}{64} - \frac{128}{64} = - \frac{3}{64} < 0 \stack{\tx{Fall 1}}{\Ra} \eklamm{\frac{5}{4}, \frac{3}{2}}$\\
		$\vdots$
\end{enumerate}
$M = \frac{1 + 2}{2}$
\end{beispiel}

\begin{bemerkung}
\mbox{}\par
\begin{enumerate}
\item Das Konstruktionsverfahren im obigen Beweis nennt man auch Intervallhalbierungsmethode \ac{bzw.} Bisektionsmethode.
\item Der Satz \vref{satz:Wertannahme_stetiger_Funktionen_S3_6} Kann erweitert werden zur Aussage, dass jeder Wert zwischen $f(a)$ und $f(b)$ angenommen wird.
		\Insert{MA-23.04.2009-IMG-2}
\end{enumerate}
\end{bemerkung}

\section{Aufgabe 3.3}
Zeigen Sie:
\begin{enumerate}[label=\alph*)]
\item Es gibt mindestens einen Schnittpunkt der beiden Funktionen $f$: $[0, \un) \ra \R$, $x \mapsto f(x) = x^2$ und $g$: $[0, \un) \ra \R$, $x \mapsto g(x) = 2$.
\item Es gibt mindestens eine L�sung der Gleichung $x^2 = \cosx{}$.
\end{enumerate}
\Todo{Ref zur L�sung}

\begin{fsatz}[Existenz von Maximum/Minimum]
Sei $f$: $D \ra \R$ stetig. Dann nimmt $f$ auf einem Intervall $[a, b] \subseteq D$ sein Maximum und sein Minimum an.\\
Daher $\exists \underline{x}, \overline{x} \in [a, b]$ mit\\
\[\forall x \in [a, b]:~ f(x) \klgl f(x) \klgl f(\overline{x})\]
\Mark{Satz 3.7}
\end{fsatz}

\begin{fbeweis}
-
\end{fbeweis}

\begin{bemerkung}
\mbox{}\par
\begin{enumerate}
\item Es ist wichtig, dass das Intervall abgeschlossen ist.
		\begin{beispiel}[F�r Probleme bei offenen Intervallen]
		\begin{align*}
		f:~(0, &1) \ra &&\R\\
		&x \mapsto &&\frac{1}{x}
		\end{align*}
		\Insert{MA-23.04.2009-IMG-3}
		\end{beispiel}

\item Es kann auch mehrere oder gar $\un$-viele Maxima/Minima geben
		\begin{beispiel}
		\begin{align*}
		f:~(0, &4\pi) \ra &&\R\\
		&x \mapsto &&\cosx{}
		\end{align*}
		\Insert{MA-23.04.2009-IMG-4}
		\[\overline{x} \in \gklamm{0, 2\pi, 4\pi}, \underline{x} \in \gklamm{\pi, 3\pi}\]
		\end{beispiel}

		\begin{beispiel}
		\begin{align*}
		f:~(-1, &2) \ra &&\R\\
		&x \mapsto &&f(x) = \begin{cases}0 & \tx{f�r } -1 \klgl x \klgl 0\\x & \tx{f�r } 0 < x < 1\\1 & \tx{f�r } 1 \klgl x \klgl 2\end{cases}
		\end{align*}
		\end{beispiel}
		\Insert{MA-23.04.2009-IMG-5}

\item Falls $f$ zudem differenzierbar ist, kann man die Orte der Maxima/Minima bestimmen
\end{enumerate}
\end{bemerkung}

\section{L�sungen}
\subsection{Aufgabe 3.3}
\begin{enumerate}[label=\alph*)]
\item $\exists$ Schnittpunkt von $f(x) = x^2$ und $g(x) = 2$ in $[0, \un)$\\
		\Lra $\exists x \in [0, \un)$: $f(x) = 2$
		\begin{description}
		\item[Variante 1] Anpassen des Beweises von Satz \vref{satz:Wertannahme_stetiger_Funktionen_S3_6} auf den Zwischenwert $2$ statt $0$.\\
				Folge von Intervallen mit $1, 2$ wie gehabt.\\
				\begin{enumerate}[label=\arabic*), start=3]
				\item $f(x_n) \klgl 2$, $f(y_n) > 2$
				\begin{description}
				\item[(IA)] $[x_0, y_0] = [0, 2]$ $f(0) = 0 < 2$, $f(2) = 4 > 2$
				\item[(IS)] \dots
						\begin{description}
						\item[Fall 1] $f(m) \klgl 2$
						\item[Fall 2] $f(m) > 2$
						\end{description}
				\end{description}
				\end{enumerate}

		\item[Variante 2] \ac{z.z.} $\exists x \in [0, \un)$ mit $x^2 = 2$ ($\Lra x^2 - 2 = 0$)
				\[h(x) = x^2 - 2\]
				da $h(0) = -2 < 0$, $h(2) = 2 > 0$, $h$ stetig besagt das Zwischenwertsatz, dass ein $\xi \in [0, 2]$ exisitiert mit $h(\xi) = 0$ ($\Lra f(\xi) = g(\xi) = 2$)
		\end{description}

\item $\exists x \in \R$: $x^2 = \cosx{}$ ($\Lra x^2 - \cosx{} = 0$)
		\begin{align*}
		&h(x) = x^2 - \cosx{}, [a, b] = \eklamm{0, \frac{\pi}{2}}\\
		&h(0) = -1 < 0, h \rkl{\frac{\pi}{2}} = \frac{\pi^2}{4} > 0
		\end{align*}
		Zwischenwertsatz auf $h$ anwenden.
\end{enumerate}
\Todo{Ref zur Aufgabe}

%
\part{Differentialrechnung}
\Mark{Part 4}
\chapter{Differenzierbarkeit}
\Mark{Chapter 4.1}
$f$: $D \ra \R$ stetig dann bedeutet dies, dass $f$ in $D$ keine \ggq{Spr�nge} hat\\
$f$: $D \ra \R$ differenzierbar bedeutet, dass $f$ in $D$ keine Spr�nge und keine \ggq{Knicke} hat

\begin{fdefinition}[Differenzierbarkeit]
Die Funktion $f$: $D \ra \R$ hei�t im Punkt $x_0 \in D$ \indexb{differenzierbar}, falls der Grenzwert
\[f' (x_0) \ceq \lim_{x \ra x_0} \frac{f(x) - f(x_0)}{x - x_0}\]
existiert $(f' (x_0) \in \R)$.\\
Die Zahl $f'(x_0)$ hei�t dann \indexb{Ableitung} von $f$ an der Stelle $x_0$.\\
Die Funktion $f$: $D \ra \R$ hei�t differenzierbar in $A \subseteq D$, falls $f$ in jedem $x_0 \in A$ differenzierbar ist.
\Mark{Definition 4.1}
\end{fdefinition}
\section*{Geometrische Interpretation}
\Insert{MA-28.04.2009-1}
Mit $h = \Delta x$ stellt der Differenzenquotient
\[\frac{f(x_0 + \Delta x) - f(x_0)}{\delta x} = \frac{f(x_0 + h) - f(x_0)}{h} = \frac{\Delta y}{\Delta x} = \frac{\Delta f}{\Delta x}\]
die Steigung der Sekante (Sehne) zwischen $f(x_0)$ und $f(x_0 + \Delta x)$ dar.\\
F�r $h \ra 0$ \ac{bzw.} $\Delta x \ra 0$ geht die Sekante in die Tangente in $x_0$ �ber und Steigung der Tangente ist
\begin{align*}
f'(x_0) &= \frac{df(x)}{dx}_{x = x_0} = \frac{df}{dx}_{(x_0)} = \lim_{\Delta x \ra 0} \frac{f(x + \Delta x) - f(x_0)}{\Delta x}\\
&=\tanx{\alpha}
\end{align*}

\begin{bemerkung}
\mbox{}\par
\begin{enumerate}
\item In Anlehnung an d en Differenzenquotient $\frac{\Delta f}{\Delta x}$ nennt man $\frac{df}{dx}$ \indexb{Differenzialquotient}.
\item $h$ \ac{bzw.} $\Delta x$ muss nicht $> 0$ sein.
\item Alternative Definition der Differenzierbarkeit\\
		Sei $D \subseteq \R$ und $x_0 \in D$, so dass mindestens eine Folge $(x_n)_{n \in \N}$ mit $x_n \in D \backslash \gklamm{x_0}$ existiert mit $\lim_{n \ra \un} (x_n) = x_0$.\\
		Eine Funktion $f$: $D \ra \R$ hei�t differenzierbar in $x_0 \in D$, falls es eine Konstante $a \in \R$ und eine Funktion $r$: $D \ra \R$ mit $\lim_{x \ra x_0} r(x) = 0$ gibt, sodass
		\[f(x) = f(x_0) + a \mal (x - x_0) + r(x) \mal (x - x_0)~~x \in D\]
		In diesem Fall ist $a = f' (x_0)$

\item Die Gleichung der Tangente an $f$ in $x_0$ lautet
		\[\ub{T(f, x)}{= T(x) = T_f(x)} = f(x_0) + f' (x_0) \mal (x - x_0)\]

\item aus 3, 4 folgt, dass sich lokal (also in der N�he eines gegebenen Punktes $x_0$) der \ggq{Zuwachs} der Funktion $f$ durch den Zuwachs der Tangente $T_f$ ($a_n$ $f in x_0$) darstellen l�sst.
		\begin{align*}
		\Delta f &= f(x) - f(x_0)\\
		&= f(x_0) + f'(x_0) (x - x_0) + r(x) \mal (x - x_0) - f(x_0)\\
		&= f'(x_0) \mal (x - x_0) + r(x) \mal (x - x_0)\\
		&\approx f'(x_0) \mal (x - x_0)\\
		&= T_f(x) - T_f(x_0)\\
		\Ra& df = f'(x) dx \tx{ Differential}
		\end{align*}
		\ac{d.h.} wenn man um die Infinitesimale Gr��e $dx$ von $x_0$ weggeht, �ndert sich der Funktionswert um $df = f'(x_0) dx$
\end{enumerate}
\end{bemerkung}

\begin{beispiel}
\mbox{}\par
\begin{enumerate}
\item $f(x) = 4x - 5$\\
		Ist $f$ in $x_0 = 3$ differenzierbar?\\
		Ist $f$ in $\R$ differenzierbar?
		\begin{align*}
		\lim_{x \ra x_0} \frac{f(x) - f(x_0)}{x - x_0} &= \lim_{h \ra 0} \frac{f(x_0 + h) - f(x_0)}{h}\\
		&= \lim_{h \ra 0} \frac{\rkl{4 (x_0 + h) - 5} - (4 x_0 - 5)}{h}\\
		&= \lim_{h \ra 0} \frac{4x_0 + 4h - 5 - 4x_0 + 5}{h} = \lim_{h \ra 0} 4 = 4
		\end{align*}
		Da der Grenzwert existiert ohne dass wir den Wert von $x_0$ gebraucht h�tten, ist $f$ in ganz $\R$ differenzierbar.
\end{enumerate}
\end{beispiel}

\section{Aufgabe 4.1}
Ermitteln sie die Ableitung der Funktion $f(x) = x^2$ direkt mit Hilfe der Definition\\
\Todo{Link zu L�sung}
\section{L�sungen}
\subsection{Aufgabe 4.1}
\Todo{Link zur Aufgabe}
$f(x) = x^2$ gesucht $f'(x)$
\begin{align*}
\lim_{x \a x_0} \frac{f(x) - f(x_0)}{x - x_0} &= \lim_{h \ra 0} \frac{f(x_0 + h) - f(x_0)}{h}\\
&= \lim_{h \ra 0} \frac{(x_0 + h)^2 - x_0^2}{h} = \lim_{h \ra 0} \frac{x_0^2 + 2x_0 h + h^2 - x_0^2}{h}\\
&= \lim_{h \ra 0} (2 x_0 + h) = 2 x_0 = f'(x_0)\\
\Ra & f'(x) = 2x
\end{align*}

\begin{beispiel}
\mbox{}\par
\begin{enumerate}[start=2]
\item $f(x) = e^x$ gesucht $f'(x)$
		\begin{align*}
		\lim_{x \ra x_0} \frac{f(x) - f(x_0)}{x - x_0} &= \lim_{h \ra 0} \frac{f (x_0 + h) - f(x_0)}{h}\\
		&= \lim_{h \ra 0} \frac{e^{x_0 + h} - e^{x_0}}{h} = \lim_{h \ra 0} \frac{e^{x_0} e^h - e^{x_0}}{h}\\
		&= e^{x_0} \mal \lim_{h \ra 0} \frac{e^h - 1}{h} = e^{x_0} \lim_{h \ra 0} \frac{1}{h} \rkl{\sum_{k = 0}^{\un} \frac{h^k}{k!} - 1}\\
		&= e^{x_0} \mal \lim_{h \ra 0} \rkl{\frac{1}{h} \sum_{k = 1}^{\un} \frac{h^k}{k!}} = e^{x_0} \ub{\lim_{h \ra 0} \sum_{k = 1}^{\un} \frac{h^{k - 1}}{k!}}{= \textcircled{\star}}\\
		\textcircled{\start} &\klgl \lim_{h \ra 0} \sum_{k = 1}^{\un} \frac{\betrag{h}^{k - 1}}{k!} = \lim_{h \ra 0} \sum_{k = 0}^{\un} \frac{\betrag{h}^k}{(k + 1)!} \klgl \lim_{h \ra 0} \sum_{k = 0}^{\un} \frac{\betrag{h}^k}{k!}\\
		&= \lim_{h \ra 0} e^{\betrag{h}} = 1\\
		\textcircled{\start} &= \lim_{h \ra 0} \frac{h^k}{(k + 1)!} = \lim_{h \ra 0} \rkl{1 - \sum_{k = 1}^{\un} \frac{\betrag{h}^k}{(k + 1)!}}\\
		&= \lim_{h \ra 0} \rkl{1 - \sum_{k = 1}^{\infty} \frac{\betrag{h}}{2} \mal \frac{\betrag{h}}{3} \mal \frac{\betrag{h}}{4} \dots \frac{\betrag{h}}{(k + 1)}} \grgl \lim_{h \ra \infty} \rkl{1 - \sum_{k = 1}^{\infty} \rkl{\frac{\betrag{h}}{2}}^k}\\
		&= \lim_{h \ra 0} \rkl{1 - \frac{\betrag{h}}{2} \mal \sum_{k = 0}^{\un} \rkl{\frac{\betrag{h}}{2}}^k}\\
		&= \lim_{h \ra 0} \rkl{1 - \frac{\betrag{h}}{2} \mal \frac{1}{1 - \frac{\betrag{h}}{2}}}\\
		&= 1 - \lim_{h \ra 0} \frac{\betrag{h}}{2} \mal \lim_{h \ra 0} \frac{1}{1 - \frac{\betrag{h}}{2} = 1 - 0 \mal 1 = 1\\
		\Ra& \frac{d}{dx} \rkl{e^x}_{x = x_0} = e^{x_0} = f'(x_0}\\
		\Ra& f'(x) = e^x
		\end{align*}
\end{enumerate}
\end{beispiel}
%\part{�bungsbl�tter}
%\chapter{Erg�nzungsbl�tter}
\section{Erg�nzungsblatt 03}
\subsection{Aufgabe 12}

\subsection{Aufgabe 13}
\subsubsection{L�sung}
\renewcommand{\labelenumi}{\alph{enumi})}
\begin{enumerate}
\item 
\item 
\end{enumerate}
\renewcommand{\labelenumi}{\arabic{enumi}.}

\subsection{Aufgabe 14}

\subsection{Aufgabe 15}

\subsection{Aufgabe 16}

\subsection{Aufgabe 17}
\subsubsection{Aufgabenstellung}
Seien $A, B$ beliebige Mengen, Beweisen oder widerlegen sie die folgenden Aussagen:
\renewcommand{\labelenumi}{\alph{enumi})}
\begin{enumerate}
\item $\mathscr{P}(A \cup B) \subseteq (\mathscr{P}(A) \cup \mathscr{P}(B))$
\item $\mathscr{P}(A \cup B) = (\mathscr{P}(A) \cup \mathscr{P}(B))$
\item $\mathscr{P}(A \cup B) \supseteq (\mathscr{P}(A) \cup \mathscr{P}(B))$
\item $\mathscr{P}(A \times B) \subseteq (\mathscr{P}(A) \times \mathscr{P}(B))$
\end{enumerate}
\renewcommand{\labelenumi}{\arabic{enumi}.}

\subsubsection{Aufgaben L�sung}
\renewcommand{\labelenumi}{zu \alph{enumi})}
\begin{enumerate}
\item TODO
\item TODO
\item TODO
\item TODO
\end{enumerate}
\renewcommand{\labelenumi}{\arabic{enumi}.}

\chapter{Aufgaben}
\section{E27}
\subsection{Aufgabenstellung}
$\mathbb{P}_2 = \gklamm{a t^2 + bt + c | a, b, c \in \mathbb{R}} B = \gklamm{1, t, t^2}$
\begin{enumerate}[label=\alph*)]
\item $p_1(t) = 1 - 3 t^1 \Ra (p_1)_{\mathcal{B}} = \vektor{1}{0}{-3}$\\
		$p_2(t) = 2 + t - 5 t^2 \Ra (p_2)_{\mathcal{B}} = \vektor{2}{1}{-5}$\\
		$p_3(t) = 1 + 2t \Ra (p_3)_{\mathcal{B}} = \vektor{1}{2}{0}$\\
\item zu zeigen $B = \gklamm{p_1(t), p_2(t), p_3(t)}$ ist eine Basis von $\mathbb{P}_2$ es reicht zu zeigen $p_1(t), p_2(t), p_3(t)$ linear unabh�ngig $\Leftrightarrow$ zu zeigen $(p_1)_{\mathcal{B}}, (p_2)_{\mathcal{B}}, (p_3)_{\mathcal{B}}$ linear unabh�ngig
\end{enumerate}

\subsection{Aufgabenl�sung}
$\gauss{ccc|c}{1 & 2 & 1 & 0\\0 & 1 & 2 & 0\\-3 & -5 & 0 & 0} \underrightarrow{\RM{3} = \RM{3} + 3 \RM{1}}$\\
$\underrightarrow{\RM{3} = \RM{3} - \RM{2}}$\\
$\underrightarrow{\RM{1} = \RM{1} - \RM{3}}_{\RM{2} = \RM{2} - 2 \RM{3}}$\\
$\underrightarrow{\RM{1} = \RM{1} - 2 \RM{2}}$\\
Alternative zu zeigen $\begin{vmatrix}1 & 2 & 1\\0 & 1 & 2\\-3 & -5 & 0\end{vmatrix} \neq 0$\\
Alternativ zeigen $1, t, t^2$ ist Linearkombination von $p_1(t), p_2(t), p_3(t)$

\section{Aufgabe E28}
\begin{enumerate}
\item $\mathbb{R}^3$
\item $\mathbb{P}_2 (\mathbb{R})$
\item $v_1 = \spanx{\gklamm{\varvektor{c}{1\\0\\0\\0}, \varvektor{c}{0\\1\\0\\0}, \varvektor{c}{0\\0\\1\\0}}}$
\item $v_2 = \gklamm{\varvektor{cc}{a & b\\c & 0} \vert a, b, c \in \mathbb{R}}$
\item $v_3 = \gklamm{ax^3 + bx^5 + cx^6 \vert a, b, c \in \mathbb{R}}$
\end{enumerate}

\section{Aufgabe E29}
$\vec{x} \in H \Lra es ex s,t \in \mathbb{R} \mal s \vec{v_1} + t \vec{v_2} = \vec{x}$
\[\gauss{cc|c}{3 & -1 & 3\\6 & 0 & 12\\2 & 1 & 7} \underrightarrow{\RM{1} \lra \RM{2}} \gauss{cc|c}{6 & 0 & 12\\3 & -1 & 3\\2 & 1 & 7}\]
\[\underrightarrow{\RM{1} = \RM{1} \mal \frac{1}{6}} \gauss{cc|c}{1 & 0 & 2\\3 & -1 & 3\\2 & 1 & 7}\]
\[\underrightarrow{\RM{2} = \RM{2} - 3 \RM{1}}_{\RM{3} = \RM{3} - 2 \RM{2}} \gauss{cc|c}{1 & 0 & 2\\0 & -1 & -3\\0 & 1 & 3} \underrightarrow{\RM{3} = \RM{3} + \RM{2}} \gauss{cc|c}{1 & 0 & 2\\0 & -1 & -3\\0 & 0 & 0}\]
\[\ra \gauss{cc|c}{1 & 0 & 2\\0 & 1 & 3\\0 & 0 & 0}\]
\[\Ra \vec{x} \in H, \vec{x_B} = \varvektor{c}{2\\3}\]

\section{Aufgabe E30}
8

%\chapter{�bungsbl�tter}
\section{L�sungen}
\subsection{�bung 1.1}
\label{sec:Uebungsblaetter_1_1L}
%L�sung zu Aufgabe \vref{sec:Uebungsblaetter_1_1}.

\begin{itemize}
\item 3 Wagen 1. Klasse
\item 5 Wagen 2. Klasse
\item 2 Gep�ckwagen
\end{itemize}

\begin{enumerate}[label=\alph*)]
\item $\binom{10}{3} \mal \binom{7}{5} \mal \binom{2}{2} = \frac{10 \mal 9 \mal 8}{3!} \mal \frac{7 \mal 6 \mal 5 \mal 4 \mal 3}{5!} \mal \frac{2 \mal 1}{2!} = \frac{10!}{3! \mal 5! \mal 2!}$\\
		$\binom{10}{3}$: Anzahl m�glicher Pl�tze f�r 1. Klasse auszusuchen\\
		$\binom{7}{5}$: Anzahl M�glichkeiten im Anschluss Pl�tze f�r 2. Klasse festzulegen\\
		$\binom{2}{2}$: Letzte 2 Pl�tze im Gep�ckwagen

\item Wagen wie in a), alle 2. Klasse Wagen am St�ck

		%MA-02.04.2009-IMG-1
		\begin{tikzpicture}[decoration=brace]
		\draw[thick, ->] (0,0) -- (6,0) node[below] {$n$};
		\draw[-|] (2,-0.2) node[below] {$n - 1$} -- (2,1.8) node[left]{$S_{n - 1}$};
		\draw[-|] (3,-0.2) node[below] {$n$} -- (3,2.8) node[above]{$S_1$};
		\draw[dashed] (2,1.8) -- (5,1.8);
		\draw[dashed] (3,2.8) -- (5,2.8);
		\draw[decorate] (2.5,2.8) -- (2.5,1.8) node[right, midway] {$\epsilon$};
		\draw (3,2.8) -- (2.5,2.8);
		
		\draw[draw=red, <->] (5,2.8) node[right] {$S$} -- (5,1.8) node[right,midway] {$> \frac{\epsilon}{2}$};
		\draw[draw=blue, <->] (6,0) node[right] {$S$} -- (6,2.8) node[right,midway] {$> \frac{\epsilon}{2}$};
		\end{tikzpicture}

		6 M�glichkeiten den Platz des ersten 2. Klasse Wagens festzulegen\\
		$\binom{5}{3} \mal \binom{2}{2}$ Anzahl M�glichkeiten die Pl�tze f�r die anderen Wagen festzulegen\\
		$6 \mal \binom{5}{3} \mal \binom{2}{2} = 6 \mal \frac{5 \mal 4 \mal 3}{1 \mal 2 \mal 3} = 60$
\end{enumerate}

\subsection{�bung 1.2 (Abwandlung)}
\label{sec:Uebungsblaetter_1_2AL}
%L�sung zu Aufgabe \vref{sec:Uebungsblaetter_1_2A}.

\begin{itemize}
\item 12 stellige Zahl mit Ziffern 1-9
\item max zwei 9
\item Summe der ersten beiden 16
\end{itemize}

\subsubsection*{ersten beiden Stellen}
$\left.88\right\}$ \textcircled{1}\\
$\left.\begin{array}{l}
79\\
97
\end{array}\right\}$ \textcircled{2}

\begin{enumerate}[label=\textcircled{\arabic*}]
\item M�glichkeiten noch $2 \times 9$, $1 \times 9$, $0 \times 9$ in den 10 Stellen\\
		$\binom{10}{2} \mal 8 + \binom{10}{1} \mal 8^9 + \binom{10}{0} \mal 8^{10}$\\
		$8$ ist die Anzahl der M�glichkeiten pro Stelle\\
		$\space^8$ ist die Anzahl der Stellen

\item M�glichkeiten noch $1 \times 9$ oder $0 \times 9$ in 10 Stellen\\
		$\binom{10}{1} \mal 8^9 + \binom{10}{0} \mal 8^{10}$\\
		\Ra Totale Anzahl = $1 \mal \rkl{\rkl{\binom{10}{2} \mal 8^8 + \binom{10}{1} \mal 8^9 + \binom{10}{0} \mal 8^{10}} + 2 \mal \rkl{\binom{10}{1} \mal 8^9 + \binom{10}{0} \mal 8^{10}}}$
\item 
\end{enumerate}

\subsection{�bung 1.3}
\label{sec:Uebungsblaetter_1_3L}
%L�sung zu Aufgabe \vref{sec:Uebungsblaetter_1_3}.

zu zeigen $\forall n \in \N_0: \sum_{k = 0}^n \binom{m + k}{k} = \binom{m + n + 1}{n}$
\begin{description}
\item[(IA)] $n = 0$
		\[\left.\begin{array}{l}
		\tx{LS} = \sum_{k = 0}^0 \binom{m + k}{k} = \binom{m + 0}{0} = 1\\
		\tx{RS} = \binom{m + 0 + 1}{0} = 1
		\end{array}\right\} \checkmark\]

\item[(IS)] zu zeigen $\forall n \in \N: \underbrace{\sum_{k = 0}^n \binom{m + k}{k} = \binom{m + n + 1}{n}}_{= \tx{(IV)}}$ \Ra $\sum_{k = 0}^{n + 1} \binom{m + k}{k} = \binom{m + n + 2}{n + 1}$
		\begin{align*}
		&\sum_{k = 0}^{n + 1} \binom{m + k}{k} = \binom{m + n + 1}{n + 1} + \sum_{k = 0}^n \binom{m + k}{k}\\
		\stack{=}{\tx{(IV)}} & \underbrace{\binom{m + n + 1}{n + 1}}_{\binom{m + n + 1}{m}} + \binom{m + n + 1}{n} \underbrace{=}_{\tx{Satz 1.3 iii)}} \binom{m + n + 2}{n + 1}
		\end{align*}\qed
\end{description}

\subsection{�bung 1.4}
\label{sec:Uebungsblaetter_1_4L}
%L�sung zu Aufgabe \vref{sec:Uebungsblaetter_1_4}.

F�r welche $(x, y) \in \mb{R}^2$ gilt:
\[xy \klgl x^2 + y^2\]
\textit{Behauptung:} Gilt f�r alle $(x, y) \in \mb{R}^2$
\begin{description}
\item[Fall 1] $xy \klgl 0$ (\ac{d.h.} $(x \klgl 0 \und y \grgl 0) \oder (x \klgl 0 \und y \grgl 0)$\\
		\begin{align*}
		&x y \klgl 0 \tx{ und } \underbrace{x^2}_{\grgl 0} + \underbrace{y^2}_{\grgl 0} \grgl 0\\
		\Ra &xy \klgl x^2 + y^2
		\end{align*}

\item[Fall 2] $xy > 0 (\Lra (x > 0 \und y > 0) \oder (x < 0 \und y < 0)$
		\begin{align*}
		&xy \stack{\klgl}{!} x^2 + y^2~~~~\vert - xy\\
		\Lra &0 \stack{\klgl}{!} x^2 - xy + y^2
		\La &x^2 - xy + y^2 > x^2 - x - \underbrace{xy}_{> 0} + y^2 = (x - y)^2 \grgl 0
		\end{align*}\qed
\end{description}

\subsection{�bung 1.5}
\label{sec:Uebungsblaetter_1_5L}
%L�sung zu Aufgabe \vref{sec:Uebungsblaetter_1_5}.

$\betrag{\frac{3 - 2x}{2 + x}} \klgl 4$ $x \in \mb{R} \backslash \gklamm{-1}$\\
\Lra $-4 \klgl \frac{3 - 2x}{1 + x} \klgl 4$
\begin{description}
\item[Fall 1] $1 + x > 0 (\Lra x > - 1)$
		\begin{align*}
		\Lra &-4(1 + x) \klgl 3 - 2x \klgl 4(1 + x)\\
		\Lra &-4 -4x \klgl 3 - 2x \klgl 4 + 4x\\
		\Lra &\matrixp{-4 -4x \klgl 3 - 2x \vert + 4x - 3\\\Lra - 7 \klgl 2x \vert : 2\\\Lra -\frac{7}{2} \klgl x} \und \matrixp{3 - 2x \klgl 4 + 4x \vert +2x - 4\\- 1 \klgl 6x \vert :6\\-\frac{1}{6} \klgl x}\\
		&L_1 = \left[- \frac{1}{6}, \infty\right)
		\end{align*}

\item[Fall 2] $1 + x < 0$ ($\Lra x < - 1$)
		\begin{align*}
		\Lra &-4(1 + x) \grgl 3 - 2x \grgl 4(1 + x)\\
		\Lra &-4 -4x \grgl 3 - 2x \grgl 4 + 4x\\
		\Lra &\matrixp{-4 -4x \grgl 3 - 2x \vert + 4x - 3\\\Lra - 7 \grgl 2x \vert : 2\\\Lra -\frac{7}{2} \grgl x} \und \matrixp{3 - 2x \grgl 4 + 4x \vert +2x - 4\\- 1 \grgl 6x \vert :6\\-\frac{1}{6} \grgl x}\\
		&L_2 = \left(\infty, -\frac{7}{2}\right]
		\end{align*}
\end{description}
\[L = L_1 \cup L_2 ) \left(-\infty, -\frac{7}{2}\right] \cup \left[-\frac{1}{6}, \infty\right)\]

\subsection{�bung 2.1}
$a_{n + 1} = \frac{1}{2} \rkl{a_n + \frac{x_0}{a_n}} ~~ n \in \N$
\begin{enumerate}[label=\alph*)]
\item \ac{z.z.} $a_n > \sqrt{x_0}$
		\begin{description}
		\item[(IA)] $n = 1$ $a_1 > \sqrt{x_0}$ erf�llt nach Voraussetzung
		\item[(IS)] \ac{z.z.} $a_n > \sqrt{x_0} \Ra a_{n + 1} > \sqrt{x_0}$
				\begin{align*}
				a_{n + 1} &= \frac{1}{2} \rkl{a_n + \frac{x_0}{a_n}} = \ub{\frac{1}{2a_n}}{> 0} \rkl{a_n^2 + x_0}\\
				&= \frac{1}{2a_n} \rkl{a_n^2 - 2a_n \sqrt{x_0} + \rkl{\sqrt{x_0}}^2 + 2a_n \mal \sqrt{x_0}}\\
				&= \frac{1}{2a_n} \rkl{a_n^2 - 2a_n \sqrt{x_0} + \rkl{\sqrt{x_0}}^2} + \sqrt{x_0}\\
				&= \ub{\frac{1}{2a_n}}{>0} \ub{\rkl{\ub{a_n - \sqrt{x_0}}{> 0}}^2} + \sqrt{x_0} > \sqrt{x_0}\\
				\Ra& (a_n)_{n \in \N} \tx{ monoton fallend}
				\end{align*}
		\end{description}

\item \ac{z.z.} $(a_n)_{n \in \N}$ monoton fallend
		\begin{align*}
		a_{n + 1} - a_n &= \frac{1}{2} \rkl{a_n + \frac{x_0}{a_n}} - a_n = \frac{x_0}{2a_n} - \frac{1}{2} a_n\\
		&= \frac{x_0 - a_n^2}{2a_n} = \ub{\frac{1}{2a_n}}{>0} \ub{\rkl{\sqrt{x_0} - a_n}}{<0} \mal \ub{\rkl{\sqrt{x_0} + a_n}}{>0} < 0
		\end{align*}

\item \ac{z.z.} $\lim_{n \ra \infty} a_n = \sqrt{x_0}$
		\begin{align*}
		&a = \lim_{n \ra \un} a_n = \lim_{n \ra \un} a_n + 1 = \lim_{n \ra \un} \rac{1}{2} \rkl{a_n + \frac{x_0}{a_n}} = \frac{1}{2} \rkl{a + \frac{x_0}{a}}\\
		\Lra& \rac{1}{2} a = \frac{x_0}{2a} \Vert a \mal 2\\
		\Lra& a^2 = x_0\\
		\stack{a > 0}{\Lra}& a = \sqrt{x_0}
		\end{align*}

\item $a_1 = 2$, $x_0 = 2$\\
		\[a_2 = \frac{1}{2} \rkl{a_1 + \frac{2}{a_1}} = \frac{1}{2} \rkl{2 + \frac{2}{2}} = \frac{3}{2}\]
		\[a_3 = \frac{1}{2} \rkl{a_2 + \rac{2}{a_2}} = \frac{1}{2} \rkl{\frac{3}{2} + \rac{2 \mal 2}{3}} = \rac{9 + 8}{2 \mal 6} = \frac{17}{12}\]
		\begin{align*}
		\betrag{a - \sqrt{2}} &= a_3 - \sqrt{2} = \frac{a_3^2 - 2}{a_3 + \sqrt{2}} \klgl \frac{a_3^2 - 2}{a_3}\\
		&= a_3 - \frac{2}{a_3} = \frac{17}{12} - \frac{2 \mal 12}{17} = \frac{289 - 288}{12 \mal 17} = \frac{1}{12 \mal 17}\\
		\Ra& \sqrt{2} \in \eklamm{\ub{\frac{17}{12}}{a_3} - \frac{1}{12 \mal 17}, \ub{\frac{17}{12}}{a_3}}
		\end{align*}
\end{enumerate}

\subsection{�bung 2.2}
\begin{enumerate}[label=\alph*)]
\item $a_n = n - \sqrt{n^2 - b}$
		\begin{align*}
		\lim_{n \ra \un} a_n &= \frac{\rkl{n - \sqrt{n^2 - 5n}} \mal \rkl{n + \sqrt{n^2 - 5n}}}{n + \sqrt{n^2 - 5n}}\\
		&= \lim_{n \ra \un} \frac{n^2 - \rkl{n^2 - 5n}}{n + \sqrt{n^2 - 5n}} = \lim_{n \ra \un} \frac{5n}{n + \sqrt{n^2 - 5n}}\\
		&= \lim_{n \ra \un} \frac{5}{1 + \sqrt{\frac{n^2}{n^2} - \frac{5n}{n^2}}} = \lim_{n \ra \un} \frac{5}{1 + \sqrt{1 - \frac{5}{n}}} = \frac{5}{2}
		\end{align*}

\item $b_n = \frac{\sinx{n}}{n}$
		\begin{align*}
		0 &\klgl \lim_{n \ra \un} \betrag{b_n} = \lim_{n \ra \un} \betrag{\frac{\sinx{n}}{n}} \klgl \lim_{n \ra \un} \frac{1}{n} = 0\\
		&\Ra \lim_{n \ra \un} b_n = 0
		\end{align*}

\item $c_n = n \rkl{1 - \rkl{1 - \frac{1}{n}}^{42}}$
		\begin{align*}
		\lim_{n \ra \un} n \rkl{1 - \rkl{1 - \frac{1}{n}}^{42}} &= \lim_{n \ra \un} n \rkl{1 - \sum_{k = 0}^{42} \binom{42}{k} \rkl{-\frac{1}{n}}^k}\\
		&= \lim_{n \ra \un} n \rkl{\sum_{k = 1}^{42} \binom{42}{k} \mal \rkl{-1}^{k - 1} \rkl{\frac{1}{n}}^k}\\
		&= \lim_{n \ra \un} \sum_{k = 1}^{42} \binom{42}{k} \mal \rkl{-1}^{k - 1} \rkl{\frac{1}{n}}^{k - 1}\\
		&= \binom{42}{1} + \sum_{k = 2}^{42} \mal (-1)^{k - 1} \mal \rkl{\ub{\lim_{n \ra \un} \rkl{\frac{1}{n}}}{=0}}^{k - 1} = \binom{42}{1} = 42
		\end{align*}

\item $d_n = \rkl{1 + \frac{1}{n}}^{n^2}$\\
		Behauptung: $(d_n)_{n \in \N}$ ist divergent (genauer $\lim_{n \ra \un} d_n = \un$)
		\begin{align*}
		&\rkl{1 + \frac{1}{2}}^{n^2} = \rkl{\rkl{1 + \frac{1}{n}}^n}^n \ub{\grgl}{\tx{f�r $n$ gen�gend gro�}} \rkl{\ub{\rkl{1 + \frac{1}{n}}^n}{\overrightarrow{n \ra \un} e}}^k\\
		\Ra& \lim_{n \ra \un} \rkl{1 + \frac{1}{n}}^{n^2} \ub{\grgl}{\tx{f�r $\forall k \in \N$}} \rkl{\lim_{n \ra \un} \rkl{1 + \frac{1}{n}}^n}^k = e^k\\
		&\forall K \in \R \exists k \in \N:~e^k > K
		\end{align*}
		\Ra $(d_n)$ ist nicht nach oben beschr�nkt und damit divergent

\item $e_n = \sqrt[n]{\sqrt[n]{n!}}$
		\begin{align*}
		&\lim_{n \ra \un} e_n = \lim_{n \ra \un} \sqrt[n]{\sqrt[n]{n!}} = \tx{\textcircled{$\star$}}\\
		&\tx{\textcircled{$\star$}} \grgl \lim_{n \ra \un} \sqrt[n]{\sqrt[n]{1}}  = \lim_{n \ra \un} \sqrt[n]{1} = 1\\
		&\tx{\textcircled{$\star$}} \klgl \lim_{n \ra \un} \sqrt[n]{\sqrt[n]{n^n}} = \lim_{n \ra \un} \sqrt[n]{n} = 1\\
		\Ra& \lim_{n \ra \un} \sqrt[n]{\sqrt[n]{n!}} = 1
		\end{align*}
\end{enumerate}

\subsection{�bung 2.3}
\begin{enumerate}[label=\alph*)]
\item gesucht $p, q \in \N$ mit $\frac{p}{qq} = 0,\overline{4711}$
		\begin{align*}
		0,\overline{4711} &= \frac{4711}{10^4} + \frac{4711}{10^8} + \frac{4711}{10^{12}} + \dots\\
		&= 4711 \sum_{k = 1}^{\un} \rkl{\frac{1}{10^4}}^k = 4711 \frac{1}{10^4} \sum_{k = 0}^{\un} \rkl{\frac{1}{10^4}}^k\\
		&= 4711 \frac{1}{10^4} \mal \frac{1}{1 - \frac{1}{10^4}} = 4711 \mal \frac{1}{10^4 - 1}\\
		&= \frac{4711}{9999} \Ra p = 4711, q = 9999
		\end{align*}

\item $\frac{p}{q} = 0,1230\overline{443}$
		\begin{align*}
		\frac{p}{q} &= 0,1230\overline{443} = \frac{123}{1000} + \frac{443}{10^7} + \frac{443}{10^{10}} + \frac{443}{10^{13}} + \dots\\
		&= \frac{123}{10^3} + \frac{443}{10^7} \sum_{k = 0}^{\un} \rkl{\frac{1}{10^3}}^k = \frac{123}{10^3} + \frac{443}{10^7} \mal \frac{1}{1 - \frac{1}{10^3}}\\
		&= \frac{123}{10^3} + \frac{443}{10^4} \mal \frac{1}{999} = \frac{123 \mal 9990 + 443}{9990000}\\
		&= \frac{1229213}{9990000}
		\end{align*}
\end{enumerate}

\subsection{�bung 2.4}
$a_n = $ In der $n$-ten Minute zur�ckgelegte Anteil am Gummiband\\
$a_1 = \frac{25cm}{100cm} = \frac{1}{4}$\\
$a_2 = \frac{25cm}{200cm} = \frac{1}{8}$\\
$a_3 = \frac{25cm}{300cm} = \frac{1}{12}$\\
$a_n = \frac{1}{4n}$
bis und mit $n$-ter Minute zur�ckgelegter Anteil
\[\sum_{n = 1}^{N} a_n = \sum_{n = 1}^{N} \frac{1}{4n} = \frac{1}{4} \ub{\sum_{n = 1}^{N} \frac{1}{n}}{\tx{harmonische Reihe}}\]
\Ra die Ameise erreicht das Ende
\[\frac{1}{4} \sum_{n = 1}^{N} \frac{1}{n} \grgl 1\]
kleinstes $N$, welches die Bedingung erf�llt\\
\Ra $N = 31min$

\subsection{�bung 2.5}
\begin{enumerate}[label=\alph*)]
\item $\sum_{h = 0}^{\un} \frac{m!}{m^m}$ konvergent oder divergent?\\
		Majorantenkriterium
		\[
		\sum_{m = 0}^{\un} \frac{m!}{m^m} = \sum_{m = 0}^{\un} \rkl{\ub{\frac{m}{m}}{=1} \mal \ub{\frac{(m - 1)}{m}}{<1} \mal \ub{\frac{3}{m}}{<1} \mal \frac{2}{m} \mal \frac{1}{m}} \klgl 1 + \ub{2 \mal \sum_{m = 1}^{\un} \frac{1}{m^2}}{\tx{konvergent}}
		\]

\item $\sum_{h = 0}^{\un} \frac{m!}{m^m}$ konvergent oder divergent?\\
		notwendiges Kriterium 
		\[
		\sum_{m = 0}^{\un} \frac{m!}{m^m} = \sum_{m = 0}^{\un} \rkl{\ub{\frac{m}{m}}{=1} \mal \ub{\frac{(m - 1)}{m}}{<1} \mal \ub{\frac{3}{m}}{<1} \mal \frac{2}{m} \mal \frac{1}{m}} \klgl 1 + \ub{2 \mal \sum_{m = 1}^{\un} \frac{1}{m^2}}{\tx{konvergent}}
		\]

\item $\sum_{n = 1}^{\infty} \frac{\sinx{n}}{n^2}$ konvergent/divergent?\\
		Majorantenkriterium
		\[\betrag{\sum_{n = 1}^{\un} \frac{\sinx{x}}{n^2}} \klgl \sum_{n = 1}^{\un} \betrag{\frac{\sinx{n}}{n^2}} \klgl \sum_{n = 1}^{\un} \frac{1}{n^2}\]
		verallgemeinerte harmonische Reihe mit\\
		$\alpha = 2$ \Ra konvergent
\end{enumerate}

\subsection{�bung 2.6}
\begin{enumerate}[label=\alph*)]
\item $\sum_{n = 0}^{\un} \ub{\rkl{\frac{1}{2^n} + \frac{1}{3^n}}}{= a_n} \rkl{x - \ub{0}{= x_0}}^n$\\
		Quotientenmethode
		\begin{align*}
		\lim_{n \ra \un} \betrag{\frac{a_n}{a_{n + 1}}} &= \lim_{n \ra \un} \rkl{\rkl{\frac{1}{2^n}} + \frac{1}{3^n}} \mal \rkl{\frac{1}{2^{n + 1}} + \frac{1}{3^{n + 1}}^{-1}}\\
		&= \lim_{n \ra \un} \frac{\frac{1}{2^n} + \frac{1}{3^n}}{\frac{1}{2^{n + 1}} + \frac{1}{3^{n + 1}}} = \lim_{n \ra \un} \frac{\frac{3^n + 2^n}{6^n}}{\frac{3^{n + 1} + 2^{n + 1}}{6^{n + 1}}} = \lim_{n \ra \un} \frac{\rkl{3^n + 2^n} \mal 6^{n + 1}}{6^n \rkl{3^{n + 1} + 2^{n + 1}}}\\
		&= 6 \mal \lim_{n \ra \un} \frac{3^n + 2^n}{3^{n + 1} + 2^{n + 1}} = 6 \mal \lim_{n \ra \un} \frac{1 + \rkl{\frac{2}{3}}^n}{3 + 2 \mal \rkl{\frac{2}{3}}^n}\\
		&= 6 \mal \frac{lim_{n \ra \un} \rkl{1+ \rkl{\frac{2}{3}}^n}}{\lim_{n \ra \un} \rkl{3 + 2 \rkl{\frac{2}{3}}^n}} = 6 \mal \frac{1}{3} = 2
		\end{align*}
		Konvergenzradius $r = q = 2$

\item $\sum_{k = 0}^{\un} \frac{\rkl{2k}!}{\rkl{k!}^k} \rkl{x + 2}^k = \sum_{k = 0}^{\un} \ub{\frac{\rkl{2k}!}{\rkl{k!}^k}}{=a_k} \mal \rkl{x - \ub{\rkl{-2}}{= x_0}}^k$\\
		Wurzelmethode
		\begin{align*}
		0 \klgl \lim_{k \ra \un} \sqrt[k]{\betrag{a_k}} &= \lim_{k \ra \un} \sqrt[k]{\betrag{\frac{(2k)!}{(k!)^k}}} = \lim_{k \ra \un} \sqrt[k]{(2k)!} \mal \frac{1}{\sqrt[k]{k!^k}}\\
		&= \lim_{k \ra \un} \frac{1}{k!} \mal \sqrt[k]{(2k)!} \klgl \lim_{k \ra \un} \frac{1}{k!} \sqrt[k]{(2k)^{2k}} = \lim_{k \ra \un} \frac{(2k)^2}{k!} = 0
		\end{align*}
		\Ra Konvergenzradius $r = \un$
\item 
\item 

\end{enumerate}

%\chapter{Testbl�tter}
\section{1. Test zur Vorlesung Mathematik 2}
\subsection{\texorpdfstring{$a, b \in \mb{R}$. Kreuzen Sie ALLE richtigen Antworten an.}{a, b element R. Kreuzen Sie ALLE richtigen Antworten an.}}
\begin{enumerate}[label=\alph*)]
\item Die Aussage $a > b \Ra \frac{a}{b} < 1$ gilt falls:
		\begin{multicols}{4}
		\doublespacing
		$\square$\hspace{0.5cm}$b > 0$\\
		$\square$\hspace{0.5cm}$b \neq 0$\\
		$\square$\hspace{0.5cm}$b = 1$\\
		$\square$\hspace{0.5cm}$b < 0$
		\end{multicols}

\item Die Aussage $a^2 > 0$ ist �quivalent zu:
		\begin{multicols}{2}
		\doublespacing
		$\square$\hspace{0.5cm}$\betrag{a} > 0$\\
		$\square$\hspace{0.5cm}$(a > 0) \oder (a < 0)$\\
		$\square$\hspace{0.5cm}$-\betrag{a} > 0$\\
		$\square$\hspace{0.5cm}$(a > 0) \und (a < 0)$
		\end{multicols}

\item Die Aussage $\betrag{a} < b$ ist �quivalent zu:
		\begin{multicols}{2}
		\doublespacing
		$\square$\hspace{0.5cm}$-b < a < b$\\
		$\square$\hspace{0.5cm}$(a < b) \oder (a > -b)$\\
		$\square$\hspace{0.5cm}$-b < -a < b$\\
		$\square$\hspace{0.5cm}$0 < a < b$\\
		$\square$\hspace{0.5cm}$(a < b) \und (a > -b)$
		\end{multicols}
\end{enumerate}

\subsection{Kreuzen Sie ALLE richtigen Antworten an.}
\begin{enumerate}[label=\alph*)]
\item Welche der folgenden Folgen sind beschr�nkt
		\begin{multicols}{4}
		\doublespacing
		$\square$\hspace{0.5cm}$((-1)^n)_{n \in \mb{N}_0}$\\
		$\square$\hspace{0.5cm}$\rklamm{\frac{1}{n}}_{n \in \N}$\\
		$\square$\hspace{0.5cm}$(-n^3)_{n \in \N}$\\
		$\square$\hspace{0.5cm}$(2 \mal n)_{n \in \N_0}$
		\end{multicols}

\item Gilt f�r zwei Folgen $(A_n)_{n \in \N}$ und $(b_n)_{n \in \N}$ dass $\lim_{n \ra \infty} a_n = \infty$ und $\lim_{n \ra \infty} b_n = \infty$, dann gilt \textbf{m�glicherweise} $\lim_{n \ra \infty} (a_n - b_n) = \dots$
		\begin{multicols}{3}
		\doublespacing
		$\square$\hspace{0.5cm}$\infty$\\
		$\square$\hspace{0.5cm}$-\infty$\\
		$\square$\hspace{0.5cm}$0$\\
		$\square$\hspace{0.5cm}$2$\\
		$\square$\hspace{0.5cm}$n$\\
		$\square$\hspace{0.5cm}existiert nicht
		\end{multicols}

\item $(a_n)_{n \in \N}$ sei eine Folge. Dann gilt:
		\begin{multicols}{2}
		\doublespacing
		$\square$\hspace{0.5cm}$a_n = O(n^3) \Ra a_n = o(n^3)$\\
		$\square$\hspace{0.5cm}$a_n = O(n^4) \Ra \forall p \klgl 4$: $a_n = O(n^p)$\\
		$\square$\hspace{0.5cm}$a_n = o(n^3) \Ra a_n = O(n^3)$\\
		$\square$\hspace{0.5cm}$a_n = O(n^4) \Ra \forall p \grgl 4$: $a_n = O(n^p)$
		\end{multicols}

\item F�r die Folge $(a_n)_{n \in \N}$ mit $a_1 = \frac{1}{8}$, $a_{n + 1} \ceq a_n^2 + a_1$ (f�r $n \in \N$) gilt:
		\begin{multicols}{2}
		\doublespacing
		$\square$\hspace{0.5cm}$(a_n)_{n \in \N}$ ist monoton wachsend\\
		$\square$\hspace{0.5cm}$(a_n)_{n \in \N}$ ist nach oben beschr�nkt\\
		$\square$\hspace{0.5cm}$(a_n)_{n \in \N}$ ist konvergent\\
		$\square$\hspace{0.5cm}$(a_n)_{n \in \N}$ ist monoton fallend\\
		$\square$\hspace{0.5cm}$(a_n)_{n \in \N}$ ist nach unten beschr�nkt\\
		$\square$\hspace{0.5cm}$(a_n)_{n \in \N}$ ist divergent
		\end{multicols}
\end{enumerate}

\subsection{Kreuzen Sie ALLE richtigen Antworten an.}
\begin{enumerate}[label=(\alph*)]
\item Wieviele M�glichkeiten gibt es aus den Ziffern 1-8 eine f�nfstellige Zahl zu bilden, wenn jede Ziffer nur einmal vorkommen darf?
		\begin{multicols}{3}
		\doublespacing
		$\square$\hspace{0.5cm}$\binom{8}{5}$\\
		$\square$\hspace{0.5cm}$4 \mal 5 \mal 6 \mal 7 \mal 8$\\
		$\square$\hspace{0.5cm}$\binom{8}{5} \mal 5!$\\
		$\square$\hspace{0.5cm}$\frac{8^5}{5!}$\\
		$\square$\hspace{0.5cm}$8^5$\\
		$\square$\hspace{0.5cm}$\frac{8!}{5!}$
		\end{multicols}

\item 4 Kochb�cher, 5 Physikb�cher und 6 Chemieb�cher sollen auf einem Regal nebeneinander gestellt werden. Auf wie viele Arten kann man das tun, wenn B�cher des gleichen Stoffgebietes nebeneinander gestellt werden sollen und alle B�cher verschieden sind?
		\begin{multicols}{3}
		\doublespacing
		$\square$\hspace{0.5cm}$3! \mal 4! \mal 5! \mal 6!$\\
		$\square$\hspace{0.5cm}$\frac{15!}{4! \mal 5! \mal 6!}$\\
		$\square$\hspace{0.5cm}$12 \mal \binom{11}{6}$\\
		$\square$\hspace{0.5cm}$12 \mal 7 \mal 2$\\
		$\square$\hspace{0.5cm}$\binom{15}{4} \mal \binom{11}{6} \mal \binom{5}{5}$\\
		$\square$\hspace{0.5cm}$15 \mal (6! + 5! + 4!)$
		\end{multicols}
\end{enumerate}

\subsection{Bestimmen Sie alle L�sungen der folgenden Ungleichung:}
\[\betrag{\frac{x - 3}{x - 1}} > 3\]

%</Text>%%%%%%%%%%%%%%%%%%%%%%%%%%%%%%%%%%%%%%%%%%%%%%%%%%%%%%%%%%%%%%%%%%%%%%%
\backmatter
%Literaturverzeichnis
%\newpage
%\addcontentsline{toc}{part}{Literaturverzeichnis}
%\bibliography{literatur}{}

%Stichwortverzeichnis
\newpage
\renewcommand{\indexname}{Stichwortverzeichnis}
\addcontentsline{toc}{part}{Stichwortverzeichnis}
\printindex

\end{document}


